\chapter{力学的拉格朗日表述}

\begin{problem}{9.1.2}
如图 9.2 所示, 一个均质的、半径为 R 的圆盘沿水平的 x 轴做纯滚动, 一根长 2l 的均质细棒与圆盘保持无滑动接触, 一端沿 x 轴滑动. 运动时, 圆盘与棒保持在同一竖直平面内, 选取适当的坐标, 写出约束关系, 并说明描述系统需用多少个独立坐标.
\begin{center}
\begin{tikzpicture}[>=Stealth, thick]
    \def\th{32}
    \def\R{1.05}
    \pgfmathsetmacro{\xc}{\R/tan(\th/2)}
    \def\rodlen{4.6}
    \def\rodthick{0.07}

    % Horizontal x-axis
    \draw (-0.3,0) -- ({\xc+\R+0.8},0) node[below left=1pt, font=\small] {$x$};

    % Angle theta_1
    \draw[semithick] (0.85,0) arc (0:\th:0.85);
    \node[font=\small] at ({1.10*cos(\th/2)}, {1.10*sin(\th/2)}) {$\theta_1$};

    % Homogeneous disk of radius R
    \coordinate (C) at (\xc,\R);
    \coordinate (T) at ({\xc*cos(\th)}, {\xc*sin(\th)});
    \draw (C) circle (\R);
    \fill (C) circle (1.3pt);
    \draw[->, semithick] (C) -- (T) node[midway, right=2pt, font=\small] {$R$};

    % Angular velocity indicator
    \draw[->, semithick] ($(C) + ({1.28*cos(24)}, {1.28*sin(24)})$) arc (24:-24:1.28);
    \node[right=2pt, font=\small] at ($(C) + (1.28,0)$) {$\dot{\theta}_2$};

    % Thin rod of length 2l
    \begin{scope}[rotate=\th]
        \draw[fill=white] (0,0) rectangle (\rodlen,\rodthick);

        % Dimension 2l above the rod
        \draw[thin] (0,\rodthick+0.08) -- (0,\rodthick+0.78);
        \draw[thin] (\rodlen,\rodthick+0.08) -- (\rodlen,\rodthick+0.78);
        \draw[<->, thin] (0,\rodthick+0.62) -- (\rodlen,\rodthick+0.62)
            node[midway, fill=white, inner sep=1.5pt, font=\small] {$2l$};

        % Dimension s along the rod to contact point
        \draw[thin] (2.60,-0.08) -- (2.60,-0.65);
        \draw[thin] (\xc,-0.08) -- (\xc,-0.65);
        \draw[<->, thin] (2.60,-0.50) -- (\xc,-0.50)
            node[midway, fill=white, inner sep=1pt, font=\small] {$s$};
    \end{scope}

    % Rod center of mass label (upright)
    \node[font=\small] at ({1.68*cos(\th)+0.34*sin(\th)}, {1.68*sin(\th)-0.34*cos(\th)}) {$(x_1, y_1)$};
\end{tikzpicture}

\end{center}
\end{problem}
\clearpage


\begin{problem}{9.1.3}
一根长为 $2a$ 、质量为 $m$ 的杆 $BC$ 用一根未伸长时长为 $b$ 、劲度系数为 $k$ 的系在杆的 $B$ 端的弹簧悬于固定点 $A$ ，如图 9.3 所示。系统的运动限于包含杆与弹簧的铅直平面内，选择适当的广义坐标，求作用于杆上的广义力分量。
\begin{center}
\begin{tikzpicture}[>=Stealth, thick]
    % Ceiling with hatching
    \draw (-1.2,4.2) -- (1.2,4.2);
    \foreach \x in {-1.05,-0.80,...,1.05} {
        \draw[semithick] (\x,4.2) -- (\x+0.15,4.38);
    }
    \node[left=2pt, font=\small] at (0,4.02) {$A$};

    % Spring k
    \draw (0,4.2) -- (0,3.85);
    \draw[decorate, decoration={coil, aspect=0.45, segment length=4.5pt, amplitude=4.2pt}]
        (0,3.85) -- (0,2.55);
    \node[right=4pt, font=\small] at (0.12,3.2) {$k$};
    \draw (0,2.55) -- (0,2.25);

    % Uniform rod BC
    \draw (-0.14,0.25) rectangle (0.14,2.25);
    \node[left=2pt, font=\small] at (-0.14,2.20) {$B$};
    \node[left=2pt, font=\small] at (-0.14,0.30) {$C$};
    \node[left=3pt, font=\small] at (-0.14,1.25) {$m$};

    % Dimension 2a
    \draw[thin] (0.28,2.25) -- (0.80,2.25);
    \draw[thin] (0.28,0.25) -- (0.80,0.25);
    \draw[<->, thin] (0.58,2.25) -- (0.58,0.25)
        node[midway, fill=white, inner sep=1.5pt, font=\small] {$2a$};
\end{tikzpicture}

\end{center}
\end{problem}
\clearpage


\begin{problem}{9.1.4}
一个在均匀重力场中运动的质点,如用球坐标来描述质点的运动,取竖直向上方向为极轴,求重力的三个广义力分量.
\end{problem}
\clearpage


\begin{problem}{9.1.5}
一个质量为 $m$ 、半径为 $a$ 的均质薄圆筒，在另一个质量为 $M$ 、半径为 $2a$ 的均质薄圆筒内部做无滑滚动，后者又在水平面上做无滑滚动。选大圆筒的角位移 $\theta$ 以及两圆筒的轴构成的平面与铅直面的夹角 $\varphi$ 为广义坐标，求作用于系统的所有力的广义力分量 $Q_{\theta}$ 和 $Q_{\varphi}$ 。
\end{problem}
\clearpage


\begin{problem}{9.1.6}
如图 9.5 所示, 一个质量为 $m$ 、半径为 $a$ 的均质圆柱体, 在一个质量为 $M$ 的木块中割出一个半径为 $b$ 的半圆柱形空心槽内做纯滚动, 木块又由一个劲度系数为 $k$ 的弹簧支承着, 可沿竖直导轨做无摩擦运动. 取木块的竖直向上的位移 $x$ 和圆柱中心的角位移 $\theta$ 为广义坐标, $x$ 和 $\theta$ 的零点迭在系统的平衡位置. 求作用于系统的所有力的广义力分量.
\begin{center}
\begin{tikzpicture}[>=Stealth, thick]
    \def\b{1.75}
    \def\a{0.75}
    \def\th{52}
    \pgfmathsetmacro{\ax}{(\b-\a)*sin(\th)}
    \pgfmathsetmacro{\ay}{-(\b-\a)*cos(\th)}

    % Outer vertical guide frame
    \draw (-2.30,0.70) -- (-2.30,-3.65) -- (2.30,-3.65) -- (2.30,0.70);

    % Block M with semicircular trough of radius b
    \draw (-2.15,0) -- (-1.75,0) arc (180:360:\b) -- (2.15,0) -- (2.15,-2.25) -- (-2.15,-2.25) -- cycle;
    \node[font=\small] at (-1.65,-1.85) {$M$};

    % Spring k supporting block M
    \draw (0,-3.65) -- (0,-3.32);
    \draw[decorate, decoration={zigzag, segment length=5.5pt, amplitude=3.2pt}]
        (0,-3.32) -- (0,-2.58);
    \draw (0,-2.58) -- (0,-2.25);
    \node[right=4pt, font=\small] at (0.05,-2.95) {$k$};

    % Central vertical axis x
    \draw[->, semithick] (0,-2.25) -- (0,0.95) node[left=1pt, font=\small] {$x$};

    % Center B and trough radius b
    \coordinate (B) at (0,0);
    \node[above left=1pt, font=\small] at (B) {$B$};
    \draw[->, semithick] (B) -- ({-\b*cos(22)}, {-\b*sin(22)})
        node[midway, above left=0pt, font=\small] {$b$};

    % Cylinder m of radius a centered at A
    \coordinate (A) at (\ax,\ay);
    \draw (A) circle (\a);
    \fill (A) circle (1.2pt);
    \node[right=2pt, font=\small] at (A) {$A$};
    \node[font=\small] at ($(A)+(0.20,0.34)$) {$m$};

    % Angle theta and radius a
    \draw[semithick] (B) -- (A);
    \draw[semithick] (0,-0.48) arc (270:{270+\th}:0.48);
    \node[font=\small] at ({0.68*cos(270+\th/2)}, {0.68*sin(270+\th/2)}) {$\theta$};

    \draw[->, semithick] (A) -- ($(A)+({\a*cos(-128)}, {\a*sin(-128)})$)
        node[midway, below right=0pt, font=\small] {$a$};
\end{tikzpicture}

\end{center}
\end{problem}
\clearpage


\begin{problem}{9.1.7}
$xy$ 平面内的任何点可改用 $c, y$ 为坐标或 $c, x$ 为坐标，其中 $c$ 满足 $xy = c$ . 如重力沿 $y$ 轴负向，求质量为 $m$ 的质点所受重力的广义力.
\end{problem}
\clearpage


\begin{problem}{9.1.8}
一个质量为 $m$ 的质点系在不可伸长的绳子上，绳子穿过板 $B$ 的小孔栓于固

定点, 板以 $s = A \sin \omega t$ 沿竖直轴 $y$ 做上下运动, 其中 $A, \omega$ 为常量, 质点在图9.6所示的 $xy$ 平面内运动, 求质点受到的广义力(不计小孔处的摩擦力).
\begin{center}
\begin{tikzpicture}[>=Stealth, thick]
    \def\th{32}
    \def\rlen{2.05}
    \def\yplate{-0.95}

    % Ceiling with hatching
    \draw (-1.6,1.15) -- (1.6,1.15);
    \foreach \x in {-1.45,-1.15,...,1.45} {
        \draw[semithick] (\x,1.15) -- (\x+0.18,1.33);
    }

    % Vertical y-axis and horizontal x-axis
    \draw[->, semithick] (0,1.15) -- (0,-3.15) node[left=1pt, font=\small] {$y$};
    \draw[->, semithick] (0,0) -- (2.10,0) node[above left=1pt, font=\small] {$x$};
    \node[left=2pt, font=\small] at (0,0) {$O$};

    % Moving plate B with central hole
    \draw (-1.65,\yplate+0.14) -- (-0.07,\yplate+0.14) -- (-0.07,\yplate) -- (-1.65,\yplate);
    \draw (1.50,\yplate+0.14) -- (0.07,\yplate+0.14) -- (0.07,\yplate) -- (1.50,\yplate);
    \node[above=2pt, font=\small] at (-1.40,\yplate+0.14) {$B$};

    % Displacement s of plate B
    \draw[thin] (1.55,\yplate) -- (1.90,\yplate);
    \draw[<->, thin] (1.72,0) -- (1.72,\yplate) node[midway, right=2pt, font=\small] {$s$};

    % String and particle m
    \coordinate (H) at (0,\yplate);
    \coordinate (M) at ({\rlen*sin(\th)}, {\yplate-\rlen*cos(\th)});
    \draw (H) -- (M);
    \draw[fill=white] (M) circle (0.14);
    \node[right=3pt, font=\small] at (M) {$m$};

    % Angle theta and string length r
    \draw[semithick] (0,\yplate-0.60) arc (270:{270+\th}:0.60);
    \node[font=\small] at ({0.80*sin(\th/2)}, {\yplate-0.80*cos(\th/2)}) {$\theta$};
    \node[above right=1pt, font=\small] at ($(H)!0.48!(M)$) {$r$};
\end{tikzpicture}

\end{center}
\end{problem}
\clearpage


\begin{problem}{9.1.9}
图 9.7 中木块 A、B 淹没在黏性液体中，黏性力与速度成正比，比例系数分别为 $\alpha_{1}$ 、 $\alpha_{2}$ ，用图中 y、 $y_{3}$ 为广义坐标，求广义黏性力.
\begin{center}
\begin{tikzpicture}[>=Stealth, thick]
    % Ceiling with hatching
    \draw (-2.2,0) -- (2.2,0);
    \foreach \x in {-2.05,-1.80,...,2.05} {
        \draw[semithick] (\x,0) -- (\x+0.15,0.16);
    }

    % Spring k and pulley
    \draw (0,0) -- (0,-0.25);
    \draw[decorate, decoration={zigzag, segment length=5pt, amplitude=3pt}]
        (0,-0.25) -- (0,-0.75);
    \draw (0,-0.75) -- (0,-1.55);
    \node[right=4pt, font=\small] at (0.04,-0.50) {$k$};

    \coordinate (P) at (0,-1.55);
    \draw (P) circle (0.62);
    \fill (P) circle (1.4pt);

    % Container and viscous liquid surface
    \draw (-2.42,-1.92) -- (-2.22,-2.12) -- (-2.22,-4.30) -- (2.22,-4.30) -- (2.22,-2.12) -- (2.42,-1.92);
    \draw[semithick] (-2.22,-2.22) -- (2.22,-2.22);
    \node[font=\small] at (-1.42,-4.00) {黏性液体};

    % Left mass m1 (A)
    \coordinate (A) at (-0.62,-3.32);
    \draw (-0.62,-1.55) -- (A);
    \draw[fill=white] (A) circle (0.31);
    \node[font=\small] at (A) {$m_1$};
    \node[font=\small] at ($(A)+(0.36,0.28)$) {$\alpha_1$};
    \node[font=\small] at ($(A)+(0.38,-0.28)$) {$A$};

    % Right mass m2 (B)
    \coordinate (B) at (0.62,-3.72);
    \draw (0.62,-1.55) -- (B);
    \draw[fill=white] (B) circle (0.31);
    \node[font=\small] at (B) {$m_2$};
    \node[font=\small] at ($(B)+(0.36,0.28)$) {$\alpha_2$};
    \node[font=\small] at ($(B)+(0.38,-0.28)$) {$B$};

    % Coordinate dimensions y1, y3, y, y2
    \draw[thin] (-1.85,-3.32) -- (-0.98,-3.32);
    \draw[<->, thin] (-1.68,0) -- (-1.68,-3.32)
        node[midway, fill=white, inner sep=1.5pt, font=\small] {$y_1$};

    \draw[thin] (-1.32,-1.55) -- (-0.72,-1.55);
    \draw[<->, thin] (-1.16,-1.55) -- (-1.16,-3.32)
        node[pos=0.62, fill=white, inner sep=1.5pt, font=\small] {$y_3$};

    \draw[thin] (0.72,-1.55) -- (1.42,-1.55);
    \draw[<->, thin] (1.25,0) -- (1.25,-1.55)
        node[midway, fill=white, inner sep=1.5pt, font=\small] {$y$};

    \draw[thin] (0.98,-3.72) -- (1.95,-3.72);
    \draw[<->, thin] (1.78,0) -- (1.78,-3.72)
        node[midway, fill=white, inner sep=1.5pt, font=\small] {$y_2$};
\end{tikzpicture}

\end{center}
\end{problem}
\clearpage


\begin{problem}{9.1.10}
图 9.8 中棒形磁铁和两个单个磁极均受到铁质薄板的黏性阻力, 阻力与相对速度成正比, 比例系数分别为 $\alpha_{1}, \alpha_{2}, \alpha_{3}$ . 设磁棒、磁极和薄板均只做竖直方向的运动, 运动中, 磁棒、磁极均与薄板保持接触. 求广义黏性力 $Q_{y_{1}}$ 和 $Q_{y_{2}}$ .
\begin{center}
\begin{tikzpicture}[>=Stealth, thick]
    % Ceiling with hatching
    \draw (-2.5,0) -- (2.5,0);
    \foreach \x in {-2.35,-2.10,...,2.35} {
        \draw[semithick] (\x,0) -- (\x+0.15,0.16);
    }

    % Two upper pulleys with suspension brackets
    \foreach \sx in {-1.20, 1.20} {
        \draw (\sx,0) -- (\sx,-1.0);
        \draw[fill=white] (\sx-0.07,0) rectangle (\sx+0.07,-1.0);
        \draw (\sx,-1.0) circle (0.45);
        \fill (\sx,-1.0) circle (1.3pt);
    }

    % Iron plate (铁质薄板)
    \draw (-2.35,-1.65) rectangle (2.35,-5.95);
    \node[font=\small, align=center] at (1.88,-4.55) {铁\\质\\薄\\板};

    % Outer cords from upper pulleys to iron plate
    \draw[semithick] (-1.65,-1.0) -- (-1.65,-1.65);
    \draw[semithick] (1.65,-1.0) -- (1.65,-1.65);

    % Dimension y1
    \draw[<->, thin] (0,0) -- (0,-1.65)
        node[midway, fill=white, inner sep=1.5pt, font=\small] {$y_1$};

    % Bar magnet (磁棒) and inner cords
    \draw[semithick] (-0.75,-1.0) -- (-0.75,-2.05);
    \draw[semithick] (0.75,-1.0) -- (0.75,-2.05);
    \draw[fill=white] (-1.75,-2.05) rectangle (1.75,-2.35);
    \node[left=2pt, font=\small] at (-1.75,-2.20) {$\mathrm{S}$};
    \node[right=2pt, font=\small] at (1.75,-2.20) {$\mathrm{N}$};
    \node[below=1pt, font=\small] at (-0.75,-2.35) {$\alpha_1$};
    \draw[thin] (0.65,-2.35) -- (0.95,-2.72) node[below, font=\small] {磁棒};

    % Lower pulley suspended from bar magnet
    \draw[fill=white] (-0.07,-2.35) rectangle (0.07,-3.55);
    \draw (0,-3.55) circle (0.50);
    \fill (0,-3.55) circle (1.3pt);

    % Two magnetic poles A and B
    \draw[semithick] (-0.50,-3.55) -- (-0.50,-4.90);
    \draw[fill=white] (-0.76,-4.90) rectangle (-0.24,-5.40);
    \node[font=\small] at (-0.50,-5.15) {$N$};
    \node[above right=0pt, font=\small] at (-0.25,-4.90) {$A$};
    \node[below left=0pt, font=\small] at (-0.72,-5.34) {$\alpha_2$};

    \draw[semithick] (0.50,-3.55) -- (0.50,-4.55);
    \draw[fill=white] (0.24,-4.55) rectangle (0.76,-5.05);
    \node[font=\small] at (0.50,-4.80) {$N$};
    \node[above right=0pt, font=\small] at (0.70,-4.55) {$B$};
    \node[below right=0pt, font=\small] at (0.74,-5.00) {$\alpha_3$};

    \node[font=\small] (pole) at (0.05,-5.68) {磁极};
    \draw[thin] (-0.24,-5.28) -- (pole.north west);
    \draw[thin] (0.50,-5.05) -- (pole.north east);

    % Dimension y2
    \draw[thin] (-1.45,-3.55) -- (-0.60,-3.55);
    \draw[thin] (-1.45,-4.90) -- (-0.84,-4.90);
    \draw[<->, thin] (-1.25,-3.55) -- (-1.25,-4.90)
        node[midway, fill=white, inner sep=1.5pt, font=\small] {$y_2$};
\end{tikzpicture}

\end{center}
\end{problem}
\clearpage


\begin{problem}{9.1.11}
图 9.9 中 $m_{1}$ 、 $m_{2}$ 限于在光滑的水平线上运动， $m_{1}$ 、 $m_{2}$ 均固连着一个减震器，活塞 $P_{1}$ 、 $P_{2}$ 和 $m_{3}$ 固连于水平棒上，活塞的作用力大小与相对速度的 n 次方成正比，方向与相对速度方向相反。求广义力 $Q_{x_{1}}$ 、 $Q_{q_{1}}$ 、 $Q_{q_{2}}$ （弹簧力用势能处理）。
\begin{center}
\begin{tikzpicture}[>=Stealth, thick]
    % Coordinate axes
    \draw[->, semithick] (0,0) -- (0,2.05) node[left=1pt, font=\small] {$y$};
    \draw[->, semithick] (0,0) -- (6.0,0) node[below left=1pt, font=\small] {$x$};

    % Spring k1
    \draw[semithick] (0,0.32) -- (0.18,0.32);
    \draw[semithick, decorate, decoration={zigzag, segment length=4.5pt, amplitude=2.6pt}]
        (0.18,0.32) -- (0.72,0.32);
    \draw[semithick] (0.72,0.32) -- (0.90,0.32);
    \node[above=2pt, font=\small] at (0.45,0.44) {$k_1$};

    % Block m1 and damper alpha1 / P1
    \draw[fill=white] (0.90,0) rectangle (2.20,0.72);
    \node[font=\small] at (1.55,0.36) {$m_1$};
    \node[below=2pt, font=\small] at (0.90,0) {$x_1$};

    \draw[fill=white] (1.15,0.72) rectangle (1.95,1.25);
    \node[font=\small] at (1.40,0.98) {$\alpha_1$};
    % Piston P1
    \draw[semithick] (1.65,0.72) -- (1.65,1.32);
    \draw[semithick] (1.95,0.72) -- (1.95,1.32);
    \draw[semithick] (1.58,1.32) -- (2.02,1.32);
    \foreach \yy in {0.84,1.08,1.16,1.24} {
        \draw[thin] (1.65,\yy) -- (1.95,\yy);
    }
    \node[above=1pt, font=\small] at (1.80,1.32) {$P_1$};

    % Spring k2
    \draw[semithick] (2.20,0.32) -- (2.32,0.32);
    \draw[semithick, decorate, decoration={zigzag, segment length=4.5pt, amplitude=2.6pt}]
        (2.32,0.32) -- (2.73,0.32);
    \draw[semithick] (2.73,0.32) -- (2.85,0.32);
    \node[above=2pt, font=\small] at (2.52,0.44) {$k_2$};

    % Horizontal rod from P1 to m3
    \draw[semithick] (1.95,1.02) -- (2.85,1.02);
    \draw[semithick] (1.95,1.12) -- (2.85,1.12);

    % Block m3
    \draw[fill=white] (2.85,0) rectangle (3.35,1.50);
    \node[font=\small] at (3.10,0.75) {$m_3$};
    \node[below=2pt, font=\small] at (2.85,0) {$x_3$};

    % Spring k3
    \draw[semithick] (3.35,0.32) -- (3.47,0.32);
    \draw[semithick, decorate, decoration={zigzag, segment length=4.5pt, amplitude=2.6pt}]
        (3.47,0.32) -- (3.88,0.32);
    \draw[semithick] (3.88,0.32) -- (4.00,0.32);
    \node[above=2pt, font=\small] at (3.68,0.44) {$k_3$};

    % Horizontal rod from m3 to P2
    \draw[semithick] (3.35,1.02) -- (4.25,1.02);
    \draw[semithick] (3.35,1.12) -- (4.25,1.12);

    % Block m2 and damper alpha2 / P2
    \draw[fill=white] (4.00,0) rectangle (5.45,0.72);
    \node[font=\small] at (4.72,0.36) {$m_2$};
    \node[below=2pt, font=\small] at (4.00,0) {$x_2$};

    \draw[fill=white] (4.25,0.72) rectangle (5.05,1.25);
    \node[font=\small] at (4.80,0.98) {$\alpha_2$};
    % Piston P2
    \draw[semithick] (4.25,0.72) -- (4.25,1.32);
    \draw[semithick] (4.55,0.72) -- (4.55,1.32);
    \draw[semithick] (4.18,1.32) -- (4.62,1.32);
    \foreach \yy in {0.84,1.08,1.16,1.24} {
        \draw[thin] (4.25,\yy) -- (4.55,\yy);
    }
    \node[above=1pt, font=\small] at (4.40,1.32) {$P_2$};

    % Dimension q2 (top)
    \draw[thin] (0.90,1.62) -- (0.90,2.02);
    \draw[thin] (2.85,1.62) -- (2.85,2.02);
    \draw[<->, thin] (0.90,1.82) -- (2.85,1.82) node[midway, above=1pt, font=\small] {$q_2$};

    % Dimension q1 (bottom)
    \draw[thin] (0.90,-0.45) -- (0.90,-0.88);
    \draw[thin] (4.00,-0.45) -- (4.00,-0.88);
    \draw[<->, thin] (0.90,-0.72) -- (4.00,-0.72) node[midway, above=1pt, font=\small] {$q_1$};
\end{tikzpicture}

\end{center}
\end{problem}
\clearpage


\begin{problem}{9.1.13}
作用在 $xy$ 平面上的哑铃的黏性阻力与其速度成正比，比例系数分别为 $\alpha_{1}$ 、

$\alpha_{2}$ ，以图9.11所示的 $x_{1}$ 、 $y_{1}$ 、 $\varphi$ 为广义坐标，求广义力.
\begin{center}
\begin{tikzpicture}[>=Stealth, thick]
    % Coordinate axes
    \draw[->, semithick] (0,0) -- (0,3.25) node[left=1pt, font=\small] {$y$};
    \draw[->, semithick] (0,0) -- (4.30,0) node[below left=1pt, font=\small] {$x$};

    % Dumbbell endpoints
    \coordinate (P2) at (1.20,0.85);
    \coordinate (P1) at (3.00,2.65);

    % Horizontal dashed line and angle phi
    \draw[dashed, semithick] (P2) -- (2.75,0.85);
    \draw (P2) -- (P1) node[midway, above left=1pt, font=\small] {$l$};
    \draw[semithick] ($(P2)+(0.82,0)$) arc (0:45:0.82);
    \node[font=\small] at ($(P2)+({1.02*cos(22.5)}, {1.02*sin(22.5)})$) {$\varphi$};

    % Mass 2 and Mass 1 circles
    \draw[fill=white] (P2) circle (0.18);
    \node[above left=2pt, font=\small] at (P2) {$\alpha_2$};
    \node[below=5pt, font=\small] at (P2) {$(x_2, y_2)$};

    \draw[fill=white] (P1) circle (0.18);
    \node[above left=2pt, font=\small] at (P1) {$\alpha_1$};
    \node[right=5pt, font=\small] at (P1) {$(x_1, y_1)$};
\end{tikzpicture}

\end{center}
\end{problem}
\clearpage


\begin{problem}{9.1.14}
平板在 $xy$ 平面上运动, 坐标系 $O'x'y'$ 固连于平板, 用 $O'$ 点在静坐标系 $Oxy$ 

中的坐标 x、y 及 $x'$ 、x 轴的夹角 $\varphi$ 表示平板的位形，面元 $dx'dy'$ 受到的黏性阻力为

\[
\mathrm{d} f _ {x} = - a \dot {x} \mathrm{d} x ^ {\prime} \mathrm{d} y ^ {\prime} \mathrm{d} f _ {y} = - a \dot {y} \mathrm{d} x ^ {\prime} \mathrm{d} y ^ {\prime}
\]

其中 $\dot{x},\dot{y}$ 分别是 $\dot{x} (x^{\prime},y^{\prime})$ 、 $\dot{x} (x^{\prime},y^{\prime})$ 的缩写，不要与 $O^{\prime}$ 点的速度分量 $\dot{x},\dot{y}$ 发生混淆.求平板受到的广义力.
\end{problem}
\clearpage


\begin{problem}{9.1.15}
若上题所述的平板是边长为 $2a$ 和 $2b$ 的矩形板，取 $O'$ 在其中心， $x', y'$ 轴平行于矩形的边, $x'$ 轴平行于边长为 $2\alpha$ 的边, $y'$ 轴平行于边长为 $2b$ 的边,求黏性阻力的广义力.
\end{problem}
\clearpage


\begin{problem}{9.1.16}
一个半径为 r 的薄圆盘在一块涂油的平板上运动, 设圆盘的面元受到油的黏性阻力, 每单位面积、每单位速度受到的黏性力为 $\alpha$ , 证明广义黏性力为

\[
Q _ {x} = - \alpha A \dot {x}, Q _ {y} = - \alpha A \dot {y}, Q _ {\varphi} = - \frac {1}{2} \alpha A r ^ {2} \dot {\varphi}
\]

其中 $A=\pi r^{2}$ ，x、y 是盘心的位置坐标， $\varphi$ 表示圆盘的角位置.
\end{problem}
\clearpage


\begin{problem}{9.1.17}
若 9.1.10 题中磁铁与铁质薄板间的阻力与相对速度的平方成正比, 求 $Q_{y_1}$ 和 $Q_{y_2}$ .
\end{problem}
\clearpage


\begin{problem}{9.1.18}
图 9.13 中木块 A、B 固连于长度为 l 的轻棒的两端，在木块 C 和 D 上沿棒的方向滑动，木块 C 在水平面上也可沿棒的方向滑动，D 是固定的。图中画出了各接触面间的摩擦因数 $\mu_{1}$ 、 $\mu_{2}$ 、 $\mu_{3}$ ，套在棒上的木块 E 可沿棒无摩擦滑动。木块 A、B、C、D、E 的质量分别为 $m_{1}$ 、 $m_{2}$ 、 $m_{3}$ 、 $m_{4}$ 、 $m_{5}$ ，如运动中 $\dot{x}_{1} > \dot{x}_{3} > 0$ ，求这些摩擦力的广义力 $Q_{x_{1}}$ 、 $Q_{q_{1}}$ 和 $Q_{q_{2}}$ 。
\begin{center}
\begin{tikzpicture}[>=Stealth, line width=0.6pt, font=\small]
    % Left vertical reference wall and horizontal ground
    \draw (0,0) -- (0,2.85);
    \draw (0,0) -- (6.15,0);

    % Ground hatching under fixed block D
    \foreach \x in {4.35, 4.50, ..., 5.55} {
        \draw[line width=0.45pt] (\x,0) -- (\x-0.12,-0.14);
    }

    % Block C (m3)
    \draw (0.65,0) rectangle (1.90,0.52);
    \node at (1.275,0.26) {$m_3$};
    \node[anchor=west] at (1.90,0.44) {$C$};
    \node[anchor=west] at (1.90,0.13) {$\mu_3$};

    % Block D (m4)
    \draw (4.30,0) rectangle (5.60,0.52);
    \node at (4.95,0.26) {$m_4$};
    \node[anchor=west] at (5.60,0.44) {$D$};

    % Block A (m1)
    \draw (1.15,0.52) rectangle (1.65,1.20);
    \node at (1.40,0.86) {$m_1$};
    \node[anchor=south west, inner sep=1pt] at (1.66,1.10) {$A$};
    \node[anchor=west] at (1.72,0.64) {$\mu_1$};

    % Block B (m2)
    \draw (5.05,0.52) rectangle (5.55,1.20);
    \node at (5.30,0.86) {$m_2$};
    \node[anchor=south west, inner sep=1pt] at (5.56,1.10) {$B$};
    \node[anchor=west] at (5.60,0.64) {$\mu_2$};

    % Light rod connecting A and B (dashed inside E)
    \draw (1.65,0.92) -- (2.80,0.92);
    \draw (1.65,0.82) -- (2.80,0.82);
    \draw[dashed] (2.80,0.92) -- (3.30,0.92);
    \draw[dashed] (2.80,0.82) -- (3.30,0.82);
    \draw (3.30,0.92) -- (5.05,0.92);
    \draw (3.30,0.82) -- (5.05,0.82);

    % Block E (m5)
    \draw (2.80,0.42) rectangle (3.30,1.50);
    \node at (3.05,0.62) {$m_5$};
    \node[anchor=south west, inner sep=1pt] at (3.32,1.36) {$E$};

    % Vertical extension lines
    \draw[line width=0.45pt] (0.65,0.52) -- (0.65,1.98);
    \draw[line width=0.45pt] (1.40,1.20) -- (1.40,2.85);
    \draw[line width=0.45pt] (3.05,1.50) -- (3.05,2.30);
    \draw[line width=0.45pt] (5.30,1.20) -- (5.30,2.85);

    % Dimension / coordinate arrows
    \draw[<->] (0,2.60) -- (1.40,2.60) node[midway, above] {$x_1$};
    \draw[<->] (1.40,2.60) -- (5.30,2.60) node[midway, above] {$l$};
    \draw[<->] (1.40,2.14) -- (3.05,2.14) node[midway, above] {$q_1$};
    \draw[<->] (0,1.80) -- (0.65,1.80) node[midway, above] {$x_3$};
    \draw[<->] (0.65,1.80) -- (1.40,1.80) node[midway, above] {$q_2$};
\end{tikzpicture}

\end{center}
\end{problem}
\clearpage


\begin{problem}{9.1.19}
一个半径为 $R$ 的圆环在一个粗糙的平面上运动, 用环心的坐标 $x, y$ 及圆环的角位移 $\varphi$ 为广义坐标, 设 $f$ 是单位长度的环受到的摩擦力的大小, 求环受到的广义摩擦力 (不要求算出积分).

\[
\begin{array}{r l} \text {   解   考虑图   9.14   中位于   } & \alpha \sim \alpha + \mathrm{d} \alpha \text {   的环元的速度   } \\ \nu (\alpha) = & \dot {x} i + \dot {y} j + \dot {\varphi} k \times [ R \cos (\varphi + \alpha) i + R \sin (\varphi + \alpha) j ] \\ & = [ \dot {x} - R \dot {\varphi} \sin (\varphi + \alpha) ] i + [ \dot {y} + R \dot {\varphi} \cos (\varphi + \alpha) ] j \end{array}
\]

此环元受到的摩擦力为

\[
- f \frac {v (\alpha)}{v (\alpha)} R d \alpha = - f \frac {[ \dot {x} - R \dot {\varphi} \sin (\varphi + \alpha) ] i + [ \dot {y} + R \dot {\varphi} \cos (\varphi + \alpha) ] j}{\left\{\left[ \dot {x} - R \dot {\varphi} \sin (\varphi + \alpha) \right] ^ {2} + [ \dot {y} + R \dot {\varphi} \cos (\varphi + \alpha) ] ^ {2} \right\} ^ {1 / 2}} R d \alpha
\]

此环元的位矢为

\[
\begin{array}{r l} & {\bm {r} (\alpha) = [ x + R \mathrm{cos} (\varphi + \alpha) ] \bm {i} + [ y + R \mathrm{sin} (\varphi + \alpha) ] \bm {j}} \\ & {\delta \bm {r} (\alpha) = [ \delta x - R \mathrm{sin} (\varphi + \alpha) \delta \varphi ] \bm {i} + [ \delta y + R \mathrm{cos} (\varphi + \alpha) \delta \varphi ] \bm {j}} \end{array}
\]

摩擦力对此环元做的虚功为

\[
\begin{array}{r l} & {- f \frac {\boldsymbol {v} (\alpha) \cdot \delta \boldsymbol {r} (\alpha)}{\boldsymbol {v} (\alpha)} R d \alpha} \\ & {= - f \frac {[ \dot {x} - R \dot {\varphi} \sin (\varphi + \alpha) ] (\delta x - R \sin (\varphi + \alpha) \delta \varphi)}{[ [ \dot {x} - R \dot {\varphi} \sin (\varphi + \alpha) ] ^ {2} + [ y + R \dot {\varphi} \cos (\varphi + \alpha) ] ^ {2} ] ^ {1 / 2}} R d \alpha} \\ & {- f \frac {[ \dot {y} + R \dot {\varphi} \cos (\varphi + \alpha) ] (\delta y + R \cos (\varphi + \alpha) \delta \varphi)}{[ [ \dot {x} - R \dot {\varphi} \sin (\varphi + \alpha) ] ^ {2} + [ \dot {y} + R \dot {\varphi} \cos (\varphi + \alpha) ] ^ {2} ] ^ {1 / 2}} R d \alpha} \\ & {= - f \frac {[ \dot {x} - R \dot {\varphi} \sin (\varphi + \alpha) ] \delta x + [ \dot {y} + R \dot {\varphi} \cos (\varphi + \alpha) ] \delta y}{[ [ \dot {x} - R \dot {\varphi} \sin (\varphi + \alpha) ] ^ {2} + [ \dot {y} + R \dot {\varphi} \cos (\varphi + \alpha) ] ^ {2} ] ^ {1 / 2}} R d \alpha} \\ & {- f \frac {\{R ^ {2} \dot {\varphi} + R [ \dot {y} \cos (\varphi + \alpha) - \dot {x} \sin (\varphi + \alpha) ] \} \delta \varphi}{[ [ \dot {x} - R \dot {\varphi} \sin (\varphi + \alpha) ] ^ {2} + [ \dot {y} + R \dot {\varphi} \cos (\varphi + \alpha) ] ^ {2} ] ^ {1 / 2}} R d \alpha} \end{array}
\]

摩擦力对圆环做的虚功为

\[
\begin{array}{r l} \delta W & = - f \int_ {0} ^ {2 \pi} \frac {\boldsymbol {v} (\alpha) \cdot \delta \boldsymbol {r} (\alpha)}{v (\alpha)} R d \alpha \\ & = - f \int_ {0} ^ {2 \pi} \frac {\dot {x} - R \dot {\varphi} \sin (\varphi + \alpha)}{\left\{\left[ \dot {x} - R \dot {\varphi} \sin (\varphi + \alpha) \right] ^ {2} + \left[ \dot {y} + R \dot {\varphi} \cos (\varphi + \alpha) \right] ^ {2} \right\} ^ {1 / 2}} R d \alpha \delta x \\ & - f \int_ {0} ^ {2 \pi} \frac {\dot {y} + R \dot {\varphi} \cos (\varphi + \alpha)}{\left\{\left[ \dot {x} - R \dot {\varphi} \sin (\varphi + \alpha) \right] ^ {2} + \left[ \dot {y} + R \dot {\varphi} \cos (\varphi + \alpha) \right] ^ {2} \right\} ^ {1 / 2}} R d \alpha \delta y \\ & - f \int_ {0} ^ {2 \pi} \frac {R ^ {2} \dot {\varphi} + R [ \dot {y} \cos (\varphi + \alpha) - \dot {x} \sin (\varphi + \alpha) ]}{\left\{\left[ \dot {x} - R \dot {\varphi} \sin (\varphi + \alpha) \right] ^ {2} + \left[ \dot {y} + R \dot {\varphi} \cos (\varphi + \alpha) \right] ^ {2} \right\} ^ {1 / 2}} R d \alpha \delta \varphi \\ Q _ {x} & = - f \int_ {0} ^ {2 \pi} \frac {\dot {x} - R \dot {\varphi} \sin (\varphi + \alpha)}{\left\{\left[ \dot {x} - R \dot {\varphi} \sin (\varphi + \alpha) \right] ^ {2} + \left[ \dot {y} + R \dot {\varphi} \cos (\varphi + \beta) \right] ^ {2} \right\} ^ {1 / 2}} R d \alpha \end{array}
\]

所以

\[
\begin{array}{l} Q _ {y} = - f \int_ {0} ^ {2 \pi} \frac {\dot {y} + R \dot {\varphi} \cos (\varphi + \alpha)}{\left\{\left[ \dot {x} - R \dot {\varphi} \sin (\varphi + \alpha) \right] ^ {2} + \left[ \dot {y} + R \dot {\varphi} \cos (\varphi + \alpha) \right] ^ {2} \right\} ^ {1 / 2}} R d \alpha \\ Q _ {\varphi} = - f \int_ {0} ^ {2 \pi} \frac {R ^ {2} \dot {\varphi} + R [ \dot {y} \cos (\varphi + \alpha) - \dot {x} \sin (\varphi + \alpha) ]}{\left\{\left[ \dot {x} - R \dot {\varphi} \sin (\varphi + \alpha) \right] ^ {2} + \left[ \dot {y} + R \dot {\varphi} \cos (\varphi + \alpha) \right] ^ {2} \right\} ^ {1 / 2}} R d \alpha \end{array}
\]
\begin{center}
\begin{tikzpicture}[>=Stealth, line width=0.6pt, font=\small]
    % Coordinate axes
    \draw[->] (0,0) -- (4.5,0) node[below] {$x$};
    \draw[->] (0,0) -- (0,4.0) node[left] {$y$};

    % Ring center and circle
    \coordinate (O) at (2.45,1.85);
    \draw (O) circle (1.45);

    % Dashed horizontal reference line
    \draw[dashed] (O) -- ++(1.45,0);

    % Body-fixed axes x', y' and radius at angle phi + alpha
    \draw (O) -- ++(32:2.05) node[above right, inner sep=1pt] {$x'$};
    \draw (O) -- ++(122:2.05) node[above left, inner sep=1pt] {$y'$};
    \draw (O) -- ++(66:1.45);

    % Labels
    \node[below left, inner sep=1pt] at ($(O)+(-0.02,-0.04)$) {$O'$};
    \node[below right, inner sep=1pt] at ($(O)+(0.22,-0.08)$) {$(x,y)$};
    \node at ($(O)+(16:0.78)$) {$\varphi$};
    \node at ($(O)+(49:0.75)$) {$\alpha$};
\end{tikzpicture}

\end{center}
\end{problem}
\clearpage


\begin{problem}{9.1.20}
一个均质的三脚架, 三只脚长度相等, 置于光滑的水平面上, 为使它不致滑下, 用一绳套住三脚架的底部. 若平衡时, 每只脚与铅直线的夹角为 $\theta$ . 试用虚功原理证明: 绳中张力 $T$ 与三脚架重量 $W$ 之比为 $\frac{T}{W} = \frac{1}{6\sqrt{3}}\tan \theta$ .
\begin{center}
\begin{tikzpicture}[>=Stealth, line width=0.6pt, font=\small]
    % Coordinates: O is centroid of triangle ABC
    \coordinate (O) at (0,0);
    \coordinate (A) at (-2.10,0.35);
    \coordinate (B) at (0.65,-1.20);
    \coordinate (C) at (1.45,0.85);
    \coordinate (D) at (0,2.55);
    \coordinate (M) at ($(B)!0.5!(C)$);

    % Coordinate axes
    \draw[->] (O) -- (0,3.35) node[left] {$y$};
    \draw[->] (O) -- (2.45,1.44) node[above left, inner sep=2pt] {$x$};

    % Tripod legs and base rope triangle
    \draw (D) -- (A);
    \draw (D) -- (B);
    \draw (D) -- (C);
    \draw (A) -- (B) -- (C) -- cycle;
    \draw (A) -- (M);

    % Angle theta between DA and DO
    \pic[draw, angle radius=0.65cm, "$\theta$", angle eccentricity=1.40] {angle = A--D--O};

    % Labels
    \node[above left, inner sep=2pt] at (D) {$D$};
    \node[left] at (A) {$A$};
    \node[below] at (B) {$B$};
    \node[above, inner sep=3pt] at (C) {$C$};
    \node[below, inner sep=2pt] at ($(O)+(0,-0.04)$) {$O$};
\end{tikzpicture}

\end{center}
\end{problem}
\clearpage


\begin{problem}{9.1.21}
长度均为 l 的四根轻棒用光滑铰链构成一菱形 ABCD, AB 和 AD 两边支在同一水平线相距为 2a 的两个钉子上, BD 间用一轻绳连接, C 点处挂一重量为 W 的重物, 若 A 的顶角为 2$\alpha$. 如图 9.16 所示. 试用虚功原理求绳中的张力.
\begin{center}
\begin{tikzpicture}[>=Stealth, line width=0.6pt, font=\small]
    % Rhombus vertices
    \coordinate (A) at (0,2.40);
    \coordinate (B) at (-1.60,0.20);
    \coordinate (C) at (0,-2.00);
    \coordinate (D) at (1.60,0.20);
    \coordinate (O) at (0,1.30);

    % Peg contact coordinates at y = 1.30 (a = 0.80)
    \coordinate (PL) at (-0.80,1.30);
    \coordinate (PR) at (0.80,1.30);

    % Rhombus bars and horizontal rope BD
    \draw (A) -- (B) -- (C) -- (D) -- cycle;
    \draw (B) -- (D);

    % Tension arrows along BD
    \draw[->] (B) -- (-0.65,0.20);
    \node[below] at (-0.65,0.18) {$T$};
    \draw[->] (D) -- (0.65,0.20);
    \node[below] at (0.65,0.18) {$T$};

    % Coordinate axes
    \draw[->] (A) -- (0,-2.95) node[right] {$y$};
    \draw[->] (PL) -- (2.15,1.30) node[below] {$x$};

    % Weight arrow W below C
    \draw[->] (0,-2.00) -- (0,-2.65);
    \node[right] at (0,-2.48) {$W$};

    % Pegs (small circles tangent beneath AB and AD)
    \draw[fill=white] (-0.75,1.22) circle (0.065);
    \draw[fill=white] (0.75,1.22) circle (0.065);

    % Dimension a
    \draw[line width=0.45pt] (0.75,1.15) -- (0.75,0.72);
    \draw[<->] (0,0.88) -- (0.75,0.88) node[midway, above] {$a$};

    % Dimension l parallel to AB
    \coordinate (nAB) at (-0.809,0.588); % outward unit normal to AB
    \coordinate (Aoff) at ($(A)+0.55*(nAB)$);
    \coordinate (Boff) at ($(B)+0.55*(nAB)$);
    \draw[line width=0.45pt] (A) -- ($(A)+0.67*(nAB)$);
    \draw[line width=0.45pt] (B) -- ($(B)+0.67*(nAB)$);
    \draw[<->] (Aoff) -- (Boff) node[midway, above left, inner sep=2pt] {$l$};

    % Hinge dots
    \foreach \p in {A,B,C,D} {
        \fill (\p) circle (0.055);
    }

    % Labels
    \node[above] at (A) {$A$};
    \node[left] at (B) {$B$};
    \node[below left, inner sep=2pt] at (C) {$C$};
    \node[right] at (D) {$D$};
    \node[below left, inner sep=2pt] at (O) {$O$};
    \node at (-0.23,1.72) {$\alpha$};
\end{tikzpicture}

\end{center}
\end{problem}
\clearpage


\begin{problem}{9.1.22}
线密度为 $\sigma$ 、长度为 $l$ 的均质链条两端分别连接重量分别为 $P$ 和 $Q$ 的两个

球, 放在半径为 R 的光滑半圆柱面上, 平衡时链条和两球处在与圆柱轴线垂直的铅直面上, 试用虚功原理求平衡位置.
\end{problem}
\clearpage


\begin{problem}{9.1.23}
两球 A 和 B 分别重 P 和 Q，用长为 2l 的轻杆连接后放在一光滑的、半径为 R 的球壳内，求平衡时杆与水平线的交角 $\alpha$ 和 A、B 处的约束力.
\begin{center}
\begin{tikzpicture}[>=Stealth, line width=0.6pt, font=\small]
    \def\R{2.20}
    \coordinate (O) at (0,0);
    \coordinate (A) at (218:\R);
    \coordinate (B) at (346:\R);
    \coordinate (M) at ($(A)!0.5!(B)$);

    % Spherical bowl arc
    \draw (173:\R) arc (173:367:\R);

    % Axes
    \draw[->] (O) -- (3.10,0) node[below] {$x$};
    \draw[->] (O) -- (0,-2.75) node[right] {$y$};

    % Rod AB
    \draw (A) -- (B);

    % Radial dashed lines and normal forces
    \draw[dashed] (O) -- (A);
    \draw[->] (A) -- ($(A)!0.42!(O)$) node[above left, inner sep=2pt] {$N_A$};
    \draw[dashed] (O) -- (B);
    \draw[->] (B) -- ($(B)!0.42!(O)$);
    \node[above, inner sep=2pt] at ($(B)!0.26!(O)+(0.04,0.06)$) {$N_B$};

    % Dashed perpendicular from O to AB and dashed horizontal from A
    \draw[dashed] (O) -- (M);
    \draw[dashed] (A) -- ++(1.15,0) coordinate (AH);

    % Gravity arrows P and Q
    \draw[->] (A) -- ++(0,-0.92);
    \node[right] at ($(A)+(0,-0.68)$) {$P$};
    \draw[->] (B) -- ++(0,-0.88) node[midway, right] {$Q$};

    % Angle arcs
    \pic[draw, angle radius=0.80cm, "$\alpha$", angle eccentricity=1.30] {angle = AH--A--B};
    \pic[draw, angle radius=0.72cm, "$\beta$", angle eccentricity=1.28] {angle = O--B--A};

    % Dots and labels
    \fill (A) circle (0.055);
    \fill (B) circle (0.055);
    \node[above] at (O) {$O$};
    \node[left] at (A) {$A$};
    \node[right] at (B) {$B$};
\end{tikzpicture}

\end{center}
\end{problem}
\clearpage


\begin{problem}{9.1.25}
图 9.21 所示的是犁的后轮机构. 直杆 AB 和曲杆 $CO_{2}D$ 通过铰链 B 和 C 分别与连杆 BC 相连接, 连杆与两杆的夹角分别为 $\varphi_{1}$ 和 $\varphi_{2}$ , 当机构在图示位置时, BC 是水平的, $O_{1}$ 、A 两点的铅直距离为 h, $O_{2}$ 、D 两点的水平距离为 H, $O_{1}$ 、 $B = r_{1}$ , $O_{2}C = r_{2}$ , 在 A 点有一水平力 $F_{A}$ 作用, 在 D 点有一铅直力 $F_{D}$ 作用, 此时系统处于平衡, 求此两力的比值. 解 考虑在平衡位置附近有一虚位移, 如图 9.22 所示, BC 杆与水平线的夹角为 $\psi$ , 自固定点 $O_{1}$ 、 $O_{2}$ 分别取图示的 x 轴 (水平向左) 和 y 轴 (竖直向上). 与主动力做虚功有关的坐标为 $x_{A}$ 和 $y_{D}$ ,

\[
\begin{array}{l} x _ {A} = - O _ {1} A \cos (\varphi_ {1} - \psi) \\ y _ {D} = - O _ {2} D \sin \theta \end{array}
\]

\[
\begin{array}{l} \delta x _ {A} = O _ {1} A \sin (\varphi_ {1} - \psi) \delta (\varphi_ {1} - \psi) \\ \delta y _ {D} = - O _ {2} D \cos \theta \delta \theta \end{array}
\]

$O_{1}$ 、 $O_{2}$ 是固定点，两点的水平距离不变，

\[
- r _ {1} \mathrm{cos} (\varphi_ {1} - \psi) + B C \mathrm{cos} \psi + r _ {2} \mathrm{cos} (\varphi_ {2} + \psi) = \text {常量}
\]

曲杆的角度 $\angle CO_2D$ 不变，

\[
\theta + \pi - (\varphi_ {2} + \psi) = \text { 常   量 }
\]

对上述两个约束关系进行变分,可得虚位移间满足的关系

\[
r _ {1} \sin (\varphi_ {1} - \psi) \delta (\varphi_ {1} - \psi) - B C \sin \psi \delta \psi - r _ {2} \sin (\varphi_ {2} + \psi) \delta (\varphi_ {2} + \psi) = 0\tag{1}
\]

\[
\delta \theta - \delta (\varphi_ {2} + \psi) = 0\tag{2}
\]

在平衡位置, $\psi=0$ ,但 $\delta\psi\neq0,\delta(\varphi_{1}-\psi)\neq\delta\varphi_{1}$ ,由式(1),得到在平衡位置有

\[
r _ {1} \sin \varphi_ {1} \delta (\varphi_ {1} - \psi) - r _ {2} \sin \varphi_ {2} \delta (\varphi_ {2} + \psi) = 0\tag{3}
\]

\[
\delta x _ {A} = O _ {1} A \sin \varphi_ {1} \delta (\varphi_ {1} - \psi) = h \delta (\varphi_ {1} - \psi)
\]

\[
\delta y _ {D} = - O _ {2} D \cos \theta \delta \theta = - H \delta \theta = - H \delta (\varphi_ {2} + \psi)
\]

这里用了式(2).

在平衡位置,主动力做的虚功为零,

\[
\begin{array}{r l} & {\delta W = F _ {A} \delta x _ {A} + F _ {D} \delta y _ {D}} \\ & {\qquad = F _ {A} h \delta (\varphi_ {1} - \psi) - F _ {D} H \delta (\varphi_ {2} + \psi) = 0} \\ & {\frac {F _ {D}}{F _ {A}} = \frac {h \delta (\varphi_ {1} - \psi)}{H \delta (\varphi_ {2} + \psi)} = \frac {h r _ {2} \sin \varphi_ {2}}{H r _ {1} \sin \varphi_ {1}}} \end{array}
\]

这里用了式(3).
\begin{center}
\begin{tikzpicture}[>=Stealth, line width=0.6pt, font=\small]
    % Key coordinates
    \coordinate (B) at (0,0);
    \coordinate (C) at (2.80,0);
    \coordinate (O1) at (0.918,1.311);
    \coordinate (A) at (1.513,2.161);
    \coordinate (O2) at (3.632,-1.188);
    \coordinate (D) at (4.882,-1.523);

    % Mechanism bars
    \draw (B) -- (A);
    \draw (B) -- (C);
    \draw (C) -- (O2) -- (D);

    % Forces
    \draw[->] (A) -- ++(-0.78,0) node[left] {$F_A$};
    \draw[->] (D) -- ++(0,0.78) node[above, inner sep=2pt] {$F_D$};

    % Dashed reference lines for angles
    \draw[dashed] (C) -- ++(0.85,0) coordinate (CH);
    \draw[dashed] (O2) -- ++(0.95,0) coordinate (O2H);

    % Angle arcs
    \pic[draw, angle radius=0.48cm, "$\varphi_1$", angle eccentricity=1.42] {angle = C--B--O1};
    \pic[draw, angle radius=0.48cm, "$\varphi_2$", angle eccentricity=1.42] {angle = O2--C--CH};
    \pic[draw, angle radius=0.72cm, "$\theta$", angle eccentricity=1.26] {angle = D--O2--O2H};

    % Fixed support at O1 (tilted)
    \draw (O1) ++(0.04,-0.03) -- ++(0.06,-0.25) -- ++(0.20,0.14) -- cycle;
    \draw ($(O1)+(0.04,-0.32)$) -- ($(O1)+(0.36,-0.10)$);
    \foreach \t in {0.1,0.3,0.5,0.7,0.9} {
        \coordinate (P) at ($($(O1)+(0.04,-0.32)$)!\t!($(O1)+(0.36,-0.10)$)$);
        \draw[line width=0.45pt] (P) -- ++(-0.08,-0.10);
    }

    % Fixed support at O2 (horizontal)
    \draw (O2) -- ++(-0.14,-0.24) -- ++(0.28,0) -- cycle;
    \draw ($(O2)+(-0.24,-0.24)$) -- ($(O2)+(0.24,-0.24)$);
    \foreach \dx in {-0.20,-0.10,0,0.10,0.20} {
        \draw[line width=0.45pt] ($(O2)+(\dx,-0.24)$) -- ++(-0.08,-0.10);
    }

    % Hinge circles
    \foreach \p in {B,C,O2} {
        \draw[fill=white] (\p) circle (0.055);
    }
    \draw[fill=white] ($(O1)+(0.05,-0.04)$) circle (0.05);

    % Dimensions r1, r2, h, H
    \coordinate (n1) at (-0.819,0.574);
    \draw[line width=0.45pt] (B) -- ($(B)+0.60*(n1)$);
    \draw[line width=0.45pt] (O1) -- ($(O1)+0.60*(n1)$);
    \draw[<->] ($(B)+0.48*(n1)$) -- ($(O1)+0.48*(n1)$) node[midway, above left, inner sep=1pt] {$r_1$};

    \coordinate (n2) at (-0.819,-0.574);
    \draw[line width=0.45pt] (C) -- ($(C)+0.55*(n2)$);
    \draw[line width=0.45pt] (O2) -- ($(O2)+0.55*(n2)$);
    \draw[<->] ($(C)+0.44*(n2)$) -- ($(O2)+0.44*(n2)$) node[midway, below left, inner sep=1pt] {$r_2$};

    \draw[line width=0.45pt] (1.88,2.161) -- (2.22,2.161);
    \draw[line width=0.45pt] (1.88,1.311) -- (2.22,1.311);
    \draw[<->] (2.05,1.311) -- (2.05,2.161) node[midway, right] {$h$};

    \draw[line width=0.45pt] (3.632,-1.68) -- (3.632,-2.02);
    \draw[line width=0.45pt] (4.882,-1.60) -- (4.882,-2.02);
    \draw[<->] (3.632,-1.88) -- (4.882,-1.88) node[midway, below] {$H$};

    % Labels
    \node[right] at (A) {$A$};
    \node[below left, inner sep=2pt] at (B) {$B$};
    \node[above] at (C) {$C$};
    \node[right] at (D) {$D$};
    \node[left] at ($(O1)+(-0.08,0.06)$) {$O_1$};
    \node[above right, inner sep=2pt] at (O2) {$O_2$};
\end{tikzpicture}

\end{center}
\begin{center}
\begin{tikzpicture}[>=Stealth, line width=0.6pt, font=\small]
    % Key coordinates
    \coordinate (O1) at (1.50,1.50);
    \coordinate (B) at (0.439,0.439);
    \coordinate (A) at (2.207,2.207);
    \coordinate (O2) at (4.00,-1.20);
    \coordinate (C) at (3.638,-0.04);
    \coordinate (D) at (5.113,-1.650);

    % Coordinate axes x (left from O1) and y (up from O2)
    \draw[->] (O1) -- (-0.25,1.50) node[below] {$x$};
    \draw[->] (O2) -- (4.00,0.85) node[right] {$y$};

    % Mechanism bars
    \draw (B) -- (A);
    \draw (B) -- (C);
    \draw (C) -- (O2) -- (D);

    % Forces
    \draw[->] (A) -- ++(-0.80,0) node[left] {$F_A$};
    \draw[->] (D) -- ++(0,0.75) node[above] {$F_D$};

    % Dashed reference lines
    \draw[dashed] (B) -- ++(1.45,0) coordinate (BH);
    \draw[dashed] (C) -- ++(1.15,0) coordinate (CH);
    \draw[dashed] (C) -- ($(C)!-0.36!(B)$) coordinate (BCext);
    \draw[dashed] (O2) -- ++(0.95,0) coordinate (O2H);

    % Angle arcs
    \pic[draw, angle radius=0.58cm, "$\varphi_1$", angle eccentricity=1.36] {angle = C--B--A};
    \pic[draw, angle radius=1.12cm] {angle = C--B--BH};
    \node[right] at ($(B)+(1.16,-0.09)$) {$\psi$};
    \pic[draw, angle radius=0.95cm] {angle = BCext--C--CH};
    \node[right] at ($(C)+(0.98,-0.08)$) {$\psi$};
    \pic[draw, angle radius=0.50cm, "$\varphi_2$", angle eccentricity=1.42] {angle = O2--C--BCext};
    \pic[draw, angle radius=0.70cm, "$\theta$", angle eccentricity=1.26] {angle = D--O2--O2H};

    % Fixed support at O1
    \draw (O1) ++(0.05,-0.04) -- ++(0.06,-0.24) -- ++(0.20,0.14) -- cycle;
    \draw ($(O1)+(0.05,-0.31)$) -- ($(O1)+(0.36,-0.09)$);
    \foreach \t in {0.1,0.3,0.5,0.7,0.9} {
        \coordinate (P) at ($($(O1)+(0.05,-0.31)$)!\t!($(O1)+(0.36,-0.09)$)$);
        \draw[line width=0.45pt] (P) -- ++(-0.08,-0.10);
    }

    % Fixed support at O2
    \draw (O2) -- ++(-0.14,-0.24) -- ++(0.28,0) -- cycle;
    \draw ($(O2)+(-0.24,-0.24)$) -- ($(O2)+(0.24,-0.24)$);
    \foreach \dx in {-0.20,-0.10,0,0.10,0.20} {
        \draw[line width=0.45pt] ($(O2)+(\dx,-0.24)$) -- ++(-0.08,-0.10);
    }

    % Hinge circles
    \draw[fill=white] ($(O1)+(0.05,-0.04)$) circle (0.05);
    \draw[fill=white] (O2) circle (0.055);

    % Labels
    \node[right] at (A) {$A$};
    \node[below left, inner sep=2pt] at (B) {$B$};
    \node[above left, inner sep=2pt] at (C) {$C$};
    \node[right] at (D) {$D$};
    \node[above left, inner sep=2pt] at ($(O1)+(-0.04,0.04)$) {$O_1$};
    \node[left] at ($(O2)+(-0.06,0.02)$) {$O_2$};
    \node[above left, inner sep=1pt] at ($(B)!0.55!(O1)$) {$r_1$};
    \node[left] at ($(C)!0.52!(O2)$) {$r_2$};
\end{tikzpicture}

\end{center}
\end{problem}
\clearpage


\begin{problem}{9.1.26}
图 9.23 所示为一拔桩装置, 在木桩的上端 A 点上系一绳, 将绳的另一端固定在 C 点, 又在绳上一点 B 系另一绳, 此绳另一端固定在 E 点, 然后在绳的 D 点用力向下拉, 这时绳的 BD 段是水平的, AB 段是铅垂的, CB 段与铅垂线的夹角为 $\alpha_{1}$ , DE 段与水

平线的夹角为 $\alpha_{2}$ ，向下拉力为 F，求 AB 绳作用在桩上的拉力 T.
\begin{center}
\begin{tikzpicture}[>=Stealth, line width=0.6pt, font=\small]
    % Ground line and hatching
    \draw (-0.35,0) -- (4.55,0);
    \foreach \x in {-0.25, 0.00, ..., 4.45} {
        \draw[line width=0.45pt] (\x,0) -- (\x-0.16,-0.18);
    }

    % Gantry frame (double-walled inverted U)
    \draw (0.20,0) -- (0.20,3.30) -- (4.00,3.30) -- (4.00,0)
          (0.36,0) -- (0.36,3.14) -- (3.84,3.14) -- (3.84,0);

    % Wooden pile (masks ground/hatching)
    \draw[fill=white] (2.66,-0.32) rectangle (2.84,0.85);

    % Rope coordinates
    \coordinate (A) at (2.75,0.85);
    \coordinate (B) at (2.75,2.08);
    \coordinate (C) at (3.34,3.14);
    \coordinate (D) at (1.38,2.08);
    \coordinate (E) at (0.36,2.78);

    % Ropes
    \draw (A) -- (B) -- (C);
    \draw (D) -- (B);
    \draw (D) -- (E);

    % Dashed angle references
    \draw[dashed] (B) -- ++(0,0.68);
    \draw[dashed] (D) -- ++(-0.68,0);

    % Vertical downward arrows F and y
    \draw[->] (D) -- ++(0,-0.78) node[below] {$F$};
    \draw[->] (C) -- ++(0,-1.18) node[right] {$y$};

    % Labels
    \node[above right, inner sep=2pt] at (2.82,0.82) {$A$};
    \node[right] at (B) {$B$};
    \node[below right, inner sep=2pt] at (C) {$C$};
    \node[above right, inner sep=2pt] at (D) {$D$};
    \node[above right, inner sep=2pt] at (E) {$E$};
    \node at (2.96,2.74) {$\alpha_1$};
    \node at (0.86,2.24) {$\alpha_2$};
\end{tikzpicture}

\end{center}
\end{problem}
\clearpage


\begin{problem}{9.1.27}
重 $P_{1}$ 和 $P_{2}$ 的两物体连接在不可伸长的轻绳的两端，分放在倾角为 $\alpha, \beta$ 的斜面上，绳子绕过两个定滑轮和一个动滑轮，动滑轮的轴上挂一个重 $P$ 的物体，如图9.25所示。不计摩擦和滑轮质量，求平衡时 $P_{1}, P_{2}$ 的值。
\begin{center}
\begin{tikzpicture}[>=Stealth, line width=0.6pt, font=\small]
    \def\r{0.36}
    \coordinate (C1) at (-0.72,2.22);
    \coordinate (C2) at (0.72,2.22);
    \coordinate (C3) at (0,0.72);

    % Tangent points for inclined ropes (alpha = 42 deg, beta = 48 deg)
    \coordinate (T1) at ($(C1)+(132:\r)$);
    \coordinate (T2) at ($(C2)+(42:\r)$);

    % Left incline (42 deg) and rounded wedge
    \draw (-2.80,0.58) -- (-1.62,1.65);
    \draw (-2.55,0.81) -- (-0.78,0.81) -- (-0.78,1.62)
          .. controls (-0.78,1.86) and (-1.20,1.88) .. (-1.62,1.65);
    \draw (C1) -- (-1.05,1.60);

    % Right incline (48 deg) and rounded wedge
    \draw (2.58,0.41) -- (1.58,1.52);
    \draw (2.22,0.81) -- (0.78,0.81) -- (0.78,1.62)
          .. controls (0.78,1.86) and (1.18,1.86) .. (1.58,1.52);
    \draw (C2) -- (1.05,1.60);

    % Angle labels inside wedges
    \draw (-1.90,0.81) arc (0:42:0.65);
    \node at (-1.66,1.03) {$\alpha$};
    \draw (1.57,0.81) arc (180:132:0.65);
    \node at (1.35,1.03) {$\beta$};

    % Vertical ropes between pulleys
    \draw (-0.36,2.22) -- (-0.36,0.72);
    \draw (0.36,2.22) -- (0.36,0.72);

    % Block P1 on left incline
    \coordinate (B1) at ($(T1)+(222:0.84)$);
    \draw (T1) -- (B1);
    \begin{scope}[shift={(B1)}, rotate=42]
        \draw[fill=white] (-0.18,-0.18) rectangle (0.18,0.18);
    \end{scope}
    \draw[->] (B1) -- ++(0,-0.62);
    \node[right] at ($(B1)+(0.02,-0.50)$) {$P_1$};

    % Block P2 on right incline
    \coordinate (B2) at ($(T2)+(-48:0.84)$);
    \draw (T2) -- (B2);
    \begin{scope}[shift={(B2)}, rotate=-48]
        \draw[fill=white] (-0.18,-0.18) rectangle (0.18,0.18);
    \end{scope}
    \draw[->] (B2) -- ++(0,-0.62);
    \node[left] at ($(B2)+(-0.02,-0.50)$) {$P_2$};

    % Hanging block P on movable pulley C3
    \draw (C3) -- (0,-0.02);
    \draw[fill=white] (-0.22,-0.34) rectangle (0.22,-0.02);
    \draw[->] (0,-0.18) -- (0,-0.75) node[right] {$P$};

    % Pulleys (drawn last with white fill for clean stem/rope junctions)
    \draw[fill=white] (C1) circle (\r);
    \draw (C1) -- (-1.02,1.66);
    \draw[fill=white] (C2) circle (\r);
    \draw (C2) -- (1.02,1.66);
    \draw[fill=white] (C3) circle (\r);
    \draw (C3) -- (0,0.36);

    % Coordinate labels x1, x2, x3 and angular velocity phi dot
    \node[below left, inner sep=1pt] at (-2.76,0.58) {$x_1$};
    \node[below right, inner sep=1pt] at (2.54,0.42) {$x_2$};
    \draw[->] (0,2.65) -- (0,1.82) node[below] {$x_3$};
    \draw[->] (0.48,0.88) arc (20:-45:0.50);
    \node[right] at (0.48,0.58) {$\dot{\varphi}$};
\end{tikzpicture}

\end{center}
\end{problem}
\clearpage


\begin{problem}{9.1.28}
图 9.26 所示的装置中, 鼓轮 I 和 II 上分别有力矩 $M_{1}$ 、 $M_{2}$ 作用, 物体 A 重 P, 放在斜面上, 它与斜面间的摩擦因数为 $\mu$ , 斜面倾角为 $\alpha$ , 鼓轮的半径分别为 $r_{1}$ 和 $r_{2}$ , 动滑轮 B 重 Q. 试用虚功原理求平衡时 $M_{1}$ 、 $M_{2}$ 需满足的关系. 忽略所有轴承处的摩擦.
\begin{center}
\begin{tikzpicture}[>=Stealth, line width=0.6pt, font=\small]
    % Top fixed pulley
    \coordinate (Pt) at (1.50,4.05);
    \draw (Pt) circle (0.50);
    % Ceiling and hatching
    \draw (0.85,4.80) -- (2.15,4.80);
    \foreach \x in {0.95, 1.15, ..., 2.05} {
        \draw[line width=0.45pt] (\x,4.80) -- (\x+0.12,4.94);
    }
    % Top pulley hanger bracket
    \draw[fill=white] (1.44,3.99) rectangle (1.56,4.80);
    \fill (Pt) circle (0.025);

    % Drum I
    \coordinate (D1) at (2.50,2.85);
    \draw (D1) circle (0.50);
    \node at (2.50,2.52) {$\mathrm{I}$};
    % Drum I stand & floor
    \draw (2.22,2.44) -- (2.08,2.16);
    \draw (2.78,2.44) -- (2.92,2.16);
    \draw (1.95,2.16) -- (3.05,2.16);
    \foreach \x in {2.05, 2.20, ..., 3.00} {
        \draw[line width=0.45pt] (\x,2.16) -- (\x-0.12,2.02);
    }
    % Moment M1
    \draw[->] ($(D1)+(18:0.63)$) arc (18:88:0.63);
    \node[above right, inner sep=1pt] at ($(D1)+(50:0.65)$) {$M_1$};

    % Movable pulley B
    \coordinate (B) at (1.00,2.35);
    \draw (B) circle (0.50);
    \draw[fill=white] (0.94,1.73) rectangle (1.06,2.97);
    \fill (B) circle (0.025);
    \node[left] at (0.45,2.47) {$B$};

    % Drum II
    \coordinate (D2) at (2.00,0.95);
    \draw (D2) circle (0.50);
    \node at (2.00,0.62) {$\mathrm{II}$};
    % Drum II stand & floor
    \draw (1.72,0.54) -- (1.58,0.26);
    \draw (2.28,0.54) -- (2.42,0.26);
    \draw (1.45,0.26) -- (2.55,0.26);
    \foreach \x in {1.55, 1.70, ..., 2.50} {
        \draw[line width=0.45pt] (\x,0.26) -- (\x-0.12,0.12);
    }
    % Moment M2
    \draw[->] ($(D2)+(32:0.63)$) arc (32:102:0.63);
    \node[above right, inner sep=1pt] at ($(D2)+(60:0.64)$) {$M_2$};

    % Lower-left guide pulley
    \coordinate (Cg) at (0.12,1.40);
    \draw (Cg) circle (0.38);
    \draw[fill=white] ($(Cg)+(-0.04,-0.04)$) -- ++(-0.42,0.42) -- ++(0.08,0.08) -- ++(0.42,-0.42) -- cycle;
    \fill (Cg) circle (0.025);

    % Ropes
    \draw (2.00,4.05) -- (2.00,2.85);
    \draw (1.00,4.05) -- (1.00,2.97);
    \draw (1.50,2.35) -- (1.50,0.95);
    \draw (0.50,2.35) -- (0.50,1.40);

    % Incline (alpha = 22 deg) and ground
    \draw (-1.90,0) -- (-0.42,0);
    \foreach \x in {-1.80, -1.65, ..., -0.45} {
        \draw[line width=0.45pt] (\x,0) -- (\x-0.12,-0.14);
    }
    \draw (-1.90,0) -- (-0.25,0.667);
    \draw (-1.45,0) arc (0:22:0.45);
    \node at (-1.28,0.12) {$\alpha$};

    % Block A on incline and rope to guide pulley
    \coordinate (Tg) at ($(Cg)+(-68:0.38)$);
    \begin{scope}[shift={(-1.10,0.323)}, rotate=22]
        \draw (0,0) rectangle (0.44,0.32);
        \coordinate (Arope) at (0.44,0.16);
    \end{scope}
    \draw (Arope) -- (Tg);
    \node[above] at (-1.00,0.72) {$A$};
\end{tikzpicture}

\end{center}
\end{problem}
\clearpage


\begin{problem}{9.1.29}
均质长方形薄板长为 $l$ 、宽为 $b$ 、质量为 $m$ ，宽度为 $b$ 的两对边，一边靠在光滑的墙上，一边搁在固定的光滑的柱面上，如图 9.27 所示。如果要使板在任意位置处于平衡，柱面应取何种形状，并求柱面和墙对板的约束力。
\begin{center}
\begin{tikzpicture}[>=Stealth, scale=1.15]
  % Axes
  \draw[->] (0,0) -- (2.9,0) node[below left=1pt] {$x$};
  \draw[->] (0,0) -- (0,3.55) node[left=1pt] {$y$};
  \node[below left=1pt] at (0,0) {$O$};

  % Wall hatching along y-axis (left side)
  \foreach \y in {0.2,0.45,...,3.0} {
    \draw (0,\y) -- (-0.18,\y-0.18);
  }

  % Smooth curved surface y = 0.192*x^2 tangent to x-axis at O
  \draw[thick, domain=0:2.6, smooth, variable=\x] plot ({\x}, {0.192*\x*\x});
  % Hatching beneath the curved surface
  \foreach \x in {0.35,0.65,...,2.45} {
    \pgfmathsetmacro{\yc}{0.192*\x*\x}
    \draw (\x,\yc) -- (\x-0.16,\yc-0.24);
  }

  % Thin rectangular plate AB at angle theta = 30 deg from vertical
  % Inner-left top corner touches y-axis at P1=(0, 2.87)
  % Length L = 2.86, thickness w = 0.16
  \coordinate (P1) at (0, 2.87);
  \coordinate (P2) at ($(P1) + (2.86*0.5, -2.86*0.866)$);
  \coordinate (P3) at ($(P2) + (0.16*0.866, 0.16*0.5)$);
  \coordinate (P4) at ($(P1) + (0.16*0.866, 0.16*0.5)$);
  \draw[thick, fill=white] (P1) -- (P2) -- (P3) -- (P4) -- cycle;

  % Labels A, B, C, l, theta
  \node[above right=1pt] at (P4) {$A$};
  \node[above=3pt, xshift=4pt] at (P3) {$B$};
  \coordinate (Cpt) at ($0.5*(P1)+0.5*(P2)+0.5*(0.16*0.866,0.16*0.5)$);
  \fill (Cpt) circle (1.2pt);
  \node[below left=1pt] at ($(0.5*0+0.5*1.43, 0.5*2.87+0.5*0.39)$) {$C$};
  \node[right=4pt] at (Cpt) {$l$};
  \node at (0.28, 2.05) {$\theta$};
\end{tikzpicture}

\end{center}
\end{problem}
\clearpage


\begin{problem}{9.1.30}
半径为 $R$ 、鼓轮半径为 $r$ 的匀质轮轴, 重量为 $Q$ , 通过鼓轮上缠绕的绳子挂一重物 $P$ , 置于倾角为 $\alpha$ 的粗糙的斜面上, 如图 9.28 所示, 轮轴与斜面间的摩擦因数为 $\mu$ , 求轮轴处于平衡的条件.
\begin{center}
\begin{tikzpicture}[>=Stealth, scale=1.05]
  \def\ang{22}
  % Inclined plane sloping down to the left
  \coordinate (Lend) at (-2.7, {-2.7*tan(\ang)});
  \coordinate (Rend) at (2.2, {2.2*tan(\ang)});
  \draw[thick, ->] (Rend) -- ($(Lend) - ({0.6*cos(\ang)}, {0.6*sin(\ang)})$) node[below right=1pt] {$x$};
  \draw[thick, ->] (Rend) -- ++({-2.7*sin(\ang)}, {2.7*cos(\ang)}) node[right=2pt] {$y$};

  % Horizontal reference line and angle alpha
  \draw (-2.7, {-2.7*tan(\ang)}) -- (-1.0, {-2.7*tan(\ang)});
  \node at (-1.58, {-2.7*tan(\ang)+0.18}) {$\alpha$};

  % Hatching under inclined plane
  \foreach \s in {-2.1,-1.8,...,2.1} {
    \pgfmathsetmacro{\px}{\s}
    \pgfmathsetmacro{\py}{\s*tan(\ang)}
    \draw (\px,\py) -- (\px-0.14,\py-0.2);
  }

  % Wheel and axle resting on the incline at s = 0
  \def\R{1.15}
  \def\r{0.58}
  \coordinate (O) at ({-\R*sin(\ang)}, {\R*cos(\ang)});
  \draw[thick] (O) circle (\R);
  \draw[thick] (O) circle (\r);
  \fill (O) circle (1.3pt);

  % Radius arrows R and r
  \draw[->] (O) -- ++(155:\R) node[midway, below left=1pt] {$R$};
  \draw[->] (O) -- ++(25:\r) node[pos=0.55, below=2pt] {$r$};

  % Weight Q of wheel-axle
  \draw[->] (O) -- ++(0, -0.88) node[right=1pt] {$Q$};

  % Vertical rope from right side of inner drum down to weight P
  \coordinate (T) at ($(O) + (\r, 0)$);
  \draw[thick] (T) -- ++(0, -2.1);
  \draw[thick, fill=white] ($(T) + (-0.28, -2.1)$) rectangle ++(0.56, -0.38);
  \draw[->] ($(T) + (0, -2.29)$) -- ++(0, -0.65) node[right=1pt] {$P$};
\end{tikzpicture}

\end{center}
\end{problem}
\clearpage


\begin{problem}{9.1.31}
图 9.29 所示为一差动滑轮, 需用多大的力 F 才能将重物 P 吊起来.
\begin{center}
\begin{tikzpicture}[>=Stealth, scale=0.82]
  % Ceiling
  \draw[thick] (-1.6, 1.8) -- (1.6, 1.8);
  \foreach \x in {-1.45,-1.2,...,1.45} {
    \draw (\x, 1.8) -- (\x+0.18, 2.02);
  }

  % Upper compound pulley centered at (0,0)
  \def\rTwo{1.20}
  \def\rOne{0.62}
  \draw[thick] (0,0) circle (\rTwo);
  \draw[thick] (0,0) circle (\rOne);
  \fill (0,0) circle (1.3pt);

  % Upper hanger bracket
  \draw[thick, fill=white] (-0.13, 1.8) -- (-0.13, 0) arc (180:360:0.13) -- (0.13, 1.8) -- cycle;
  \fill (0,0) circle (1.3pt);

  % Radius arrows r1 and r2
  \draw[->] (0,0) -- (165:\rOne);
  \node at (205:0.35) {$r_1$};
  \draw[->] (0,0) -- (55:\rTwo) node[midway, right=2pt] {$r_2$};

  % Lower movable pulley centered at (0, -3.3)
  \def\rLow{0.90}
  \coordinate (L) at (0, -3.3);
  \draw[thick] (L) circle (\rLow);
  \fill (L) circle (1.3pt);

  % Left rope: external tangent between (0,0, rTwo) and (0,-3.3, rLow)
  % sin(phi) = (rTwo - rLow)/3.3 = 0.30/3.3 -> ~5.2 deg
  \draw[thick] (185.2:\rTwo) -- ($(L) + (185.2:\rLow)$);
  % Middle-right rope: external tangent between (0,0, rOne) and (0,-3.3, rLow)
  % sin(psi) = (rLow - rOne)/3.3 = 0.28/3.3 -> ~4.9 deg
  \draw[thick] (4.9:\rOne) -- ($(L) + (-4.9:\rLow)$);

  % Right loop with force F (exact semicircular loop at bottom-right)
  \coordinate (U1) at (38:\rTwo);
  \coordinate (U2) at ($(U1) + (-52:2.65)$);
  \coordinate (D1) at (218:\rOne);
  \coordinate (D2) at ($(D1) + (-52:2.65)$);
  \draw[thick] (U1) -- (U2) arc (38:-142:{0.5*(\rTwo+\rOne)}) -- (D1);
  % Force arrow F along upper branch
  \draw[thick, ->] (U1) -- ($(U1)!0.52!(U2)$) node[above right=1pt] {$F$};

  % Lower hanger bracket and weight P
  \draw[thick, fill=white] (-0.11, -3.3) arc (180:0:0.11) -- (0.11, -4.5) arc (360:180:0.11) -- cycle;
  \fill (L) circle (1.3pt);
  \draw[thick] (0, -4.61) -- (0, -5.05);
  \draw[thick, fill=white] (-0.35, -5.05) rectangle (0.35, -5.48);
  \draw[->] (0, -5.26) -- (0, -6.05) node[right=1pt] {$P$};
\end{tikzpicture}

\end{center}
\end{problem}
\clearpage


\begin{problem}{9.1.32}
半径为 r 的均质半圆柱体 A, 置于另一

个半径为 $R$ 的固定的半圆柱体 $B$ 的顶端, 如图9.30所示. 若 $A$ 在 $B$ 上做纯滚动, 问 $R$ 与 $r$ 满足什么关系时, 顶端这个平衡位置是稳定平衡位置.
\begin{center}
\begin{tikzpicture}[>=Stealth, scale=1.0]
  % Ground
  \draw[thick] (-2.3, 0) -- (2.3, 0);
  \foreach \x in {-2.15,-1.9,...,2.15} {
    \draw (\x, 0) -- (\x-0.16, -0.2);
  }

  % Lower semi-cylinder B
  \def\R{1.9}
  \draw[thick] (\R, 0) arc (0:180:\R);
  \draw[->] (0,0) -- (140:\R) node[midway, below left=1pt] {$R$};
  \node at (1.35, 0.55) {$B$};

  % Upper inverted semi-cylinder A
  \def\r{1.1}
  \coordinate (CA) at (0, \R+\r);
  \draw[thick] ($(CA) + (-\r, 0)$) -- ($(CA) + (\r, 0)$) arc (0:-180:\r) -- cycle;
  \draw[->] (CA) -- ++(-42:\r) node[midway, below left=1pt] {$r$};
  \node at (-0.55, \R+\r-0.38) {$A$};
\end{tikzpicture}

\end{center}
\end{problem}
\clearpage


\begin{problem}{9.1.33}
如图 9.33 所示, 均质的汽车库大门 AB 质量为 m、长 2r、A 边可沿光滑水平滑槽运动, B 边两端各连一个弹簧, 弹簧原长为 r-a, 即当 $\theta=\pi$ 时, 弹簧不受力的作用, B 边与曲柄 OB 用铰链连接, OB=r, 其质量可忽略不计, OB 可绕固定点 O 转动. 为了保证大门转到铅直位置 ( $\theta=0$ ) 时平稳地关闭, 弹簧的劲度系数 k 应多大?
\begin{center}
\begin{tikzpicture}[>=Stealth, line cap=round, line join=round, scale=0.85]
  \def\r{2.0}
  \def\a{0.75}
  \def\thetaang{55}

  % Coordinates: O at (0, \r), ground at y=0, top track at y=2\r
  \coordinate (O) at (0, \r);
  \coordinate (B) at ({\r*sin(\thetaang)}, {\r - \r*cos(\thetaang)});
  % Compute x_A such that |AB| = 2r and y_A = 2r
  \pgfmathsetmacro{\dy}{2*\r - (\r - \r*cos(\thetaang))}
  \pgfmathsetmacro{\dx}{sqrt((2*\r)^2 - \dy^2)}
  \pgfmathsetmacro{\xA}{\r*sin(\thetaang) - \dx}
  \coordinate (A) at (\xA, 2*\r);
  \coordinate (C) at ($0.5*(A) + 0.5*(B)$);
  \coordinate (S) at (0, \r + \a);

  % Top horizontal track and vertical wall (double lines)
  \draw[thick] (-2.6, 2*\r) -- (0.08, 2*\r);
  \draw[thick] (-2.6, 2*\r - 0.12) -- (-0.06, 2*\r - 0.12);
  \draw[thick] (-0.06, 2*\r - 0.12) -- (-0.06, 0);
  \draw[thick] (0.08, 2*\r) -- (0.08, 0);

  % Ground dashed line and dashed semicircle of radius r centered at O
  \draw[dashed] (-2.4, 0) -- (0, 0);
  \draw[dashed] (0, 0) arc (-90:90:\r);

  % Rod OB (thin double line) and pin at O, B
  \draw[thick] (O) -- (B) node[pos=0.65, below left=1pt] {$r$};
  \filldraw[fill=white, draw=black, thick] (O) circle (2.2pt) node[left=4pt] {$O$};

  % Door AB (thin double line)
  \draw[thick] (A) -- (B);
  \filldraw[fill=white, draw=black, thick] (A) circle (2.8pt) node[below left=2pt] {$A$};
  \filldraw[fill=white, draw=black, thick] (B) circle (2.8pt) node[below right=2pt] {$B$};
  \fill (C) circle (1.5pt) node[above right=1pt] {$C$};

  % Spring from S(0, r+a) to B
  \draw[thick, decorate, decoration={zigzag, segment length=6pt, amplitude=3.5pt, pre length=6pt, post length=8pt}]
    (S) -- (B) node[pos=0.45, above right=2pt] {$k$};

  % Angle theta at O
  \coordinate (Odown) at (0, 0.4);
  \pic[draw, angle radius=0.65cm, angle eccentricity=1.35, "$\theta$"] {angle = Odown--O--B};

  % Dimension 2r on the left
  \draw[<->] (-2.2, 0) -- (-2.2, 2*\r) node[midway, fill=white, inner sep=2pt] {$2r$};

  % Upward y axis
  \draw (-1.75, \r) -- (-1.45, \r);
  \draw[->] (-1.6, \r) -- (-1.6, 2*\r + 0.55) node[above] {$y$};

  % Dimension a
  \draw (-0.65, \r) -- (-0.15, \r);
  \draw (-0.65, \r + \a) -- (-0.15, \r + \a);
  \draw[->] (-0.45, \r - 0.45) -- (-0.45, \r);
  \draw[->] (-0.45, \r + \a + 0.45) -- (-0.45, \r + \a);
  \node at (-0.45, \r + 0.5*\a) {$a$};

  % Dimension 2r parallel to AB
  \coordinate (Aoff) at ($(A) + (0.6, 0.75)$);
  \coordinate (Boff) at ($(B) + (0.6, 0.75)$);
  \draw (A) -- ($(A) + (0.75, 0.94)$);
  \draw (B) -- ($(B) + (0.75, 0.94)$);
  \draw[<->] (Aoff) -- (Boff) node[midway, above right=1pt] {$2r$};
\end{tikzpicture}
\end{center}
\end{problem}
\clearpage


\begin{problem}{9.1.34}
二铰拱受集中负载 P 和 Q 作用, 各部分尺寸如图 9.34 所示. 求铰链 B 处约束力的水平分量.
\begin{center}
\begin{tikzpicture}[>=Stealth, scale=1.05]
  % Arch rib (center (0,0), radius 2a with a=1)
  \draw[thick] (2.06, 0) arc (0:180:2.06);
  \draw[thick] (1.94, 0) arc (0:180:1.94);

  % Left support A
  \draw[thick] (-2, 0) -- (-2.65, 0);
  \draw[thick] (-2, 0) -- (-2, -0.65);
  \draw[thick] (-2.65, 0) -- (-2.65, -0.65) -- (-2, -0.65);
  \foreach \y in {-0.1,-0.3,-0.5} {
    \draw (-2.65, \y) -- (-2.82, \y-0.17);
  }
  \foreach \x in {-2.6,-2.4,-2.2,-2.0} {
    \draw (\x, -0.65) -- (\x-0.17, -0.82);
  }

  % Right support B
  \draw[thick] (2, 0) -- (2.65, 0);
  \draw[thick] (2, 0) -- (2, -0.65);
  \draw[thick] (2.65, 0) -- (2.65, -0.65) -- (2, -0.65);
  \foreach \y in {-0.1,-0.3,-0.5} {
    \draw (2.65, \y) -- (2.82, \y-0.17);
  }
  \foreach \x in {2.0,2.2,2.4,2.6} {
    \draw (\x, -0.65) -- (\x+0.17, -0.82);
  }

  % Hinge circles
  \foreach \pt in {(-2,0), (-2.65,0), (-2,-0.65), (2,0), (2.65,0), (2,-0.65), (0,2)} {
    \draw[thick, fill=white] \pt circle (2.8pt);
  }
  \node[below right=2pt] at (-2, 0) {$A$};
  \node[below left=2pt] at (2, 0) {$B$};
  \node[below=4pt] at (0, 2) {$C$};

  % Forces P and Q
  \draw[thick, ->] (-1, 3.05) -- (-1, 1.80);
  \node[right=1pt] at (-1, 2.90) {$P$};
  \draw[thick, ->] (2.75, 1.0) -- (1.80, 1.0);
  \node[right=1pt] at (2.75, 1.0) {$Q$};

  % Horizontal dimension lines (top)
  \draw (-2, 0.25) -- (-2, 2.65);
  \draw (0, 2.05) -- (0, 2.65);
  \draw (2, 1.05) -- (2, 2.65);
  \draw[<->] (-2, 2.45) -- (-1, 2.45) node[midway, fill=white, inner sep=1.5pt] {$a$};
  \draw[<->] (-1, 2.45) -- (0, 2.45) node[midway, fill=white, inner sep=1.5pt] {$a$};
  \draw[<->] (0, 2.45) -- (2, 2.45) node[midway, fill=white, inner sep=1.5pt] {$2a$};

  % Vertical dimension lines (right)
  \draw (0.08, 2) -- (2.55, 2);
  \draw (2.10, 0) -- (2.55, 0);
  \draw[<->] (2.30, 1.0) -- (2.30, 2.0) node[midway, fill=white, inner sep=1.5pt] {$a$};
  \draw[<->] (2.30, 0.0) -- (2.30, 1.0) node[midway, fill=white, inner sep=1.5pt] {$a$};
\end{tikzpicture}

\end{center}
\end{problem}
\clearpage


\begin{problem}{9.1.35}
图 9.38 所示的平面桁架中各杆重量不计, 求 AC 杆和 BC 杆的内力.
\begin{center}
\begin{tikzpicture}[>=Stealth, scale=1.05]
  % Wall
  \draw[thick] (0, -1.6) -- (0, 1.6);
  \foreach \y in {-1.5,-1.3,...,1.5} {
    \draw (0, \y) -- (-0.2, \y-0.2);
  }

  % Coordinates (a = 2.0, a/2 = 1.0)
  \coordinate (A) at (0.45, 1.0);
  \coordinate (B) at (0.45, -1.0);
  \coordinate (C) at (2.45, 0);
  \coordinate (E) at (2.45, -1.0);
  \coordinate (D) at (4.45, 0);

  % Wall brackets
  \draw[thick] (0, 1.15) -- (A) -- (0, 0.85);
  \draw[thick] (0, -0.85) -- (B) -- (0, -1.15);

  % Truss members
  \draw[thick] (A) -- (C) -- (D) -- (E) -- (B) -- (C) -- (E);
  \draw[dashed] (0, 0) -- (C);

  % Hinges
  \foreach \p in {A,B,C,D,E} {
    \draw[thick, fill=white] (\p) circle (2.6pt);
  }

  % Labels
  \node[above right=1pt] at (A) {$A$};
  \node[above right=1pt] at (B) {$B$};
  \node[above=2pt] at (C) {$C$};
  \node[above right=1pt] at (D) {$D$};
  \node[below right=1pt] at (E) {$E$};

  % Force P at D
  \draw[thick, ->] (D) -- ++(0, -0.95) node[below=1pt] {$P$};

  % Vertical dimensions on left
  \draw (-0.80, 1.0) -- (-0.2, 1.0);
  \draw (-0.80, 0) -- (-0.2, 0);
  \draw (-0.80, -1.0) -- (-0.2, -1.0);
  \draw[<->] (-0.54, 0) -- (-0.54, 1.0) node[midway, fill=white, inner sep=1.5pt] {$\frac{a}{2}$};
  \draw[<->] (-0.54, -1.0) -- (-0.54, 0) node[midway, fill=white, inner sep=1.5pt] {$\frac{a}{2}$};

  % Horizontal dimensions at bottom (extension lines clear of E and P)
  \draw (0.45, -1.25) -- (0.45, -1.85);
  \draw (2.45, -1.35) -- (2.45, -1.85);
  \draw (4.45, -1.42) -- (4.45, -1.85);
  \draw[<->] (0.45, -1.65) -- (2.45, -1.65) node[midway, fill=white, inner sep=1.5pt] {$a$};
  \draw[<->] (2.45, -1.65) -- (4.45, -1.65) node[midway, fill=white, inner sep=1.5pt] {$a$};
\end{tikzpicture}

\end{center}
\end{problem}
\clearpage


\section{达朗贝尔原理、达朗贝尔——拉格朗日方程}

\begin{problem}{9.2.1}
用达朗贝尔原理解 6.3.10 题.
\end{problem}
\clearpage


\begin{problem}{9.2.2}
用达朗贝尔原理解 7.4.2 题.
\end{problem}
\clearpage


\begin{problem}{9.2.3}
用达朗贝尔原理解 7.4.40 题.
\end{problem}
\clearpage


\begin{problem}{9.2.4}
均匀细杆 $AB$ 可绕通过其杆上的固定点 $O$ 的水平轴自由转动, 此水平轴又绕通过 $O$ 点的竖直轴做恒定角速度为 $\omega$ 的转动, 如 $OA = a, OB = b$ , 试用达朗贝尔原理求 $AB$ 杆绕竖直轴做稳定转动时它与竖直轴所成的角 $\alpha$ .
\end{problem}
\clearpage


\begin{problem}{9.2.5}
重量为 $P$ 的均质圆柱沿与水平面夹角为 $\alpha$ 的轻悬臂梁做无滑动的滚动如图9.4.3所示.用达朗贝尔原理求悬臂梁在 $O$ 点受到的约束力.
\begin{center}
\begin{tikzpicture}[>=Stealth, scale=1.0]
  % Ceiling
  \draw[thick] (-0.9, 0) -- (0.9, 0);
  \foreach \x in {-0.75,-0.5,...,0.75} {
    \draw (\x, 0) -- (\x+0.16, 0.2);
  }
  \node[below left=1pt] at (0, 0) {$A$};

  % Vertical line Ax
  \draw[thick] (0, 0) -- (0, -4.1) node[left=2pt] {$x$};

  % Spring AB at angle 28 deg from vertical
  \coordinate (B) at (1.05, -1.95);
  \draw[thick, decorate, decoration={coil, aspect=0.45, segment length=7pt, amplitude=4pt, pre length=8pt, post length=8pt}] (0,0) -- (B);
  \node[above right=2pt] at (0.52, -0.92) {$r$};
  \node at (0.28, -1.15) {$\theta$};
  \node[above right=1pt] at (B) {$B$};

  % Vertical dashed line from B and rod BC
  \draw[dashed] (B) -- ++(0, -1.75);
  \coordinate (C) at (2.45, -3.95);
  \draw[thick] (B) -- (C) node[right=1pt] {$C$};
  \coordinate (D) at ($(B)!0.58!(C)$);
  \fill (D) circle (1.5pt);
  \node[above right=1pt] at (D) {$D$};
  \node at (1.32, -2.95) {$\varphi$};
\end{tikzpicture}

\end{center}
\end{problem}
\clearpage


\begin{problem}{9.2.6}
一个质量为 $m$ 的质点通过一根无重量的不可伸长的绳子系在半径为 $R$ 、竖直放置的固定圆柱边缘，绳子紧缠在圆柱上，质点与圆柱相接触。给质点一个沿径向方向的水平冲量，使质点获得一个初速度 $v_0$ ，质点被约束在与圆柱垂直的光滑水平面内运动。用达朗贝尔原理求以后未缠在圆柱体上的绳子长度 $l$ 与时间 $t$ 的关系。
\begin{center}
\begin{tikzpicture}[>=Stealth, scale=1.0]
  \def\R{1.80}
  \coordinate (O) at (0,0);
  \draw[thick] (O) circle (\R);

  % Top polar axis (t=0)
  \draw[thick] (0, \R) -- (-2.1, \R);
  \node[above] at (-1.15, \R) {极轴（$t=0$ 时）};

  % Vertical radius R
  \draw[thick] (O) -- (0, \R) node[midway, left=1pt] {$R$};

  % Point O' at -35 deg (psi = 125 deg from top vertical)
  \coordinate (Op) at (-35:\R);
  \draw[thick] (O) -- (Op);
  % Inward radial acceleration arrow R\dot{\psi}^2 along O'O
  \draw[thick, ->] (Op) -- ($(Op)!0.48!(O)$);
  \node[below left=1pt] at ($(Op)!0.70!(O)$) {$R\dot{\psi}^2$};

  % Polar axis at O'
  \draw[thick] (Op) -- ++(-2.1, 0);
  \node[below=2pt] at ($(Op) + (-1.45, 0)$) {极轴};

  % Angle psi
  \draw (90:0.42) arc (90:-35:0.42);
  \node at (32:0.65) {$\psi$};

  % Angle theta between leftward polar axis (180 deg) and tangent O'P (55 deg)
  \draw ($(Op) + (180:0.55)$) arc (180:55:0.55);
  \node at ($(Op) + (118:0.78)$) {$\theta$};
  \node[below right=1pt] at (Op) {$O'$};

  % Tangent line: downward arrow R\ddot{\psi} and string O'P of length l=2.55
  \draw[thick, ->] (Op) -- ++(-125:0.95) node[right=2pt] {$R\ddot{\psi}$};
  \coordinate (P) at ($(Op) + (55:2.55)$);
  \draw[thick] (Op) -- (P);
  \fill (P) circle (2pt);

  % Acceleration vectors at P
  \draw[thick, ->] (P) -- ++(55:0.75) node[above=1pt] {$\ddot{l}$};
  \draw[thick, ->] (P) -- ++(-125:0.85) node[right=2pt] {$l\dot{\theta}^2$};
  \draw[thick, ->] (P) -- ++(-35:0.68) node[above right=1pt] {$2\dot{l}\dot{\theta}$};
  \draw[thick, ->] ($(P) + (-125:0.18)$) -- ++(-35:0.68) node[below right=1pt] {$l\ddot{\theta}$};
\end{tikzpicture}

\end{center}
\end{problem}
\clearpage


\section{拉格朗日方程}

\begin{problem}{9.3.1}
一个质量为 $m$ 的质点在平面上运动, 不采用极坐标 $r, \varphi$ , 而用 $r$ 和 $\sin \varphi$ 为广义坐标, 写出此质点的动能.
\end{problem}
\clearpage


\begin{problem}{9.3.2}
在一根拉紧的弹性绳上等距离(距离为 a)地连着小球,两端固定,平衡时,弹性绳是水平的(重力可忽略不计),小球仅限于垂直方向运动,位移很小,即图 9.46 中 $y_{1}$ 、 $y_{2}$ 、 $\cdots$ 、 $y_{5}$ 均为小量,平衡时绳子张力为 $\tau$ ,写出系统的势能.
\begin{center}
\begin{tikzpicture}[>=Stealth, scale=1.0]
  \def\a{0.82}
  % Axes
  \draw[thick, ->] (0,0) -- (6*\a+0.45, 0) node[right=1pt] {$x$};
  \draw[thick, ->] (0,0) -- (0, 1.55) node[left=1pt] {$y$};
  \node[below left=1pt] at (0,0) {$O$};

  % Bead coordinates
  \coordinate (P0) at (0,0);
  \coordinate (P1) at (\a, 0.88);
  \coordinate (P2) at (2*\a, 0.48);
  \coordinate (P3) at (3*\a, 0.72);
  \coordinate (P4) at (4*\a, 0.46);
  \coordinate (P5) at (5*\a, 0.72);
  \coordinate (P6) at (6*\a, 0);

  % String segments
  \draw[thick] (P0) -- (P1) -- (P2) -- (P3) -- (P4) -- (P5) -- (P6);

  % Beads, dashed vertical lines, and labels
  \foreach \k/\yk/\lb in {1/0.88/y_1, 2/0.48/y_2, 3/0.72/y_3, 4/0.46/y_4} {
    \fill (\k*\a, \yk) circle (2.2pt);
    \node[above=2pt] at (\k*\a, \yk) {$\k$};
    \draw[dashed] (\k*\a, \yk) -- (\k*\a, 0);
    \node[right=1pt, font=\footnotesize] at (\k*\a, 0.45*\yk) {$\lb$};
  }
  \fill (5*\a, 0.72) circle (2.2pt);
  \node[above=2pt] at (5*\a, 0.72) {$5$};
  \draw[dashed] (5*\a, 0.72) -- (5*\a, 0);
  \node[right=1pt, font=\footnotesize] at (5*\a, 0.28*0.72) {$y_5$};

  % Interval ticks and dimension arrows
  \foreach \k in {0,1,2,3,4,5,6} {
    \draw (\k*\a, -0.06) -- (\k*\a, -0.44);
  }
  \foreach \k in {0,1,2,3,4,5} {
    \draw[<->] (\k*\a, -0.25) -- ({(\k+1)*\a}, -0.25) node[midway, fill=white, inner sep=1.5pt] {$a$};
  }
\end{tikzpicture}

\end{center}
\end{problem}
\clearpage


\begin{problem}{9.3.3}
劲度系数分别为 $k_{i}(i = 1,2,3,4)$ 的四根弹簧，原长分别为 $l_{i0}$ ，一端连在 $P$ 点，另一端分别连在固定点 $A,B,C,D$ . 取 $x,y$ 坐标，其原点是四个弹簧系统的平衡位置， $A,B,C,D$ 和 $P$ 点的坐标分别为 $(x_{1},y_{1}),(x_{2},y_{2}),(x_{3},y_{3}),(x_{4},y_{4})$ 和 $(x,y)$ ，平衡时，即 $P$ 位于原点时，各弹簧的长度分别为 $l_{i},\alpha_{i},\beta_{i}$ 分别是 $OA,OB,OC,OD$ 的方向余弦，即

\[
\alpha_ {i} = \frac {x _ {i}}{l _ {i}}, \quad \beta_ {i} = \frac {y _ {i}}{l _ {i}}
\]

$P$ 点偏离原点很小, 即 $x, y$ 均为小量. 证明弹簧系统的势能的近似表达式为

\[
V = \sum_ {i = 1} ^ {4} \left[ \frac {1}{2} k _ {i} (x ^ {2} + y ^ {2}) - \frac {k _ {i} l _ {i 0}}{2 l _ {i}} (\beta_ {i} x - \alpha_ {i} y) ^ {2} \right]
\]
\begin{center}
\begin{tikzpicture}[>=Stealth, scale=0.95]
  % Axes
  \draw[thick, ->] (-2.9, 0) -- (3.15, 0) node[below=2pt] {$x$};
  \draw[thick, ->] (0, -1.9) -- (0, 2.0) node[left=1pt] {$y$};
  \node[below left=1pt] at (0,0) {$O$};

  % Fixed anchor points A, B, C, D and displaced point P
  \coordinate (A) at (2.55, 1.55);
  \coordinate (B) at (-2.40, 1.10);
  \coordinate (C) at (-1.75, -1.35);
  \coordinate (D) at (1.60, -1.45);
  \coordinate (P) at (0.25, 0.52);

  % Equilibrium dashed lines
  \draw[dashed] (0,0) -- (A) node[pos=0.62, below right=1pt] {$l_1$};
  \draw[dashed] (0,0) -- (B) node[pos=0.48, below left=1pt] {$l_2$};
  \draw[dashed] (0,0) -- (C) node[pos=0.64, below right=2pt] {$l_3$};
  \draw[dashed] (0,0) -- (D) node[pos=0.68, below left=1pt] {$l_4$};

  % Springs
  \tikzset{spring/.style={thick, decorate, decoration={zigzag, segment length=6pt, amplitude=2.5pt, pre length=14pt, post length=14pt}}}
  \draw[spring] (P) -- (A) node[midway, above left=2pt] {$k_1$};
  \draw[spring] (P) -- (B) node[midway, above right=1pt] {$k_2$};
  \draw[thick, decorate, decoration={zigzag, segment length=6pt, amplitude=2.5pt, pre length=24pt, post length=12pt}] (P) -- (C) node[pos=0.66, above left=2pt] {$k_3$};
  \draw[spring] (P) -- (D) node[midway, above right=1pt] {$k_4$};

  % Point P label
  \fill (P) circle (1.5pt);
  \node[above=2pt] at (P) {$P$};
  \node[right=4pt, fill=white, inner sep=1pt] at (P) {$(x,y)$};

  % Hatched wall supports at A, B, C, D
  % Wall at A (tilt ~ -50 deg)
  \draw[thick] ($(A)+(-0.18,0.24)$) -- ($(A)+(0.18,-0.24)$);
  \foreach \t in {-0.15,-0.05,0.05,0.15} {
    \draw ($(A)+(\t,-\t*1.33)$) -- ++(0.16, 0.08);
  }
  \node[above right=2pt] at (A) {$A$};
  \node[below right=2pt] at (A) {$(x_1,y_1)$};

  % Wall at B (vertical)
  \draw[thick] ($(B)+(0,-0.24)$) -- ($(B)+(0,0.24)$);
  \foreach \y in {-0.2,-0.1,0,0.1,0.2} {
    \draw ($(B)+(0,\y)$) -- ++(-0.15, -0.15);
  }
  \node[left=5pt] at (B) {$B$};
  \node[below left=6pt] at (B) {$(x_2,y_2)$};

  % Wall at C (tilt ~ -30 deg)
  \draw[thick] ($(C)+(-0.22,0.12)$) -- ($(C)+(0.22,-0.12)$);
  \foreach \t in {-0.18,-0.06,0.06,0.18} {
    \draw ($(C)+(\t,-0.55*\t)$) -- ++(-0.08, -0.15);
  }
  \node[below left=4pt] at (C) {$C$};
  \node[below=14pt] at (C) {$(x_3,y_3)$};

  % Wall at D (tilt ~ 35 deg)
  \draw[thick] ($(D)+(-0.20,-0.14)$) -- ($(D)+(0.20,0.14)$);
  \foreach \t in {-0.15,-0.05,0.05,0.15} {
    \draw ($(D)+(\t,0.7*\t)$) -- ++(0.08, -0.15);
  }
  \node[below=5pt] at (D) {$D$};
  \node[below right=4pt] at (D) {$(x_4,y_4)$};
\end{tikzpicture}

\end{center}
\end{problem}
\clearpage


\begin{problem}{9.3.4}
双摆中质量为 $m_{1}$ 的质点与劲度系数分别为 $k_{1}$ 和 $k_{2}$ 、原长分别为 $l_{10}$ 和 $l_{20}$ 的两弹簧相连，弹簧的另一端分别固定于 $A$ 和 $B, m_{1}, m_{2}$ 两质点和两弹簧处在同一竖直平面内运动，当 $m_{1}$ 处于 $O$ 点时，且 $\varphi = 0$ ，双摆处于平衡位置，弹簧处于拉紧情况。求两弹簧的势能及其在 $\theta$ 很小时的近似表达式。
\begin{center}
\begin{tikzpicture}[>=Stealth, line width=0.8pt]
  \def\thetaAng{32}
  \def\rone{2.5}
  \def\rtwo{2.25}
  \def\phiAng{42}

  % Origin and axes
  \coordinate (O) at (0,0);
  \node[below left=1pt] at (O) {$O$};
  \draw[->] (O) -- (3.1,0) node[below] {$x$};
  \draw[->] (O) -- (0,3.4) node[left] {$y$};

  % Ceiling support
  \coordinate (P0) at (0,3.0);
  \fill[pattern=north east lines] (-0.65,3.0) rectangle (0.65,3.18);
  \draw (-0.65,3.0) -- (0.65,3.0);

  % Mass m1 and rod r1
  \coordinate (M1) at ({\rone*sin(\thetaAng)}, {3.0-\rone*cos(\thetaAng)});
  \draw (P0) -- node[midway, above right=1pt] {$r_1$} (M1);
  \fill (M1) circle (2.2pt) node[above right=2pt] {$m_1$};

  % Angle theta
  \coordinate (Ydown) at (0,1.5);
  \pic[draw, angle radius=0.85cm, angle eccentricity=1.28, "$\theta$"] {angle = Ydown--P0--M1};

  % Vertical dashed line from m1 and rod r2 to m2
  \coordinate (M1down) at ($(M1)+(0,-1.45)$);
  \draw[dashed] (M1) -- (M1down);
  \coordinate (M2) at ($(M1)+({\rtwo*sin(\phiAng)}, {-\rtwo*cos(\phiAng)})$);
  \draw (M1) -- (M2);
  \node[below left=1pt] at ($(M1)!0.78!(M2)$) {$r_2$};
  \fill (M2) circle (2.2pt) node[below right=1pt] {$m_2$};

  % Angle phi
  \pic[draw, angle radius=0.58cm, angle eccentricity=1.34, "$\varphi$"] {angle = M1down--M1--M2};

  % Spring k1 and wall A
  \coordinate (A) at (3.85,0.42);
  \draw[decorate, decoration={zigzag, segment length=2.4mm, amplitude=1.5mm, pre length=4.5mm, post length=4.5mm}]
    (M1) -- node[midway, above=5pt] {$k_1$} (A);
  \begin{scope}[shift={(A)}, rotate=-12]
    \fill[pattern=north east lines] (0,-0.28) rectangle (0.16,0.28);
    \draw (0,-0.28) -- (0,0.28);
    \node[above=5pt] at (-0.1,0.25) {$(x_1, y_1)$};
    \node[below=4pt] at (0,-0.25) {$A$};
  \end{scope}

  % Spring k2 and wall B (passes well below O)
  \coordinate (B) at (-1.35,-1.65);
  \draw[decorate, decoration={zigzag, segment length=2.4mm, amplitude=1.5mm, pre length=5mm, post length=5mm}]
    (M1) -- node[midway, below right=4pt] {$k_2$} (B);
  \begin{scope}[shift={(B)}, rotate=42]
    \fill[pattern=north east lines] (-0.28,-0.16) rectangle (0.28,0);
    \draw (-0.28,0) -- (0.28,0);
  \end{scope}
  \node[above left=3pt] at ($(B)+(0.08,0.08)$) {$(x_2, y_2)$};
  \node[below=5pt] at (B) {$B$};
\end{tikzpicture}
\end{center}
\end{problem}
\clearpage


\begin{problem}{9.3.5}
9.2.6 题的系统有几个自由度? 是否保守系? 选用适当的广义坐标, 由拉格朗日方程得到运动微分方程.
\end{problem}
\clearpage


\begin{problem}{9.3.6}
一个珠子套在一条光滑的锥形螺旋线上, 在重力的作用下运动. 用柱坐标, 锥形螺旋线方程为

\[
\rho = a z, \quad \varphi = - b z
\]

a、b 为常量，z 轴竖直向上. 证明：珠子的运动微分方程为

\[
\ddot {z} \left(a ^ {2} + 1 + a ^ {2} b ^ {2} z ^ {2}\right) + a ^ {2} b ^ {2} z \dot {z} = - g
\]
\end{problem}
\clearpage


\begin{problem}{9.3.7}
三个相互啮合的齿轮 $G_{1}, G_{2}, G_{3}$ 通过三个扭转弹簧与三个圆盘 $D_{4}, D_{5}, D_{6}$ 耦合，如图9.50所示。三个齿轮的半径分别为 $r_{1}, r_{2}, r_{3}$ ，围绕对称轴的转动惯量分别为 $I_{1}$ 、

$I_{2}$ 、 $I_{3}$ ，三个圆盘绕其对称轴的转动惯量分别为 $I_{4}$ 、 $I_{5}$ 、 $I_{6}$ ，它们偏离平衡位置的转角分别为 $\theta_{1}$ 、 $\theta_{2}$ 、 $\cdots$ 、 $\theta_{6}$ ，如图9.50所示。三个扭转弹簧扭转单位角度时受到的扭力矩大小分别为 $k_{1}$ 、 $k_{2}$ 、 $k_{3}$ ，求此系统的拉格朗日函数。

\[
T = \frac {1}{2} I _ {1} \dot {\theta} _ {1} ^ {2} + \frac {1}{2} I _ {2} \dot {\theta} _ {2} ^ {2} + \frac {1}{2} I _ {3} \dot {\theta} _ {3} ^ {2} + \frac {1}{2} I _ {4} \dot {\theta} _ {4} ^ {2} + \frac {1}{2} I _ {5} \dot {\theta} _ {5} ^ {2} + \frac {1}{2} I _ {6} \dot {\theta} _ {6} ^ {2}
\]

三个齿轮相互啮合,接触点无相对运动,

\[
\begin{array}{r l r} & {r _ {1} \dot {\theta} _ {1} = r _ {2} \dot {\theta} _ {2},} & {r _ {2} \dot {\theta} _ {2} = r _ {3} \dot {\theta} _ {3}} \\ & {\dot {\theta} _ {1} = \frac {r _ {2}}{r _ {1}} \dot {\theta} _ {2},} & {\dot {\theta} _ {3} = \frac {r _ {2}}{r _ {3}} \dot {\theta} _ {2}} \end{array}
\]

取平衡位置时，

\[
\begin{array}{r l} & {\theta_ {1} = \theta_ {2} = \theta_ {3} = \theta_ {4} = \theta_ {5} = \theta_ {6} = 0} \\ & {\theta_ {1} = \frac {r _ {2}}{r _ {1}} \theta_ {2}, \quad \theta_ {3} = \frac {r _ {2}}{r _ {3}} \theta_ {2}} \\ & {T = \frac {1}{2} I _ {1} \Big (\frac {r _ {2}}{r _ {1}} \Big) ^ {2} \dot {\theta} _ {2} ^ {2} + \frac {1}{2} I _ {2} \dot {\theta} _ {2} ^ {2} + \frac {1}{2} I _ {3} \Big (\frac {r _ {2}}{r _ {3}} \Big) ^ {2} \dot {\theta} _ {2} ^ {2} + \frac {1}{2} I _ {4} \dot {\theta} _ {4} ^ {2} + \frac {1}{2} I _ {5} \dot {\theta} _ {5} ^ {2} + \frac {1}{2} I _ {6} \dot {\theta} _ {6} ^ {2}} \\ & {\qquad = \frac {1}{2} r _ {2} ^ {2} \Big (\frac {I _ {1}}{r _ {1} ^ {2}} + \frac {I _ {2}}{r _ {2} ^ {2}} + \frac {I _ {3}}{r _ {3} ^ {2}} \Big) \dot {\theta} _ {2} ^ {2} + \frac {1}{2} I _ {4} \dot {\theta} _ {4} ^ {2} + \frac {1}{2} I _ {5} \dot {\theta} _ {5} ^ {2} + \frac {1}{2} I _ {6} \dot {\theta} _ {6} ^ {2}} \\ & {V = \frac {1}{2} k _ {1} (\theta_ {4} - \theta_ {2}) ^ {2} + \frac {1}{2} k _ {2} (\theta_ {1} - \theta_ {5}) ^ {2} + \frac {1}{2} k _ {3} (\theta_ {3} - \theta_ {6}) ^ {2}} \\ & {\qquad = \frac {1}{2} k _ {1} (\theta_ {4} - \theta_ {2}) ^ {2} + \frac {1}{2} k _ {2} \Big (\frac {r _ {2}}{r _ {1}} \theta_ {2} - \theta_ {5}) ^ {2} + \frac {1}{2} k _ {3} \Big (\frac {r _ {2}}{r _ {3}} \theta_ {2} - \theta_ {6} \Big) ^ {2}} \\ & {\qquad = \frac {1}{2} k _ {1} (\theta_ {4} - \theta_ {2}) ^ {2} + \frac {1}{2} k _ {2} r _ {2} ^ {2} \Big (\frac {\theta_ {2}}{r _ {1}} - \frac {\theta_ {5}}{r _ {2}} \Big) ^ {2} + \frac {1}{2} k _ {3} r _ {2} ^ {2} \Big (\frac {\theta_ {2}}{r _ {3}} - \frac {\theta_ {6}}{r _ {2}} \Big) ^ {2}} \\ & {\quad L = T - V = \frac 12 r _ {2} ^ {2} \Big (\frac I{r _ {1} ^ {2}} + \frac I{r _ {2} ^ {2}} + \frac I{r _ {3} ^ {2}} \Big) \dot {\theta} _ {2} ^ {2} + \frac 12 I _ {4} \dot {\theta} _ {4} ^ {2} + \frac 12 I _ {5} \dot {\theta} _ {5} ^ {2} + \frac 12 I _ {6} \dot {\theta} _ {6} ^ {2}} \\ & {-   \frac 12 k _ {1} (\theta_ {4} - \theta_ {2}) ^ {2} -   \frac 12 k _ {2} r _ {2} ^ {2}   \Big (\frac {\theta_ {2}}{r _ {1}} - \frac {\theta_ {5}}{r _ {2}}   \Big) ^ {2} -   \frac 12 k _ {3} r _ {2} ^ {2}   \Big (\frac {\theta_ {2}}{r _ {3}} - \frac {\theta_ {6}}{r _ {2}}   \Big) ^ {2}} \end{array}
\]
\begin{center}
\begin{tikzpicture}[>=Stealth, line width=0.8pt, xscale=0.40, yscale=0.72]
  % Meshing gears G1, G2, G3
  \draw (0,1.2) rectangle (0.9,2.4);
  \node at (0.45,1.52) {$G_1$};
  \node[right=1pt] at (0.9,2.38) {$I_1$};

  \draw (0,-0.6) rectangle (0.9,1.2);
  \node at (0.45,-0.32) {$G_2$};
  \node[right=1pt] at (0.9,-0.38) {$I_2$};

  \draw (0,-2.0) rectangle (0.9,-0.6);
  \node at (0.45,-1.68) {$G_3$};
  \node[right=1pt] at (0.9,-1.98) {$I_3$};

  % Diameter dimension annotations: 2r1 (left), 2r2 (right), 2r3 (left)
  \draw[thin] (-1.25,2.4) -- (-0.15,2.4);
  \draw[thin] (-0.42,1.8) -- (-0.15,1.8);
  \draw[thin] (-1.25,1.2) -- (-0.15,1.2);
  \draw[->] (-0.85,2.10) -- (-0.85,2.4);
  \draw[->] (-0.85,1.50) -- (-0.85,1.2);
  \node at (-0.85,1.8) {$2r_1$};

  \draw[thin] (1.65,1.2) -- (2.65,1.2);
  \draw[thin] (1.65,-0.6) -- (2.65,-0.6);
  \draw[->] (2.25,0.64) -- (2.25,1.2);
  \draw[->] (2.25,-0.04) -- (2.25,-0.6);
  \node at (2.25,0.3) {$2r_2$};

  \draw[thin] (-1.25,-0.6) -- (-0.15,-0.6);
  \draw[thin] (-0.42,-1.3) -- (-0.15,-1.3);
  \draw[thin] (-1.25,-2.0) -- (-0.15,-2.0);
  \draw[->] (-0.85,-0.96) -- (-0.85,-0.6);
  \draw[->] (-0.85,-1.64) -- (-0.85,-2.0);
  \node at (-0.85,-1.3) {$2r_3$};

  % Top branch: G1 -- k2 -- D5 (y = 1.8)
  \draw (0.9,1.72) rectangle (3.65,1.88);
  \draw (3.65,1.8) -- (3.95,1.8);
  \draw[decorate, decoration={zigzag, segment length=2.2mm, amplitude=1.4mm}] (3.95,1.8) -- node[midway, above=5pt] {$k_2$} (5.65,1.8);
  \draw (5.65,1.8) -- (6.0,1.8);
  \draw (6.0,1.72) rectangle (6.7,1.88);
  \draw (6.7,1.1) rectangle (7.12,2.5);
  \node[right=1pt] at (7.12,2.52) {$I_5$};
  \node[right=1pt] at (7.12,1.1) {$D_5$};
  \draw (7.12,1.72) rectangle (9.1,1.88);
  \draw (9.1,1.8) -- (9.55,1.8);
  % Rotation arrows theta1, theta5
  \draw[->] (3.2,1.42) arc[start angle=-70, end angle=-290, x radius=0.25, y radius=0.38];
  \node[above=1pt] at (3.2,2.2) {$\theta_1$};
  \draw[->] (8.65,1.42) arc[start angle=-70, end angle=-290, x radius=0.25, y radius=0.38];
  \node[above=1pt] at (8.65,2.2) {$\theta_5$};

  % Middle branch: D4 -- k1 -- G2 (y = 0.3)
  \draw (-1.75,0.22) rectangle (0,0.38);
  \draw (-2.05,0.3) -- (-1.75,0.3);
  \draw[decorate, decoration={zigzag, segment length=2.2mm, amplitude=1.4mm}] (-3.75,0.3) -- node[midway, above=5pt] {$k_1$} (-2.05,0.3);
  \draw (-4.05,0.3) -- (-3.75,0.3);
  \draw (-4.75,0.22) rectangle (-4.05,0.38);
  \draw (-5.18,-1.8) rectangle (-4.75,2.4);
  \node[right=1pt] at (-4.75,2.4) {$I_4$};
  \node[right=1pt] at (-4.75,-1.8) {$D_4$};
  \draw (-6.35,0.22) rectangle (-5.18,0.38);
  % Rotation arrows theta2, theta4
  \draw[->] (-1.2,0.68) arc[start angle=70, end angle=290, x radius=0.25, y radius=0.38];
  \node[above=1pt] at (-1.2,0.72) {$\theta_2$};
  \draw[->] (-5.8,0.68) arc[start angle=70, end angle=290, x radius=0.25, y radius=0.38];
  \node[above=1pt] at (-5.8,0.72) {$\theta_4$};

  % Bottom branch: G3 -- k3 -- D6 (y = -1.3)
  \draw (0.9,-1.38) rectangle (3.65,-1.22);
  \draw (3.65,-1.3) -- (3.95,-1.3);
  \draw[decorate, decoration={zigzag, segment length=2.2mm, amplitude=1.4mm}] (3.95,-1.3) -- node[midway, above=5pt] {$k_3$} (5.65,-1.3);
  \draw (5.65,-1.3) -- (6.0,-1.3);
  \draw (6.0,-1.38) rectangle (6.7,-1.22);
  \draw (6.7,-2.15) rectangle (7.12,-0.45);
  \node[right=1pt] at (7.12,-0.43) {$I_6$};
  \node[right=1pt] at (7.12,-2.15) {$D_6$};
  \draw (7.12,-1.38) rectangle (9.1,-1.22);
  \draw (9.1,-1.3) -- (9.55,-1.3);
  % Rotation arrows theta3, theta6
  \draw[->] (3.2,-1.68) arc[start angle=-70, end angle=-290, x radius=0.25, y radius=0.38];
  \node[above=1pt] at (3.2,-0.9) {$\theta_3$};
  \draw[->] (8.65,-0.92) arc[start angle=70, end angle=290, x radius=0.25, y radius=0.38];
  \node[above=1pt] at (8.65,-0.88) {$\theta_6$};
\end{tikzpicture}
\end{center}
\end{problem}
\clearpage


\begin{problem}{9.3.8}
如图 9.51, 一个质量为 $m$ 、半径为 $a$ 的均质薄圆盘的质心 $C$ 固连一根不计质

量的、长为 $b$ 的细杆 $AC$ , 细杆垂直于盘面. 此系统放在粗糙的、倾角为 $\alpha$ 的斜面上, 杆的 $A$ 端在斜面上的 $O$ 点保持不动, 圆盘沿半径为 $R = (a^2 + b^2)^{1/2}$ 的圆周在平衡位置附近做小振动(纯滚动). 选 $O$ 点为坐标系原点, $\zeta$ 轴垂直于斜面, $\eta$ 轴沿斜面向上, 固连于圆盘的动坐标系原点也在 $O$ 点, $AC$ 方向为 $z$ 轴. 求:

(1) 系统的动能；

(2) 系统的势能；

（3）证明欧拉角 $\varphi$ 的运动方程如同一个角频率为 $\Omega = \left(\frac{4gR\sin\alpha}{a^2 + 6b^2}\right)^{1 / 2}$ 的单摆.
\begin{center}
\begin{tikzpicture}[>=Stealth, line width=0.8pt, scale=0.95]
  \def\alphaAng{28}
  \def\bLen{2.65}
  \def\aLen{0.92}
  \pgfmathsetmacro{\betaAng}{atan(\aLen/\bLen)}
  \pgfmathsetmacro{\rodAng}{180+\alphaAng-\betaAng}

  % Horizontal base and inclined plane
  \coordinate (BaseV) at (0,0);
  \coordinate (BaseH) at (2.4,0);
  \coordinate (O) at (\alphaAng:4.1);
  \draw (BaseV) -- (BaseH);
  \draw[->] ($(\alphaAng:-0.55)$) -- ($(\alphaAng:5.1)$) node[above right] {$\eta$};
  \pic[draw, angle radius=0.72cm, angle eccentricity=1.32, "$\alpha$"] {angle = BaseH--BaseV--O};
  \node[below right=2pt] at (O) {$O$};

  % Zeta axis perpendicular to incline
  \coordinate (ZetaTip) at ($(O)+(\alphaAng+90:3.1)$);
  \draw[->, postaction={decorate}, decoration={markings, mark=at position 0.38 with {\arrow{Stealth}}}]
    (O) -- (ZetaTip) node[above left] {$\zeta$};
  \node[right=3pt] at ($(O)+(\alphaAng+90:1.25)$) {$\dot{\varphi}$};

  % Rod AC (double-line bar along rodAng)
  \coordinate (C) at ($(O)+(\rodAng:\bLen)$);
  \coordinate (ZTip) at ($(O)+(\rodAng:\bLen+1.35)$);
  \begin{scope}[shift={(O)}, rotate=\rodAng]
    \draw[fill=white] (0,-0.065) rectangle (\bLen,0.065);
  \end{scope}
  \node[above left=2pt] at ($(O)+(\rodAng:0.42)+(\rodAng-90:0.06)$) {$A$};
  \node[above right=2pt] at ($(C)+(\rodAng-90:0.06)$) {$C$};

  % z-axis along rod AC with psi-dot arrow
  \draw[->, postaction={decorate}, decoration={markings, mark=at position 0.42 with {\arrow{Stealth}}}]
    (O) -- (ZTip) node[below left] {$z$};
  \node[above=4pt] at ($(O)+(\rodAng:1.58)$) {$\dot{\psi}$};

  % Thin circular disk at C (perpendicular to rod)
  \begin{scope}[shift={(C)}, rotate=\rodAng]
    \draw[fill=white] (-0.065,-\aLen) rectangle (0.065,\aLen);
  \end{scope}

  % Angles theta and beta
  \pic[draw, angle radius=1.05cm, angle eccentricity=1.22, "$\theta$"] {angle = ZetaTip--O--C};
  \coordinate (IncDown) at ($(\alphaAng:1.5)$);
  \pic[draw, angle radius=0.85cm, angle eccentricity=1.32, "$\beta$"] {angle = C--O--IncDown};

  % Dimension a (upper radius of disk)
  \coordinate (DiskTop) at ($(C)+(\rodAng-90:\aLen)$);
  \coordinate (DimATop) at ($(DiskTop)+(\rodAng:0.55)$);
  \coordinate (DimACen) at ($(C)+(\rodAng:0.55)$);
  \draw[thin] ($(DiskTop)+(\rodAng:0.1)$) -- ($(DimATop)+(\rodAng:0.1)$);
  \draw[thin] ($(C)+(\rodAng:0.1)$) -- ($(DimACen)+(\rodAng:0.1)$);
  \draw[<->] (DimATop) -- node[midway, left=1pt] {$a$} (DimACen);

  % Dimension b (length OC below rod)
  \coordinate (DimBC) at ($(C)+(\rodAng+90:0.56)$);
  \coordinate (DimBO) at ($(O)+(\rodAng+90:0.56)$);
  \draw[thin] ($(C)+(\rodAng+90:0.08)$) -- ($(DimBC)+(\rodAng+90:0.1)$);
  \draw[thin] ($(O)+(\rodAng+90:0.08)$) -- ($(DimBO)+(\rodAng+90:0.1)$);
  \draw[<->] (DimBC) -- node[midway, fill=white, inner sep=1.5pt] {$b$} (DimBO);
\end{tikzpicture}
\end{center}
\end{problem}
\clearpage


\begin{problem}{9.3.9}
一个质量为 m 的质点受到势能为 $V(r)$ (r 是球坐标) 的力作用.

(1) 用一个绕 z 轴以角速度 $\omega$ 转动的球坐标系写出其拉格朗日函数；

(2) 证明上面写出的拉格朗日函数和用上述的转动参考系引入与速度有关的广义势 (它来自惯性离轴力和科里奥利力) 的拉格朗日函数是相同的.
\end{problem}
\clearpage


\begin{problem}{9.3.10}
一个由 $x, y, z$ 轴构成的直角坐标系，相对于惯性系以匀角速 $\omega$ 绕 $z$ 轴转动。一个质量为 $m$ 的质点在势能 $V(x, y, z)$ 的力的作用下运动。试用 $x, y, z$ 为广义坐标，写出此质点的拉格朗日函数及拉格朗日方程，并证明这些方程与在惯性参考系中这个质点受到一个力 $-\nabla V$ 及另一个可用广义势 $U$ 导出的力作用下的运动微分方程是一样的，写出这个广义势。
\end{problem}
\clearpage


\begin{problem}{9.3.11}
质量为 $m$ 、荷电为 $q$ 的粒子在柱坐标中有轴对称的电场 $E = \frac{E_0}{r} e_r$ 和均匀磁场 $B = B_0 k, E_0, B_0$ 均为常量。求此粒子的拉格朗日函数和运动微分方程。
\end{problem}
\clearpage


\begin{problem}{9.3.12}
当电磁场的矢势和标势受到规范变换：

$\bm{A}$ 换为

$\phi$ 换为

\[
\begin{array}{r l} \boldsymbol {A} ^ {\prime} & = \boldsymbol {A} + \nabla \Lambda \\ \phi^ {\prime} & = \phi - \frac {\partial \Lambda}{\partial t} \end{array}
\]

时, 它们所描述的电磁场不变, 可是拉格朗日函数显然要改变. 证明: 用新的拉格朗日函数仍然给出与原来相同的运动微分方程.
\end{problem}
\clearpage


\begin{problem}{9.3.13}
一个长度为 $b$ 、摆锤质量为 $m$ 的单摆悬于原长为 $a$ 、劲度系数为 $k$ 的弹簧上，悬点限于在竖直方向运动。摆锤处于黏性介质中，受到的阻力其大小与摆锤的速率成正比，比例系数为 $\gamma$ 。用摆的悬点在系统的平衡位置下的距离 $y$ 和摆与铅垂线的夹角 $\theta$ 为广义坐标，写出耗散函数和运动微分方程。
\begin{center}
\begin{tikzpicture}[>=Stealth, line width=0.8pt]
  % Ceiling
  \fill[pattern=north east lines] (-0.9,2.2) rectangle (0.9,2.38);
  \draw (-0.9,2.2) -- (0.9,2.2);

  % Vertical suspension and spring k
  \draw (0,2.2) -- (0,1.75);
  \draw[decorate, decoration={zigzag, segment length=2.3mm, amplitude=1.6mm}]
    (0,1.75) -- node[midway, right=5pt] {$k$} (0,0.65);
  \draw (0,0.65) -- (0,0);

  % Origin O and coordinate axes
  \coordinate (O) at (0,0);
  \node[left=2pt] at (O) {$O$};
  \draw[->] (O) -- (1.9,0) node[right] {$x$};
  \draw[->] (O) -- (0,-2.85) node[left] {$y$};

  % Equilibrium position label
  \draw[->] (1.15,0.88) node[above right=-1pt] {平衡位置} -- (0.08,0.06);

  % Pivot at distance y below O and pendulum rod b
  \coordinate (P) at (0,-1.05);
  \node[left=2pt] at (0,-0.55) {$y$};
  \coordinate (Bob) at ($(P)+(-58:1.85)$);
  \draw (P) -- node[midway, above right=1pt] {$b$} (Bob);
  \fill (Bob) circle (2.3pt);
  \node[right=3pt] at (Bob) {$m$};
  \node[below=5pt] at ($(Bob)+(0.15,0)$) {$(x, y)$};

  % Angle theta
  \coordinate (Ydown) at (0,-2.5);
  \pic[draw, angle radius=0.62cm, angle eccentricity=1.35, "$\theta$"] {angle = Ydown--P--Bob};
\end{tikzpicture}
\end{center}
\end{problem}
\clearpage


\begin{problem}{9.3.14}
求 9.1.9 题所述系统的耗散函数, 并由此求 $Q_{y}, Q_{y_{3}}$ .
\end{problem}
\clearpage


\begin{problem}{9.3.15}
求 9.1.12 题所述系统的耗散函数 $R(\dot{x}_1, \dot{q}_1, \dot{q}_2)$ , 并由此求 $Q_{x_1}, Q_{q_1}, Q_{q_2}$ .
\end{problem}
\clearpage


\begin{problem}{9.3.16}
求 9.1.19 题所述系统的耗散函数.
\end{problem}
\clearpage


\begin{problem}{9.3.17}
上题的圆环改为半径为 $r$ 的圆盘， $f$ 改为相接触的单位面积所受的摩擦力的大小，证明耗散函数为

\[
R = f \int_ {0} ^ {r} \int_ {0} ^ {2 \pi} \left\{\left[ \dot {x} - r ^ {\prime} \dot {\varphi} \sin (\varphi + \alpha) \right] ^ {2} + \left[ \dot {y} + r ^ {\prime} \dot {\varphi} \cos (\varphi + \alpha) \right] ^ {2} \right\} ^ {1 / 2} r ^ {\prime} d r ^ {\prime} d \alpha
\]
\end{problem}
\clearpage


\begin{problem}{9.3.18}
假定图 9.53 中的减震器引入的力与相对速度的立方成正比, 比例系数为 b, 各对接触平面间的阻力是正比于相对速度的黏性阻力, 比例系数为 $\alpha_{1}$ 、 $\alpha_{2}$ 、 $\alpha_{3}$ . 求系统的耗散函数.
\begin{center}
\begin{tikzpicture}[>=Stealth, line width=0.8pt, scale=0.60]
  % Left wall and bottom x-axis
  \draw[->] (0,0) -- (0,2.95);
  \draw[->] (0,0) -- (10.1,0) node[below] {$x$};

  % Plank m1 and its left stop
  \draw (1.75,0.06) rectangle (9.45,0.72);
  \draw (1.75,0.72) rectangle (1.98,2.05);
  \node at (6.25,0.38) {$m_1$};

  % Spring k1
  \draw (0,0.39) -- (0.25,0.39);
  \draw[decorate, decoration={zigzag, segment length=2.1mm, amplitude=1.4mm}] (0.25,0.39) -- node[midway, above=4pt] {$k_1$} (1.5,0.39);
  \draw (1.5,0.39) -- (1.75,0.39);

  % Blocks m2 and m3
  \draw (4.15,0.78) rectangle (4.9,2.85);
  \node at (4.525,1.82) {$m_2$};

  \draw (7.75,0.78) rectangle (8.5,2.85);
  \node at (8.125,1.82) {$m_3$};

  % Spring k2
  \draw (1.98,1.5) -- (2.25,1.5);
  \draw[decorate, decoration={zigzag, segment length=2.1mm, amplitude=1.4mm}] (2.25,1.5) -- node[midway, above=4pt] {$k_2$} (3.88,1.5);
  \draw (3.88,1.5) -- (4.15,1.5);

  % Spring k3
  \draw (4.9,1.2) -- (5.2,1.2);
  \draw[decorate, decoration={zigzag, segment length=2.1mm, amplitude=1.4mm}] (5.2,1.2) -- node[midway, above=4pt] {$k_3$} (7.45,1.2);
  \draw (7.45,1.2) -- (7.75,1.2);

  % Dashpot (damper) between m2 and m3
  \draw (4.9,2.32) -- (6.85,2.32);
  \draw (4.9,2.44) -- (6.85,2.44);
  \draw (6.85,2.06) -- (6.85,2.70);
  % Cylinder open on left
  \draw (6.55,2.82) -- (6.72,2.74) -- (7.75,2.74);
  \draw (6.55,1.94) -- (6.72,2.02) -- (7.75,2.02);
  \node[above=2pt] at (7.22,2.76) {减震器};
  % Dots inside dashpot cylinder
  \foreach \px/\py in {7.02/2.20, 7.08/2.56, 7.22/2.36, 7.30/2.14, 7.38/2.60, 7.52/2.26, 7.56/2.47, 7.16/2.11, 7.42/2.40, 7.60/2.62} {
    \fill (\px,\py) circle (1.2pt);
  }

  % Position labels x1, x2, x3
  \node[below=2pt] at (1.75,0) {$x_1$};
  \node[below=2pt] at (4.525,0) {$x_2$};
  \node[below=2pt] at (7.75,0) {$x_3$};

  % Friction / damping coefficients alpha1, alpha2, alpha3
  \draw[->] (2.55,0.28) to[bend right=25] (2.1,0.06);
  \node[right=1pt] at (2.52,0.32) {$\alpha_1$};

  \draw[->] (4.72,0.42) to[bend right=25] (4.52,0.75);
  \node[right=1pt] at (4.7,0.42) {$\alpha_2$};

  \draw[->] (8.32,0.42) to[bend right=25] (8.12,0.75);
  \node[right=1pt] at (8.3,0.42) {$\alpha_3$};
\end{tikzpicture}
\end{center}
\end{problem}
\clearpage


\begin{problem}{9.3.19}
如图 9.54, 一个质量为 $M$ 、半径为 $R$ 的圆盘在一个光滑的水平面上滑动, 另有一个质量为 $m$ 、半径为 $r$ 的圆盘, 将其中心钉在前一圆盘上离中心距离 $b$ 处, 它可在前一圆盘上围绕该点无摩擦地转动. 选取适当的广义坐标, 写出该系统的拉格朗日函数, 给出一切可以直接写出的运动积分.
\begin{center}
\begin{tikzpicture}[>=Stealth, line width=0.8pt, scale=0.95]
  % Coordinate axes
  \draw[->] (0,0) -- (4.95,0) node[below] {$x$};
  \draw[->] (0,0) -- (0,4.85) node[left] {$y$};

  % Large disk M of radius R
  \coordinate (C1) at (2.7,2.45);
  \def\R{2.05}
  \draw (C1) circle (\R);
  \node at (3.6,1.25) {$M$};

  % Horizontal reference line and radius R
  \coordinate (H) at ($(C1)+(1.65,0)$);
  \draw (C1) -- (H);
  \draw[->] (C1) -- node[midway, above left=1pt] {$R$} ($(C1)+(228:\R)$);

  % Small disk m of radius r at distance b and angle phi
  \def\b{1.18}
  \def\r{0.66}
  \coordinate (C2) at ($(C1)+(125:\b)$);
  \draw (C1) -- node[midway, below left=1pt] {$b$} (C2);
  \draw (C2) circle (\r);
  \fill (C2) circle (1.8pt) node[above=2pt] {$m$};
  \draw[->] (C2) -- node[midway, below right=0pt] {$r$} ($(C2)+(192:\r)$);

  % Angle phi
  \pic[draw, angle radius=0.45cm, angle eccentricity=1.45, "$\varphi$"] {angle = H--C1--C2};
\end{tikzpicture}
\end{center}
\end{problem}
\clearpage


\begin{problem}{9.3.20}
一个均匀的实心圆柱体, 半径为 $R$ , 质量为 $M$ , 放在一个水平面上. 另一个相同的圆柱体放在它上面, 和它的最高母线相切. 两圆柱体从静止状态开始做无滑动滚动. 写出该系统的拉格朗日函数, 说明存在哪些运动积分, 并证明, 只要两圆柱保持接触, 则有

\[
\dot {\theta} ^ {2} = \frac {1 2 g (1 - \cos \theta)}{R (1 7 + 4 \cos \theta - 4 \cos^ {2} \theta)}
\]

其中 $\theta$ 是两轴所决定的平面和通过一个轴的铅垂面间的夹角.
\begin{center}
\begin{tikzpicture}[>=Stealth, line width=0.8pt, scale=0.95]
  \def\R{1.15}
  \def\thetaAng{30}

  % Horizontal plane
  \draw[->] (-2.05,0) -- (2.75,0) node[below] {$x$};

  % Lower and upper identical cylinders of mass M, radius R
  \coordinate (C1) at (0,\R);
  \coordinate (C2) at ({2*\R*sin(\thetaAng)}, {\R+2*\R*cos(\thetaAng)});

  \draw (C1) circle (\R);
  \fill (C1) circle (1.8pt);
  \node at (-0.58,0.98) {$M$};

  \draw (C2) circle (\R);
  \fill (C2) circle (1.8pt);
  \node at (0.72,3.82) {$M$};

  % Vertical reference line and center-to-center line
  \coordinate (Vtop) at (0,4.35);
  \draw (C1) -- (Vtop);
  \draw (C1) -- (C2);

  % Angle theta cleanly inside sector
  \pic[draw, angle radius=0.92cm, angle eccentricity=0.74, "$\theta$"] {angle = C2--C1--Vtop};

  % Rotation arrows phi and psi
  \draw[->] ($(C1)+(10:\R+0.28)$) arc[start angle=10, end angle=-35, radius=\R+0.28];
  \node[right=1pt] at ($(C1)+(-14:\R+0.32)$) {$\varphi$};

  \draw[->] ($(C2)+(15:\R+0.28)$) arc[start angle=15, end angle=-30, radius=\R+0.28];
  \node[right=1pt] at ($(C2)+(-8:\R+0.32)$) {$\psi$};
\end{tikzpicture}
\end{center}
\end{problem}
\clearpage


\begin{problem}{9.3.21}
一根质量为 $M$ 、长为 $2a$ 的均质棒 $OA$ 装在固定点的光滑枢轴上，另有一个质量为 $m$ 的珠子 $P$ 用一根自然长度为 $a$ 、劲度系数为 $k$ 的轻弹簧连于 $O$ 点，并套在棒上，可沿棒做无摩擦滑动。设 $\theta$ 是 $OA$ 与铅垂线间的夹角， $\varphi$ 是包含 $OA$ 的铅垂面与固定铅垂面间的夹角， $r$ 为 $OP$ 间的距离。用 $r, \theta, \varphi$ 为广义坐标，写出系统的拉格朗日函数，并给出两个运动常数。
\begin{center}
\begin{tikzpicture}[>=Stealth, line width=0.8pt]
  \def\thetaAng{34}
  \def\rodLen{3.25}
  \def\rLen{2.25}
  \pgfmathsetmacro{\rodAng}{-90+\thetaAng}

  % Vertical z-axis and horizontal x-axis at O
  \coordinate (O) at (0,0);
  \draw[->] (0,-3.05) -- (0,1.45) node[left] {$z$};
  \draw[->] (O) -- (2.35,0) node[below] {$x$};
  \node[above left=1pt] at (O) {$O$};

  % Azimuthal rotation arrow phi around z-axis
  \draw[->] (-0.42,0.88) arc[start angle=150, end angle=385, x radius=0.48, y radius=0.15];
  \node[right=2pt] at (0.5,0.86) {$\varphi$};

  % Rod OA (thin bar)
  \coordinate (A) at (\rodAng:\rodLen);
  \coordinate (P) at (\rodAng:\rLen);
  \begin{scope}[rotate=\rodAng]
    \draw[fill=white] (0,-0.06) rectangle (\rodLen,0.06);
  \end{scope}
  \node[below right=2pt] at (A) {$A$};
  \node[above right=2pt] at (\rodAng:2.75) {$M$};

  % Coiled spring around rod from O to P
  \draw[decorate, decoration={coil, aspect=0.45, segment length=2.5mm, amplitude=1.6mm, pre length=6.8mm, post length=1.5mm}]
    (O) -- node[pos=0.6, above right=4pt] {$k$} (P);

  % Bead P of mass m
  \draw[fill=white] (P) circle (0.14);
  \draw (P) circle (0.14);
  \node[above right=3pt] at (P) {$P$};
  \node[below left=2pt] at (P) {$m$};

  % Angle theta cleanly inside sector
  \coordinate (Zdown) at (0,-1.8);
  \pic[draw, angle radius=0.82cm, angle eccentricity=0.74, "$\theta$"] {angle = Zdown--O--A};

  % Dimension lines for 2a and r
  \coordinate (ExtO) at (\rodAng-90:1.25);
  \coordinate (ExtP) at ($(P)+(\rodAng-90:1.25)$);
  \coordinate (ExtA) at ($(A)+(\rodAng-90:0.94)$);
  \draw[thin] (O) -- (ExtO);
  \draw[thin] ($(P)+(\rodAng-90:0.18)$) -- (ExtP);
  \draw[thin] ($(A)+(\rodAng-90:0.08)$) -- (ExtA);

  \draw[<->] (\rodAng-90:0.82) -- node[midway, fill=white, inner sep=1.2pt] {$2a$} ($(A)+(\rodAng-90:0.82)$);
  \draw[<->] (\rodAng-90:1.14) -- node[midway, fill=white, inner sep=1.2pt] {$r$} ($(P)+(\rodAng-90:1.14)$);
\end{tikzpicture}
\end{center}
\end{problem}
\clearpage


\begin{problem}{9.3.22}
一个质量为 $m$ 的粒子在一光滑平面上运动, 它受到平面上一固定点 $P$ 的吸引力, 力的大小与离 $P$ 点的距离的平方成反比, 比例系数为 $k$ . 用极坐标写出该粒子的拉格朗日函数和运动微分方程, 并给出两个运动积分.
\end{problem}
\clearpage


\begin{problem}{9.3.23}
质量为 $m_{1}, m_{2}$ 的两粒子之间的作用力为有心力, 势能为 $V(r), r$ 为相对位矢.

（1）在质心平动参考系中写出此系统的拉格朗日函数，并由此证明能量和角动量守恒；证明运动总在一个平面内，满足开普勒第二定律(r 在相等的时间内扫过相等的面积)；

(2) 若势能 $V = \frac{1}{2} kr^2 (k$ 是一个正的常数), 且已知总能量 $E$ 和角动量 $J$ , 求出运动过程中 $r$ 的最大值和最小值满足的表达式.
\begin{center}
\begin{tikzpicture}[>=Stealth, line width=0.8pt]
  \def\phiAng{36}

  % Coordinate axes
  \coordinate (O) at (0,0);
  \draw[->] (-2.3,0) -- (2.5,0) node[below] {$x$};
  \draw[->] (0,-1.65) -- (0,1.95) node[left] {$y$};
  \node[below right=1pt] at (O) {$O$};

  % Position vectors r2 and r1
  \coordinate (R2) at (\phiAng:2.25);
  \coordinate (R1) at (\phiAng+180:1.75);
  \draw[->] (O) -- node[midway, above left=1pt] {$\boldsymbol{r}_2$} (R2);
  \node[above right=1pt] at (R2) {$m_2$};

  \draw[->] (O) -- node[midway, above left=1pt] {$\boldsymbol{r}_1$} (R1);
  \node[below left=1pt] at (R1) {$m_1$};

  % Angle phi2
  \coordinate (Xpos) at (1.5,0);
  \pic[draw, angle radius=0.65cm, angle eccentricity=1.35, "$\varphi_2$"] {angle = Xpos--O--R2};
\end{tikzpicture}
\end{center}
\end{problem}
\clearpage


\begin{problem}{9.3.24}
给定一个一维运动的拉格朗日函数

\[
L = \mathrm{e} ^ {\gamma_ {t}} \left(\frac {1}{2} m \dot {q} ^ {2} - \frac {1}{2} k q ^ {2}\right)
\]

其中 $\gamma, m, k$ 均是正值常量. 写出拉格朗日方程, 说明有无可直接写出的运动常数, 并描述这一运动.

若采用广义坐标 $s$ ，它与原广义坐标 $q$ 的变换关系为 $s = \mathrm{e}^{\frac{1}{2}\gamma t} q$ ，重新写出拉格朗日函数、拉格朗日方程，回答有无可直接写出的运动常数，并说明两个解之间的关系。
\end{problem}
\clearpage


\begin{problem}{9.3.25}
一个质量为 m、半径为 a 的均质圆环在一个质量为 M、倾角为 $\alpha$ 、置于光滑水平面上的斜面上作无滑动滚动，运动是平面的。选择表示斜面对某固定参考点的水平

距离 $x$ 和表示圆环从斜面上某参考点沿斜面滚下的距离 $s$ 为广义坐标，找出四个运动常数.
\begin{center}
\begin{tikzpicture}[>=Stealth, line width=0.8pt, scale=0.95]
  % Horizontal floor
  \draw[->] (-0.4,0) -- (5.4,0) node[below] {$x$};

  % Inclined wedge M
  \coordinate (BL) at (0.2,0.06);
  \coordinate (BR) at (4.95,0.06);
  \coordinate (TL) at (0.2,2.15);
  \draw (BL) -- (BR) -- (TL) -- cycle;
  \node at (1.55,0.68) {$M$};

  % Angle alpha
  \pic[draw, angle radius=0.85cm, angle eccentricity=1.28, "$\alpha$"] {angle = TL--BR--BL};

  % Coordinate s along incline
  \pgfmathsetmacro{\dx}{4.95-0.2}
  \pgfmathsetmacro{\dy}{0.06-2.15}
  \pgfmathsetmacro{\slen}{veclen(\dx,\dy)}
  \pgfmathsetmacro{\ux}{\dx/\slen}
  \pgfmathsetmacro{\uy}{\dy/\slen}
  % Upward normal unit vector to slope: (-\uy, \ux)
  \pgfmathsetmacro{\nx}{-\uy}
  \pgfmathsetmacro{\ny}{\ux}

  \coordinate (S1) at ($(TL)+0.1*(\dx,\dy)-0.13*(\nx,\ny)$);
  \coordinate (S2) at ($(TL)+0.64*(\dx,\dy)-0.13*(\nx,\ny)$);
  \draw[->] (S1) -- (S2) node[below left=1pt] {$s$};

  % Rolling hoop m
  \def\r{0.72}
  \coordinate (T) at ($(TL)+0.33*(\dx,\dy)$);
  \coordinate (C) at ($(T)+\r*(\nx,\ny)$);
  \draw (C) circle (\r);
  \node at ($(C)+(162:\r+0.25)$) {$m$};

  % Clockwise rotation arrow phi
  \draw[->] ($(C)+(38:\r+0.26)$) arc[start angle=38, end angle=-18, radius=\r+0.26];
  \node[right=1pt] at ($(C)+(10:\r+0.3)$) {$\varphi$};
\end{tikzpicture}
\end{center}
\end{problem}
\clearpage


\begin{problem}{9.3.26}
证明带电粒子在稳定电磁场中的能量积分是 $T + q\phi$ ，其中 $q$ 是带电粒子所带的电荷量， $\phi$ 是电磁场的标势， $T$ 是带电粒子的动能，
\end{problem}
\clearpage


\begin{problem}{9.3.27}
长 $l$ 、质量为 $M$ 的均质棒在竖直的 $xz$ 平面内运动, 其一端 $A$ 限制在直线 $z = x \tan \alpha$ 上运动 ( $\alpha$ 为直线对水平的 $x$ 轴的倾角), 用图 7.59 所示的广义坐标 $s$ 和 $\theta$ , 导出运动微分方程, 并确定作平动时的 $\theta$ 值.
\begin{center}
\begin{tikzpicture}[>=Stealth]
  \coordinate (O) at (0,0);
  \coordinate (B) at (0,0.055);
  \coordinate (Ytop) at (0,3.6);
  \coordinate (Rtop) at ($(B) + (68:3.1)$);

  % Coordinate axes
  \draw[->, thick] (-1.8,0) -- (2.4,0) node[right] {$x$};
  \draw[->, thick] (O) -- (Ytop) node[left] {$y$};
  \node[below right=1pt and -1pt] at (O) {$O$};

  % Tilted uniform rod
  \draw[thick, fill=white] (B) -- (Rtop) -- ++(-22:0.14) -- ++(-112:3.1) -- cycle;

  % Angle theta
  \pic[draw, angle radius=0.95cm, angle eccentricity=1.28, "$\theta$"] {angle = Rtop--B--Ytop};

  % Amplitude A indicators on the left
  \draw (-0.65,0.55) -- (-0.1,0.55);
  \draw (-0.65,-0.55) -- (-0.1,-0.55);
  \draw[<->] (-0.32,0.02) -- (-0.32,0.55) node[midway, left=1pt] {$A$};
  \draw[<->] (-0.32,-0.55) -- (-0.32,-0.02) node[midway, left=1pt] {$A$};
\end{tikzpicture}

\end{center}
\end{problem}
\clearpage


\begin{problem}{9.3.28}
(1) 质量为 $m$ 的均质杆两端的速度分别为 $\bm{v}_A$ 和 $\bm{v}_B$ , 证明其动能为

\[
T = \frac {1}{6} m (\boldsymbol {v} _ {A} ^ {2} + \boldsymbol {v} _ {A} \cdot \boldsymbol {v} _ {B} + v _ {B} ^ {2})
\]

(2) 一根质量为 $m$ 、长度为 $a$ 的均质杆, 其一端限于沿一光滑的竖直轴运动, 另一端在光滑的水平面上滑动, 设 $\theta$ 是杆与竖直向下的轴间的夹角, $\varphi$ 是杆所在竖直平面的角位移, 如 $t=0$ 时, $\theta=\frac{\pi}{6}, \dot{\theta}=0, \dot{\varphi}=\Omega$ , 证明在以后的运动中, 只要杆保持与水平面接触, 有下列关系:

\[
\dot {\theta} ^ {2} = \frac {1}{1 6} \Omega^ {2} (4 - \csc^ {2} \theta) + \frac {3 g}{2 \Omega} (\sqrt {3} - 2 \cos \theta)
\]
\end{problem}
\clearpage


\begin{problem}{9.3.29}
一个质量为 $m$ 的质点, 在重力的作用下自静止竖直下落时, 受到一个可用耗散函数 $R = \frac{1}{2} \gamma v^2$ 的阻力作用, 证明质点的最大速度为 $\frac{mg}{\gamma}$
\end{problem}
\clearpage


\begin{problem}{9.3.30}
自然长度为 $l_{0}$ 、劲度系数为 $k$ 的无质量弹簧一端固定，另一端和一个质量为 $m$ 的质点相连，在重力作用下，限在一个竖直平面内运动。用 $\varphi$ 和 $\lambda = \frac{r - r_0}{r_0}$ （其中 $r_0$ 是处于

平衡状态时弹簧的长度)为广义坐标.

(1) 写出拉格朗日函数；

(2) 用 $\omega_{s}^{2} = \frac{k}{m}, \omega_{p}^{2} = \frac{g}{r_{0}}$ 写出拉格朗日方程；

(3) 当 $\lambda$ 和 $\varphi$ 很小, 初始条件: $t = 0$ 时, $\varphi = 0, \lambda = A, \dot{\varphi} = \omega_p B, \dot{\lambda} = 0$ , 其中 $A, B$ 是常量, 求运动的最低级近似解;

（4）讨论运动的高一级近似，在什么条件下， $\lambda$ 的运动将发生共振？这在物理上能实现吗？
\begin{center}
\begin{tikzpicture}[>=Stealth]
  \coordinate (O) at (0,0);
  \coordinate (V) at (0,-3.5);
  \coordinate (M) at (-66:3.4);

  % Hatched ceiling
  \fill[pattern=north east lines] (-1.2,0) rectangle (1.4,0.22);
  \draw[thick] (-1.2,0) -- (1.4,0);

  % Vertical reference line
  \draw[dashed] (O) -- (V);

  % Tilted spring and bob
  \draw[thick] (O) -- (-66:0.95);
  \draw[thick, decorate, decoration={zigzag, segment length=6.5pt, amplitude=3.5pt}] (-66:0.95) -- (-66:2.75);
  \draw[thick] (-66:2.75) -- (M);
  \filldraw[fill=white, draw=black, thick] (M) circle (3.5pt);
  \node[below right=1pt and 3pt] at (M) {$m$};
  \node[right=7pt] at (-66:1.85) {$k$};

  % Angle phi
  \pic[draw, angle radius=0.72cm, angle eccentricity=1.35, "$\varphi$"] {angle = V--O--M};

  % Dimension r parallel to the spring
  \coordinate (R1) at (0.72,0);
  \coordinate (R1dir) at ($(R1) + (-66:1)$);
  \coordinate (R2) at ($(R1)!(M)!(R1dir)$);
  \draw (M) -- ($(R2) + (24:0.25)$);
  \draw[<->] (R1) -- (R2) node[midway, right=2pt] {$r$};
\end{tikzpicture}

\end{center}
\end{problem}
\clearpage


\begin{problem}{9.3.31}
两个质量均为 $m$ 的质点各用一根劲度系数为 $k$ 、原长为 $l$ 的弹簧连到质量为 $M$ 、长度为 $L$ 的均质细杆的中点，并限在细杆上做两根弹簧的长度保持相等的运动，细杆可以以任意方式自由运动，此系统是空间的孤立系统，列出运动微分方程并求解之（到作积分为止）。定性地描述系统的运动。
\end{problem}
\clearpage


\begin{problem}{9.3.32}
将一根很长的、可忽略质量的细管安上轴，可以无摩擦地绕轴在水平面内转动，一根长为 $l$ 、质量为 $M$ 的细杆无摩擦地在管内滑动（见图9.63），写出系统的拉格朗日方程；开始杆的中心位于转轴，细管以角速度 $\omega_0$ 旋转，证明杆在这个位置是不稳定的。将它稍微扰动一下，描述以后的运动，问经过很长时间之后，杆的径向速度和角速度各为多少（假定细管非常长，杆仍在管中）？
\begin{center}
\begin{tikzpicture}[>=Stealth]
  % Horizontal reference line
  \draw[dashed] (-0.9,0) -- (3.8,0);

  % Inclined tube, sliding rod, and dimensions
  \begin{scope}[rotate=32]
    \draw[thick] (-1.0,0.11) -- (3.2,0.11);
    \draw[thick] (-1.0,-0.11) -- (3.2,-0.11);
    \draw[thick, dashed] (3.2,0.11) -- (3.9,0.11);
    \draw[thick, dashed] (3.2,-0.11) -- (3.9,-0.11);

    % Sliding rod of length l
    \fill[pattern=north east lines] (2.4,-0.11) rectangle (3.3,0.11);
    \draw[thick] (2.4,-0.11) -- (2.4,0.11) (3.3,-0.11) -- (3.3,0.11);

    % Dimension x
    \draw (0,0.18) -- (0,0.85);
    \draw (2.4,0.18) -- (2.4,0.85);
    \draw[<->] (0,0.7) -- (2.4,0.7) node[midway, above=2pt] {$x$};

    % Dimension l
    \draw (2.4,-0.18) -- (2.4,-0.75);
    \draw (3.3,-0.18) -- (3.3,-0.75);
    \draw[<->] (2.4,-0.6) -- (3.3,-0.6) node[midway, below=2pt] {$l$};
  \end{scope}

  % Pivot point and angle phi
  \fill (0,0) circle (2.5pt);
  \draw (1.05,0) arc (0:32:1.05);
  \node at (16:1.35) {$\varphi$};
\end{tikzpicture}

\end{center}
\end{problem}
\clearpage


\begin{problem}{9.3.33}
一质量为 $M$ 的木块和一个半径为 $a$ 的无质量的竖直的圆周轨道刚性相连，一起置于光滑的水平桌面上，如图9.64所示，一个质量为 $m$ 的质点沿圆周轨道做无摩擦的运动.

(1) 以 $\theta$ 为一个坐标, 写出拉格朗日函数;

(2) 找出运动微分方程；

(3) 在小角度极限下, 求 $\theta(t)$ .
\begin{center}
\begin{tikzpicture}[>=Stealth]
  % Hatched horizontal table
  \fill[pattern=north east lines] (-2.0,-0.22) rectangle (3.4,0);
  \draw[thick] (-2.0,0) -- (3.4,0);

  % Circular track
  \coordinate (C) at (0,1.35);
  \draw[thick] (C) circle (1.35);
  \draw[thick] (C) circle (1.15);

  % Rigid connector and block M
  \draw[thick] (1.35,1.35) -- (2.15,1.35);
  \draw[thick] (2.15,0) rectangle (3.35,1.65);
  \node at (2.75,0.825) {$M$};

  % Radius a
  \draw[->, thick] (C) -- ++(155:1.35) node[midway, above=2pt] {$a$};

  % Particle m and angle theta
  \coordinate (V) at (0,0.2);
  \coordinate (P) at ($(C) + (-128:1.15)$);
  \draw[dashed] (C) -- (V);
  \draw[dashed] (C) -- (P);
  \fill (P) circle (2.8pt);
  \node at ($(C) + (-132:1.62)$) {$m$};

  \pic[draw, angle radius=0.6cm, angle eccentricity=1.35, "$\theta$"] {angle = P--C--V};
\end{tikzpicture}

\end{center}
\end{problem}
\clearpage


\begin{problem}{9.3.34}
质量为 $M$ 、半径为 $R$ 的两个相同的均质圆盘用三根相同的扭杆支承，如图9.65所示。扭杆的恢复力矩为 $\tau = -k\theta (k$ 为适用于此 $l$ 和 $\theta$ 的扭转常数），圆盘可围绕扭杆的竖直轴自由转动， $\theta_{1}, \theta_{2}$ 分别表示两盘偏离平衡位置的角位移，忽略扭杆的转动惯量。初始条件为 $\theta_{1}(0) = 0, \theta_{2}(0) = 0, \dot{\theta}_{1}(0) = 0, \dot{\theta}_{2}(0) = \Omega$ （常量）。要使圆盘1得到系统的全部动能需要多长时间？可以保留隐函数形式。
\begin{center}
\begin{tikzpicture}[>=Stealth]
  % Left panel: side view
  \fill[pattern=north east lines] (-1.0,3.6) rectangle (1.0,3.82);
  \draw[thick] (-1.0,3.6) -- (1.0,3.6);
  \fill[pattern=north east lines] (-1.0,-0.22) rectangle (1.0,0);
  \draw[thick] (-1.0,0) -- (1.0,0);

  % Torsion shaft segments
  \draw[thick] (-0.08,0) rectangle (0.08,0.96);
  \draw[thick] (-0.08,1.32) rectangle (0.08,2.28);
  \draw[thick] (-0.08,2.64) rectangle (0.08,3.60);

  % Two identical disks
  \draw[thick, fill=white] (-0.8,0.96) rectangle (0.8,1.32);
  \draw[thick, fill=white] (-0.8,2.28) rectangle (0.8,2.64);
  \node at (-1.04,2.46) {$1$};
  \node at (-1.04,1.14) {$2$};

  % Length labels l
  \draw[<->] (-0.58,0) -- (-0.58,0.96) node[midway, left=1pt] {$l$};
  \draw[<->] (-0.58,1.32) -- (-0.58,2.28) node[midway, left=1pt] {$l$};
  \draw[<->] (-0.58,2.64) -- (-0.58,3.60) node[midway, left=1pt] {$l$};

  % Right panel: top view
  \coordinate (D) at (2.75,1.6);
  \draw[thick] (D) circle (1.05);
  \draw (2.75,0.3) -- (2.75,3.3);
  \draw (D) -- ++(52:1.55);
  \draw (D) -- ++(28:1.65);

  % Angular displacements theta_1 and theta_2
  \draw[->] (D) ++(90:0.82) arc (90:52:0.82);
  \node at ($(D) + (71:0.56)$) {$\theta_1$};
  \draw[->] (D) ++(90:1.35) arc (90:28:1.35);
  \node at ($(D) + (59:1.60)$) {$\theta_2$};
\end{tikzpicture}

\end{center}
\end{problem}
\clearpage


\begin{problem}{9.3.35}
长 $2L$ 、质量 $M$ 的均匀细杆用一根长 $l$ 、质量不计不可伸长的绳子拴在固定的钉子上。如图 9.66 所示。在杆的自由端加一水平的恒力 $F$ ，写出这个系统的拉格朗日函数和运动微分方程； $t=0$ 时，系统从下垂位置的静止状态开始运动，求在非常短的时间内（所有角度都很小）绳子和杆与竖直方向的夹角。
\begin{center}
\begin{tikzpicture}[>=Stealth]
  % Hatched ceiling
  \fill[pattern=north east lines] (-0.9,0) rectangle (1.1,0.22);
  \draw[thick] (-0.9,0) -- (1.1,0);

  % String and uniform rod
  \draw[thick] (0,0) -- (0,-1.3);
  \draw[thick] (-0.12,-1.3) rectangle (0.12,-3.7);
  \node at (0.35,-2.5) {$M$};

  % Horizontal constant force F
  \draw[->, thick] (0.12,-3.7) -- (1.35,-3.7) node[right=2pt] {$F$};

  % Length dimensions on the left
  \draw (-0.85,-1.3) -- (-0.22,-1.3);
  \draw (-0.85,-3.7) -- (-0.22,-3.7);
  \draw[<->] (-0.65,0) -- (-0.65,-1.3) node[midway, left=1pt] {$l$};
  \draw[<->] (-0.65,-1.3) -- (-0.65,-3.7) node[midway, left=1pt] {$2L$};
\end{tikzpicture}

\end{center}
\end{problem}
\clearpage


\begin{problem}{9.3.36}
考虑一个双星系统.

(1) 用两颗星的笛卡儿坐标写出系统的拉格朗日函数；

(2) 证明其势能是坐标的 $(-1)$ 次齐次函数, 即

\[
V (\alpha \boldsymbol {r} _ {1}, \alpha \boldsymbol {r} _ {2}) = \alpha^ {- 1} V (\boldsymbol {r} _ {1}, \boldsymbol {r} _ {2})
\]

其中 $\alpha$ 是任意实数；

（3）作一个变换，它保持由拉格朗日函数得到的运动微分方程不变，由此找出联系转动周期和它的轨道有关尺寸的开普勒第三定律.
\begin{center}
\begin{tikzpicture}[>=Stealth]
  \begin{scope}[rotate=38]
    % Center of mass C marker
    \draw[thick] (1.65,-0.09) -- (1.65,0.09);

    % Position vectors from C
    \draw[->, thick] (1.58,0.04) -- (0.08,0.04) node[midway, above left=2pt and -2pt] {$\boldsymbol{r}_2$};
    \draw[->, thick] (1.72,0.04) -- (4.42,0.04) node[midway, above left=2pt and -2pt] {$\boldsymbol{r}$};

    % Vector below from m2 to m1
    \draw[thick] (0,0) -- (0,-0.52);
    \draw[->, thick] (0,-0.52) -- (4.42,-0.52) node[pos=0.68, below right=2pt and -2pt] {$\boldsymbol{r}_1$};

    % Point labels
    \node[below left=3pt and 2pt] at (0,0) {$m_2$};
    \node[above right=2pt and 2pt] at (4.5,0) {$m_1$};
    \node at (1.65,-0.28) {$C$};
  \end{scope}
\end{tikzpicture}

\end{center}
\end{problem}
\clearpage


\begin{problem}{9.3.37}
质量为 M、长为 l 的两均质细杆由光滑的铰链和一根线连接起来，置于光滑的地面上，处于静止状态，杆与水平地面的夹角为 $30^{\circ}$ ，如图 9.69 所示。铰链的质量可以不计，t=0 时，线被切断。求：

(1) 铰链击中地面时的速度；

(2) 铰链击中地面时的时间(写出积分式子即可, 不必求积).
\begin{center}
\begin{tikzpicture}[>=Stealth]
  % Hatched horizontal ground
  \fill[pattern=north east lines] (-2.7,-0.22) rectangle (2.7,0);
  \draw[thick] (-2.7,0) -- (2.7,0);

  % Two uniform rods of length l at 30 deg
  \coordinate (Lbot) at (-2.295,0);
  \coordinate (Rbot) at (2.295,0);
  \coordinate (Apex) at (0,1.325);

  \draw[thick, fill=white] (Lbot) -- (Apex) -- ++(120:0.11) -- ++(-150:2.65) -- cycle;
  \draw[thick, fill=white] (Rbot) -- (Apex) -- ++(60:0.11) -- ++(-30:2.65) -- cycle;

  % Connecting string
  \draw[thick] (-1.048,0.72) -- (1.048,0.72);
  \node[below=2pt] at (0,0.72) {线};

  % Smooth hinge at apex
  \filldraw[fill=white, draw=black, thick] (0,1.36) circle (0.11);
  \node (hinge) at (0.85,1.95) {铰链};
  \draw[->, thick] (hinge.west) to[bend right=25] (0.12,1.48);

  % Mass M label
  \node[above right=1pt and 1pt] at (1.15,0.75) {$M$};

  % 30 degree angle markings
  \draw (-1.55,0) arc (0:30:0.745);
  \node at (-1.12,0.22) {$30^\circ$};
  \draw (1.55,0) arc (180:150:0.745);
  \node at (1.12,0.22) {$30^\circ$};

  % Length l dimension along left rod
  \draw ($(Lbot) + (120:0.20)$) -- ($(Lbot) + (120:0.78)$);
  \draw ($(Apex) + (120:0.20)$) -- ($(Apex) + (120:0.78)$);
  \draw[<->] ($(Lbot) + (120:0.62)$) -- ($(Apex) + (120:0.62)$) node[midway, above left=1pt and 1pt] {$l$};
\end{tikzpicture}

\end{center}
\end{problem}
\clearpage


\begin{problem}{9.3.38}
球面摆由质量为 $m$ 的质点, 用一根长为 $l$ 的不可伸长的轻绳悬挂在一个固定点构成, 该质点被限制在球面上运动.

(1) 若使绳与铅直方向的夹角 $\theta_{0}$ 不变, 质点做此圆周运动的角速度多大?

(2) 在(1)问所述的圆轨道上的质点受一垂直于其速度的冲量, 结果在其轨道达最高点时绳与铅直方向成 $\theta_{1}$ 角, 写出当质点在最低点时绳与铅直方向的夹角 $\theta_{2}$ 所满足的方程 (不必解);

(3) 就在 $\theta_{0}$ 附近振动的振幅很小的情况, 求振动的角频率.
\end{problem}
\clearpage


\begin{problem}{9.3.39}
一弹簧摆由质量为 $m$ 的质点经原长为 $l$ 、劲度系数为 $k$ 的无质量弹簧系于固定点组成，如图9.71所示。假定系统限制在一个竖直平面内运动，导出运动微分方程，

并在对平衡位置的角位移和径向位移都很小的近似下解运动微分方程.
\begin{center}
\begin{tikzpicture}[>=Stealth]
  \coordinate (O) at (0,0);
  \coordinate (V) at (0,-3.2);
  \coordinate (M) at (-58:3.1);

  % Hatched ceiling
  \fill[pattern=north east lines] (-1.0,0) rectangle (1.3,0.22);
  \draw[thick] (-1.0,0) -- (1.3,0);

  % Vertical reference line
  \draw[dashed] (O) -- (V);

  % Tilted spring and bob
  \draw[thick] (O) -- (-58:0.98);
  \draw[thick, decorate, decoration={zigzag, segment length=6.5pt, amplitude=3.5pt}] (-58:0.98) -- (-58:2.48);
  \draw[thick] (-58:2.48) -- (M);
  \filldraw[fill=white, draw=black, thick] (M) circle (3.5pt);
  \node[right=4pt] at (M) {$m$};
  \node[above right=0pt and 6pt] at (-58:1.72) {$k$};

  % Angle phi
  \pic[draw, angle radius=0.72cm, angle eccentricity=1.35, "$\varphi$"] {angle = V--O--M};
\end{tikzpicture}

\end{center}
\end{problem}
\clearpage


\begin{problem}{9.3.40}
粒子受到力心的引力作用, 力的大小反比于它到力心的距离的三次方, 导出运动微分方程, 并解方程求轨道, 讨论轨道的性质是如何依赖于系统的参数的.
\end{problem}
\clearpage


\begin{problem}{9.3.41}
一个质量为 $m$ 的粒子在有心力

\[
F (r) = - \frac {k}{r ^ {2}} + \frac {k ^ {\prime}}{r ^ {3}} \quad (k > 0)
\]

作用下做平面运动.

(1) 写出系统的拉格朗日函数；

(2) 导出运动微分方程, 并证明轨道角动量 J 是一个运动常数;

(3) 若 $J^{2}>-mk'$ ，求轨道方程.
\end{problem}
\clearpage


\begin{problem}{9.3.42}
一个质量为 $m$ 的珠子在一个半径为 $a$ 的旋转圆环上无摩擦地滑动, 圆环绕通过竖直的直径作恒定角速度 $\omega$ 的转动, 如图 9.72 所示:

(1) 写出系统的拉格朗日函数, 并找出所有的运动积分;

(2) 确定 $\omega < \omega_0$ 和 $\omega > \omega_0\left(\omega_0 = \sqrt{\frac{g}{a}}\right)$ 时珠子的平衡位置；

(3) 这些平衡位置,哪些是稳定的?哪些是不稳定的?

(4) 计算围绕稳定平衡点的小振幅振动的角频率.
\begin{center}
\begin{tikzpicture}[>=Stealth]
  \coordinate (O) at (0,0);
  \coordinate (V) at (0,-1.65);
  \coordinate (P) at (-54:1.65);

  % Rotating circular hoop and vertical axis
  \draw[thick] (O) circle (1.65);
  \draw[thick] (0,-2.25) -- (0,2.65);

  % Rotation arrow omega
  \draw[->, thick] (-0.32,2.22) arc (140:395:0.32 and 0.11);
  \node at (0.55,2.24) {$\omega$};

  % Radius a and bead m
  \draw[thick] (O) -- (P) node[pos=0.52, above right=-1pt and 2pt] {$a$};
  \fill (P) circle (3pt);
  \node[below right=1pt and 2pt] at (P) {$m$};

  % Angle phi
  \pic[draw, angle radius=0.65cm, angle eccentricity=1.32, "$\varphi$"] {angle = V--O--P};
\end{tikzpicture}

\end{center}
\end{problem}
\clearpage


\begin{problem}{9.3.43}
质量为 $m_{1}$ 和 $m_{2}$ 的两质点被一长 l 并穿过一光滑水平面上的光滑小孔的不可伸长的轻绳连接， $m_{1}$ 在水平面上运动， $m_{2}$ 在铅垂线上运动。

(1) 要使 $m_{2}$ 在平面下 d 处保持不动, 必须给 $m_{1}$ 以什么样的初速度;

(2) 如果 $m_{2}$ 在铅直方向稍微偏离上述平衡位置, 将产生小振动, 用拉格朗日方程求振动周期.
\end{problem}
\clearpage


\begin{problem}{9.3.44}
一个质量为 $M$ 、半径为 $R$ 的均质圆球沿一质量为 $m$ 、倾角为 $\alpha$ 的三角块向下做纯滚动，三角块在光滑的水平面上自由运动，如图9.73所示。

(1) 求系统的拉格朗日函数, 写出拉格朗日方程;

(2) 设开始系统处于静止, 球心在水平面上 H 处, 求系统的运动.
\begin{center}
\begin{tikzpicture}[>=Stealth, thick]
    % Horizontal ground
    \draw (0,0) -- (5.2,0);
    \foreach \x in {0.15,0.4,...,5.05} {
        \draw[thin] (\x,0) -- (\x-0.15,-0.16);
    }

    % Triangular wedge (right angle at bottom-left)
    \coordinate (A) at (0.4,0);
    \coordinate (B) at (4.4,0);
    \coordinate (C) at (0.4,1.8);
    \draw (A) -- (B) -- (C) -- cycle;

    % Angle alpha at bottom-right corner
    \pic[draw, thin, angle radius=0.65cm, angle eccentricity=1.35, "$\alpha$"] {angle = C--B--A};

    % Sphere resting on inclined plane
    % Slope vector: (4.0, -1.8), unit normal: (0.4104, 0.9119)
    % Contact point at t = 0.25 from C to B: (1.4, 1.35)
    % Radius R = 0.65 -> Center O = (1.67, 1.94)
    \coordinate (O) at (1.67, 1.94);
    \draw (O) circle (0.65);

    % Dashed height reference line from sphere center
    \draw[dashed, thin] (O) -- (5.0, 1.94);

    % Height dimension arrow H
    \draw[<->, thin] (4.8, 0) -- (4.8, 1.94) node[midway, right] {$H$};
\end{tikzpicture}

\end{center}
\end{problem}
\clearpage


\begin{problem}{9.3.45}
质量为 $m$ 的质点被约束在一半径为 $r$ 、质量为 $M$ 的均质圆环内表面无摩擦

运动, 环在竖直的 $xy$ 平面内沿水平的 $x$ 轴做纯滚动, 如图9.75所示. $t = 0$ 时, 环是静止的, 质点在环的顶部, 具有沿 $x$ 轴方向的速度 $v_0$ , 求当质点运动到环的底部时相对于固定的 $x$ 轴的速度 $v$ , 并给出当 $\frac{m}{M} \to 0$ 及 $\frac{M}{m} \to 0$ 时的结果.
\begin{center}
\begin{tikzpicture}[>=Stealth, thick]
    % Coordinate axes
    \draw[->] (0,0) -- (4.8,0) node[below left] {$x$};
    \draw[->] (0,0) -- (0,3.6) node[left] {$y$};

    % Ring of mass M and radius r = 1.45 resting on x-axis
    \coordinate (C) at (2.7, 1.45);
    \draw (C) circle (1.45);
    \node at ($(C) + (-35:1.72)$) {$M$};

    % Vertical dashed diameter line
    \draw[dashed, thin] (2.7, 0) -- (2.7, 3.6);

    % Particle m inside ring at angle theta from top vertical
    \coordinate (P) at ($(C) + (45:1.45)$);
    \draw[thin] (C) -- (P);
    \fill (P) circle (2.2pt) node[above right] {$m$};

    % Angle theta arc from vertical to radius
    \draw[->, thin] ($(C) + (90:0.58)$) arc (90:45:0.58);
    \node at ($(C) + (67.5:0.82)$) {$\theta$};
\end{tikzpicture}

\end{center}
\end{problem}
\clearpage


\begin{problem}{9.3.46}
一质点在光滑的旋转抛物面 $r^2 = az$ 内侧滑动，证明受到的约束力的大小为

\[
\text {常量} \cdot \left(1 + \frac {4 r ^ {2}}{a ^ {2}}\right) ^ {- 3 / 2}
\]

说明约束力的方向.
\end{problem}
\clearpage


\begin{problem}{9.3.47}
一个肥皂泡的半径 r 按 $r = at + b$ (其中 a、b 为常量) 的规律随时间变化, 一个质量为 m 的质点保持在表面上运动. t = 0 时, 给质点一个沿切向的速度 v, 不考虑重力的作用, 求质点的运动规律及表面对质点的作用力 (限用带未定乘子的拉格朗日方程).
\end{problem}
\clearpage


\begin{problem}{9.3.48}
一个质量为 $m$ 的质点被约束沿一弯成半径为 $a$ 的水平圆形金属丝上运动, 摩擦因数为 $\mu$ , 开始速率为 $v_{0}$ , 若不计重力的作用, 求质点的角位置与时间的关系, 以及金属丝对质点的作用力.
\end{problem}
\clearpage


\begin{problem}{9.3.49}
一根光滑的金属丝弯成一螺旋线, 其方程为 $z = a\varphi, \rho = b, a, b$ 为常量, $\rho, \varphi, z$ 为柱坐标, 原点为引力中心, 引力大小与距离成正比, 比例系数为 $k$ , 有一质量为 $m$ 的珠子在金属丝上做无摩擦滑动. 求金属丝对珠子的作用力的 $\rho, \varphi, z$ 分量 (不考虑重力的作用).
\end{problem}
\clearpage


\begin{problem}{9.3.50}
一个质量为 $m$ 的质点在一个光滑的抛物面[其方程为 $x^{2} + y^{2} = a(a - z)$ , 其中 $a$ 是常量, $z$ 轴竖直向上]的顶端从静止开始下滑, 选柱坐标 $\rho, \varphi, z$ 为广义坐标, 写出含未定乘子的拉格朗日方程, 并求出约束力, 在何处质点离开抛物面?
\end{problem}
\clearpage


\begin{problem}{9.3.51}
一个质量为 $m$ 、半径为 $a$ 的均质圆柱体 $A$ 在一个固定的、半径为 $b$ 的圆柱体 $B$ 的顶部处于静止的状态，两圆柱体的轴平行且水平，让圆柱体 $A$ 稍偏离其平衡位置做无滑动滚动下圆柱体 $B$ . 用带未定乘子的拉格朗日方程求约束力，并找出两圆柱体分离的位置.
\end{problem}
\clearpage


\begin{problem}{9.3.52}
一个质量为 $m$ 、半径为 $a$ 的均质圆盘，受到约束，保持盘面铅直地在一完全粗糙的水平面上滚动，受到一个力 $f$ 作用，作用线通过盘心， $f = f_{1}i + f_{2}j, i, j$ 分别是水平面上的 $x, y$ 轴的单位矢量， $z$ 轴垂直于水平面。取固连于圆盘的坐标系 $x', y', z'$ ，原点位于盘心 $C, x'y'$ 轴在盘面上， $z'$ 轴垂直于盘面，取 $z', z$ 轴的夹角为 $\theta, \varphi$ 是 $i$ 和 $k \times k'(k'$ 是 $z'$ 轴的单位矢量）间的夹角， $\psi$ 是 $i'$ 和 $k \times k'$ 的夹角， $\theta, \varphi, \psi$ 是欧拉角，选质心 $C$ 的 $x, y$ 坐标和 $\varphi, \psi$ 为广义坐标，写出盘面的运动微分方程。
\end{problem}
\clearpage


\begin{problem}{9.3.53}
一个质量为 $m$ 的质点, 受到一个力 $F = F_{x}i + F_{y}j + F_{z}k$ 作用以及一个非完整约束 $\dot{x} - t\dot{y} = a$ , 其中 $a$ 为常量, 求运动微分方程.
\end{problem}
\clearpage


\section{冲击运动 机电模拟}

\begin{problem}{9.4.1}
用拉格朗日方程解7.4.33题.
\end{problem}
\clearpage


\begin{problem}{9.4.2}
用拉格朗日方程解7.4.34题.
\end{problem}
\clearpage


\begin{problem}{9.4.3}
一个由质量为 $m$ 和 $M$ 的两个质点用一根不计质量、长 $a$ 的棒构成的哑铃开始处于静止状态，突然在棒上距 $Mb$ 处受到一个垂直于棒的冲量 $I$ ，求作用后哑铃的运动状态。
\end{problem}
\clearpage


\begin{problem}{9.4.4}
质量为 $m_{1}$ 与 $m_{2}$ 的两粒子被一劲度系数为 k 的轻弹簧连接, 静止在光滑的水平面上.

(1) 突然沿从 $m_{1}$ 到 $m_{2}$ 的方向给予 $m_{1}$ 一个冲量 I，问 $m_{2}$ 在碰撞后第一次到达静止之前运行多远；

(2) 是否可能给 $m_{1}$ 一冲量使系统由静止变成无振动的转动？
\end{problem}
\clearpage


\begin{problem}{9.4.5}
两根长度为 $a$ 、质量为 $m$ 的棒 $AB$ 和 $BC$ , 在 $B$ 点无摩擦地连接, 放在光滑的水平桌面上, 开始两棒(即点 A、B、C)共线, 在 A 点施以与 ABC 线垂直的冲量 I, 求冲击后的瞬时两棒的运动状态.
\end{problem}
\clearpage


\begin{problem}{9.4.6}
质量均为 $m$ 、长度均为 $2a$ 的两根均质棒 $PQ, QR$ 在 $Q$ 点光滑铰连，开始置于光滑的桌面上成一直线，在 $PQ$ 的 $P$ 端作用一冲量 $I$ ，其方向沿垂直于 $PQ$ 的水平方向。求两棒在作用后各自质心的初速度和初角速度。
\end{problem}
\clearpage


\begin{problem}{9.4.7}
一个质量为 $m$ 、半径为 $a$ 的均质圆环可围绕边缘上一点 $O$ 自由转动, 开始处于静止状态, 突然在边缘上另一点 $P$ 作用一个冲量 $I$ , 其方向垂直于圆环所在平面. 证明: 由此冲量产生的角速度大小不会超过 $\frac{3I}{\sqrt{2}ma}$ , 并求产生这么大的角速度 $OP$ 的距离.
\end{problem}
\clearpage


\begin{problem}{9.4.8}
由长度为 $2l$ 、质量为 $m$ 的四根杆用光滑铰链连接成的构件 $ABCD$ ，开始成正方形静置于光滑的水平面上，取图9.84所示的坐标，质心位于坐标原点， $AB, CD$ 与 $x$ 轴平行， $AD, BC$ 与 $y$ 轴平行，今在 $A$ 点沿 $AD$ 方向作用一冲量 $I$ ，求作用后构件的运动状态。
\begin{center}
\begin{tikzpicture}[>=Stealth, thick]
    % Square linkage ABCD centered at origin O
    \coordinate (A) at (-1.5,-1.5);
    \coordinate (B) at (1.5,-1.5);
    \coordinate (C) at (1.5,1.5);
    \coordinate (D) at (-1.5,1.5);

    \draw (A) -- (B) -- (C) -- (D) -- cycle;

    \node[left] at (A) {$A$};
    \node[below right] at (B) {$B$};
    \node[above right] at (C) {$C$};
    \node[above left] at (D) {$D$};

    % Coordinate axes from origin O
    \draw[->] (0,0) -- (2.1,0) node[below] {$x$};
    \draw[->] (0,0) -- (0,2.1) node[left] {$y$};
    \node[below left] at (0,0) {$O$};

    % Impulse arrow I at A along AD direction
    \draw[->] (-1.5,-2.45) -- (-1.5,-1.55);
    \node[right] at (-1.5,-2.0) {$I$};
\end{tikzpicture}

\end{center}
\end{problem}
\clearpage


\begin{problem}{9.4.9}
分别用基尔霍夫定律和拉格朗日方程得到图 9.86 所示电路的电路方程.
\begin{center}
\begin{tikzpicture}[>=Stealth, thick,
    res/.style={draw, fill=white, minimum width=0.72cm, minimum height=0.22cm, inner sep=0pt},
    vres/.style={draw, fill=white, minimum width=0.22cm, minimum height=0.72cm, inner sep=0pt}]
    % Outer rectangle & middle branch
    \draw (0,0) -- (0,2.8) -- (5.0,2.8) -- (5.0,0) -- cycle;
    \draw (2.5,0) -- (2.5,2.8);

    % Junction dots
    \fill (2.5,2.8) circle (1.8pt);
    \fill (2.5,0) circle (1.8pt);

    % Left AC source epsilon_1
    \draw[fill=white] (0,1.4) circle (0.34);
    \node at (0,1.4) {$\sim$};
    \node[right, font=\small] at (0.32,1.76) {$+$};
    \node[right, font=\small] at (0.32,1.04) {$-$};
    \node[right] at (0.38,1.4) {$\varepsilon_1$};

    % Top-left resistor R_1
    \node[res, label=above:$R_1$] at (1.25,2.8) {};

    % Middle inductor L
    \draw[fill=white, draw=none] (2.3,1.5) rectangle (2.7,2.45);
    \draw[decorate, decoration={coil, aspect=0.45, segment length=2.2mm, amplitude=1.5mm}] (2.5,2.5) -- (2.5,1.45);
    \node[right] at (2.68,1.98) {$L$};

    % Middle AC source epsilon_2
    \draw[fill=white] (2.5,0.75) circle (0.34);
    \node at (2.5,0.75) {$\sim$};
    \node[right, font=\small] at (2.82,1.11) {$+$};
    \node[right, font=\small] at (2.82,0.39) {$-$};
    \node[right] at (2.88,0.75) {$\varepsilon_2$};

    % Top-right capacitor C
    \draw[fill=white, draw=none] (3.62,2.6) rectangle (3.88,3.0);
    \draw (3.66,2.58) -- (3.66,3.02);
    \draw (3.84,2.58) -- (3.84,3.02);
    \node[above] at (3.75,3.02) {$C$};

    % Right resistor R_2
    \node[vres, label=right:$R_2$] at (5.0,1.4) {};

    % Mesh current arrows
    \draw[->] (1.00,1.25) .. controls (1.28,2.15) and (2.15,2.05) .. (1.88,0.85);
    \node at (1.52,1.42) {$\dot{q}_1$};

    \draw[->] (3.55,1.25) .. controls (3.83,2.15) and (4.55,2.05) .. (4.28,0.85);
    \node at (3.95,1.42) {$\dot{q}_2$};
\end{tikzpicture}

\end{center}
\end{problem}
\clearpage


\begin{problem}{9.4.10}
用拉格朗日方程得到图 9.87 所示电路的电路方程.

提示:电流源可作为一个约束 $\dot{q}=i$ 来处理.
\begin{center}
\begin{tikzpicture}[>=Stealth, thick,
    res/.style={draw, fill=white, minimum width=0.68cm, minimum height=0.20cm, inner sep=0pt},
    vres/.style={draw, fill=white, minimum width=0.20cm, minimum height=0.68cm, inner sep=0pt}]
    % Grid of 3 loops
    \draw (0,0) -- (0,2.2) -- (5.4,2.2) -- (5.4,0) -- cycle;
    \draw (1.8,0) -- (1.8,2.2);
    \draw (3.6,0) -- (3.6,2.2);

    % Junction dots
    \foreach \x in {1.8, 3.6} {
        \fill (\x,0) circle (1.8pt);
        \fill (\x,2.2) circle (1.8pt);
    }

    % Left current source i
    \draw[fill=white] (0,1.1) circle (0.32);
    \draw (-0.32,1.1) -- (0.32,1.1);
    \node[left] at (-0.36,1.1) {$i$};

    % Vertical resistor R_1
    \node[vres, label=right:$R_1$] at (1.8,1.1) {};

    % Top capacitor C
    \draw[fill=white, draw=none] (2.58,2.0) rectangle (2.82,2.4);
    \draw (2.62,1.98) -- (2.62,2.42);
    \draw (2.78,1.98) -- (2.78,2.42);
    \node[above] at (2.7,2.42) {$C$};

    % Vertical resistor R_2
    \node[vres, label=right:$R_2$] at (3.6,1.1) {};

    % Top resistor R_3
    \node[res, label=above:$R_3$] at (4.5,2.2) {};

    % Right AC voltage source epsilon
    \draw[fill=white] (5.4,1.1) circle (0.32);
    \node at (5.4,1.1) {$\sim$};
    \node[right, font=\small] at (5.70,1.44) {$+$};
    \node[right, font=\small] at (5.70,0.76) {$-$};
    \node[right] at (5.76,1.1) {$\varepsilon$};

    % Mesh current arrows
    \draw[->] (0.62,1.00) .. controls (0.85,1.65) and (1.42,1.55) .. (1.22,0.65);
    \node at (0.96,1.06) {$\dot{q}$};

    \draw[->] (2.56,1.00) .. controls (2.78,1.65) and (3.26,1.55) .. (3.08,0.65);
    \node at (2.86,1.06) {$\dot{q}_1$};

    \draw[->] (4.36,1.00) .. controls (4.58,1.65) and (5.02,1.55) .. (4.84,0.65);
    \node at (4.65,1.06) {$\dot{q}_2$};
\end{tikzpicture}

\end{center}
\end{problem}
\clearpage


\begin{problem}{9.4.11}
用拉格朗日方程写出图 9.88 所示电路的电路方程.
\begin{center}
\begin{tikzpicture}[>=Stealth, thick,
    res/.style={draw, fill=white, minimum width=0.66cm, minimum height=0.20cm, inner sep=0pt},
    vres/.style={draw, fill=white, minimum width=0.20cm, minimum height=0.66cm, inner sep=0pt}]
    % Outer loop (broken at bottom center between coupled coils)
    \draw (2.35,0) -- (0,0) -- (0,3.4) -- (5.4,3.4) -- (5.4,0) -- (3.05,0);

    % Center branch down to splitter for L_2 and L_3
    \draw (2.7,3.4) -- (2.7,1.75);
    \draw (2.35,1.75) -- (3.05,1.75);
    \draw (2.35,1.75) -- (2.35,0);
    \draw (3.05,1.75) -- (3.05,0);

    % Junction dots
    \fill (2.7,3.4) circle (1.8pt);
    \fill (2.7,1.75) circle (1.8pt);
    \fill (2.35,0) circle (1.8pt);
    \fill (3.05,0) circle (1.8pt);

    % Left AC source epsilon
    \draw[fill=white] (0,1.7) circle (0.34);
    \node at (0,1.7) {$\sim$};
    \node[right, font=\small] at (0.32,2.06) {$+$};
    \node[right, font=\small] at (0.32,1.34) {$-$};
    \node[right] at (0.38,1.7) {$\varepsilon$};

    % Top-left R_1 and L_1
    \node[res, label=above:$R_1$] at (0.85,3.4) {};
    \draw[fill=white, draw=none] (1.45,3.2) rectangle (2.35,3.6);
    \draw[decorate, decoration={coil, aspect=0.45, segment length=2.1mm, amplitude=1.4mm}] (1.45,3.4) -- (2.35,3.4);
    \node[below] at (1.9,3.2) {$L_1$};

    % Top-right capacitor C_2
    \draw[fill=white, draw=none] (3.92,3.2) rectangle (4.18,3.6);
    \draw (3.96,3.18) -- (3.96,3.62);
    \draw (4.14,3.18) -- (4.14,3.62);
    \node[below] at (4.05,3.16) {$C_2$};

    % Right R_4 and L_4
    \node[vres, label=right:$R_4$] at (5.4,2.4) {};
    \draw[fill=white, draw=none] (5.2,0.65) rectangle (5.6,1.55);
    \draw[decorate, decoration={coil, aspect=0.45, segment length=2.1mm, amplitude=1.4mm}] (5.4,1.55) -- (5.4,0.65);
    \node[right] at (5.58,1.1) {$L_4$};

    % Bottom R_2 and R_3
    \node[res, label=above:$R_2$] at (1.15,0) {};
    \node[res, label=above:$R_3$] at (4.25,0) {};

    % Middle capacitor C_1
    \draw[fill=white, draw=none] (2.5,2.45) rectangle (2.9,2.71);
    \draw (2.48,2.67) -- (2.92,2.67);
    \draw (2.48,2.49) -- (2.92,2.49);
    \node[right] at (2.94,2.58) {$C_1$};

    % Coupled inductors L_2 and L_3 with mutual inductance M
    \draw[fill=white, draw=none] (2.15,0.45) rectangle (2.55,1.45);
    \draw[decorate, decoration={coil, aspect=0.45, segment length=2.1mm, amplitude=1.4mm}] (2.35,1.45) -- (2.35,0.45);
    \node[left] at (2.15,0.95) {$L_2$};

    \draw[fill=white, draw=none] (2.85,0.45) rectangle (3.25,1.45);
    \draw[decorate, decoration={coil, aspect=0.45, segment length=2.1mm, amplitude=1.4mm}] (3.05,0.45) -- (3.05,1.45);
    \node[right] at (3.25,0.95) {$L_3$};

    \node at (2.7,0.22) {$M$};

    % Mesh current arrows
    \draw[->] (0.92,1.55) .. controls (1.20,2.45) and (1.98,2.30) .. (1.68,1.10);
    \node at (1.92,2.05) {$\dot{q}_1$};

    \draw[->] (3.62,1.55) .. controls (3.90,2.45) and (4.75,2.30) .. (4.45,1.10);
    \node at (4.70,2.05) {$\dot{q}_2$};
\end{tikzpicture}

\end{center}
\end{problem}
\clearpage


\begin{problem}{9.4.12}
用拉格朗日方程写出图 9.89 电路(各线圈有互感 $M_{12}, M_{23}, M_{13}$ , 图中难于表明)的电路方程.
\begin{center}
\begin{tikzpicture}[>=Stealth, thick,
    res/.style={draw, fill=white, minimum width=0.68cm, minimum height=0.20cm, inner sep=0pt}]
    % Three parallel horizontal branches
    \draw (0,0) -- (0,2.8) -- (5.8,2.8) -- (5.8,0) -- cycle;
    \draw (0,1.4) -- (5.8,1.4);

    % Junction dots
    \fill (0,1.4) circle (1.8pt);
    \fill (5.8,1.4) circle (1.8pt);

    % Top branch: R_1, L_1, C_1, E_1
    \node[res, label=below:$R_1$] at (0.8,2.8) {};
    \draw[fill=white, draw=none] (1.55,2.6) rectangle (2.5,3.0);
    \draw[decorate, decoration={coil, aspect=0.45, segment length=2.2mm, amplitude=1.4mm}] (1.55,2.8) -- (2.5,2.8);
    \node[below] at (2.02,2.6) {$L_1$};
    \draw[fill=white, draw=none] (3.22,2.6) rectangle (3.44,3.0);
    \draw (3.25,2.58) -- (3.25,3.02);
    \draw (3.41,2.58) -- (3.41,3.02);
    \node[below] at (3.33,2.58) {$C_1$};
    \draw[fill=white, draw=none] (4.42,2.6) rectangle (4.66,3.0);
    \draw[ultra thick] (4.46,2.66) -- (4.46,2.94);
    \draw (4.62,2.56) -- (4.62,3.04);
    \node[below] at (4.54,2.56) {$E_1$};

    % Middle branch: C_2, L_2, R_2, E_2
    \draw[fill=white, draw=none] (0.68,1.2) rectangle (0.9,1.6);
    \draw (0.71,1.18) -- (0.71,1.62);
    \draw (0.87,1.18) -- (0.87,1.62);
    \node[below] at (0.79,1.18) {$C_2$};
    \draw[fill=white, draw=none] (1.55,1.2) rectangle (2.5,1.6);
    \draw[decorate, decoration={coil, aspect=0.45, segment length=2.2mm, amplitude=1.4mm}] (1.55,1.4) -- (2.5,1.4);
    \node[below] at (2.02,1.2) {$L_2$};
    \node[res, label=below:$R_2$] at (3.35,1.4) {};
    \draw[fill=white, draw=none] (4.42,1.2) rectangle (4.66,1.6);
    \draw[ultra thick] (4.46,1.26) -- (4.46,1.54);
    \draw (4.62,1.16) -- (4.62,1.64);
    \node[below] at (4.54,1.16) {$E_2$};

    % Bottom branch: R_3, C_3, L_3, E_3
    \node[res, label=below:$R_3$] at (0.8,0) {};
    \draw[fill=white, draw=none] (1.92, -0.2) rectangle (2.14, 0.2);
    \draw (1.95,-0.22) -- (1.95,0.22);
    \draw (2.11,-0.22) -- (2.11,0.22);
    \node[below] at (2.03,-0.22) {$C_3$};
    \draw[fill=white, draw=none] (2.85,-0.2) rectangle (3.8,0.2);
    \draw[decorate, decoration={coil, aspect=0.45, segment length=2.2mm, amplitude=1.4mm}] (2.85,0) -- (3.8,0);
    \node[below] at (3.32,-0.2) {$L_3$};
    \draw[fill=white, draw=none] (4.42,-0.2) rectangle (4.66,0.2);
    \draw[ultra thick] (4.46,-0.14) -- (4.46,0.14);
    \draw (4.62,-0.24) -- (4.62,0.24);
    \node[below] at (4.54,-0.24) {$E_3$};

    % Current arrows along wires
    \draw[->] (5.8,2.4) -- (5.8,2.0);
    \node[right] at (5.8,2.2) {$\dot{q}_1$};

    \draw[->] (5.45,1.4) -- (5.05,1.4);
    \node[above] at (5.25,1.42) {$\dot{q}_2$};

    \draw[->] (5.8,0.95) -- (5.8,0.55);
    \node[right] at (5.8,0.75) {$\dot{q}_3$};
\end{tikzpicture}

\end{center}
\end{problem}
\clearpage


\begin{problem}{9.4.13}
图 9.90 所示的电路中包含两个相同的二极管, 假定 $E_{3} = A \mid \dot{q}_{3}^{b-1} \mid \dot{q}_{3}$ , 其中 $A, b$ 是常量, $E_{2} = E_{0} \sin \omega t$ , 各线圈有互感 $M_{12}, M_{13}, M_{23}$ . 用拉格朗日方程写出电路方程.
\begin{center}
\begin{tikzpicture}[>=Stealth, thick,
    res/.style={draw, fill=white, minimum width=0.64cm, minimum height=0.18cm, inner sep=0pt},
    vres/.style={draw, fill=white, minimum width=0.18cm, minimum height=0.64cm, inner sep=0pt}]
    % Main rails and branches
    \draw (0,0) -- (0,3.6) -- (5.6,3.6);
    \draw (0,0) -- (5.6,0);
    \draw (1.5,3.6) -- (1.5,0.95);
    \draw (1.5,0.35) -- (1.5,0);

    % Junction dots
    \fill (1.5,3.6) circle (1.7pt);
    \fill (1.5,0) circle (1.7pt);
    \fill (3.9,3.6) circle (1.7pt);
    \fill (3.9,0) circle (1.7pt);

    % Left branch (x=0): R_1, C, L_1, arrow q_1 dot
    \node[vres, label=left:$R_1$] at (0,0.8) {};
    \draw[fill=white, draw=none] (-0.2,1.58) rectangle (0.2,1.80);
    \draw (-0.22,1.61) -- (0.22,1.61);
    \draw (-0.22,1.77) -- (0.22,1.77);
    \node[left] at (-0.22,1.69) {$C$};
    \draw[fill=white, draw=none] (-0.2,2.15) rectangle (0.2,2.95);
    \draw[decorate, decoration={coil, aspect=0.45, segment length=2.0mm, amplitude=1.3mm}] (0,2.95) -- (0,2.15);
    \node[left] at (-0.18,2.55) {$L_1$};
    \draw[->] (0,3.12) -- (0,3.45);
    \node[left] at (-0.05,3.28) {$\dot{q}_1$};

    % Top-left battery E_1 (long thin on left, short thick on right)
    \draw[fill=white, draw=none] (0.62,3.4) rectangle (0.86,3.8);
    \draw (0.66,3.36) -- (0.66,3.84);
    \draw[ultra thick] (0.82,3.46) -- (0.82,3.74);
    \node[below] at (0.74,3.36) {$E_1$};

    % Second vertical branch (x=1.5): q_2 dot, R_2, L_2, E_2 = E_0 sin(wt)
    \draw[->] (1.5,3.42) -- (1.5,3.05);
    \node[right] at (1.55,3.24) {$\dot{q}_2$};
    \node[vres, label=right:$R_2$] at (1.5,2.5) {};
    \draw[fill=white, draw=none] (1.3,1.25) rectangle (1.7,2.0);
    \draw[decorate, decoration={coil, aspect=0.45, segment length=2.0mm, amplitude=1.3mm}] (1.5,2.0) -- (1.5,1.25);
    \node[right] at (1.68,1.62) {$L_2$};

    % AC terminal circles for E_2
    \draw[fill=white] (1.5,0.95) circle (2.0pt);
    \draw[fill=white] (1.5,0.35) circle (2.0pt);
    \node[right, font=\small] at (1.58,0.95) {$+$};
    \node[right, font=\small] at (1.58,0.35) {$-$};
    \node[font=\footnotesize] at (1.5,0.65) {$E_2 = E_0\sin\omega t$};

    % Top-middle branch: L_3, R_3, q_3 dot
    \draw[fill=white, draw=none] (1.85,3.4) rectangle (2.6,3.8);
    \draw[decorate, decoration={coil, aspect=0.45, segment length=2.0mm, amplitude=1.3mm}] (1.85,3.6) -- (2.6,3.6);
    \node[below] at (2.22,3.4) {$L_3$};
    \node[res, label=below:$R_3$] at (3.00,3.6) {};
    \draw[->] (3.76,3.6) -- (3.48,3.6);
    \node[above] at (3.62,3.62) {$\dot{q}_3$};

    % Left vacuum tube diode at x=3.9, center (3.9, 1.55)
    \draw (3.9,1.55) circle (0.62);
    % Anode at bottom
    \draw (3.9,0) -- (3.9,1.23);
    \draw (3.68,1.23) -- (4.12,1.23);
    % Cathode & heater at top
    \draw (3.9,3.6) -- (3.9,2.85);
    \fill (3.9,2.85) circle (1.7pt);
    \draw (3.9,2.85) -- (3.9,1.67) -- (3.68,1.67);
    % Heater loop & filament battery at top-left of diode
    \draw (3.68,1.67) -- (3.68,1.83) -- (4.06,1.83) -- (4.06,2.45) -- (3.35,2.45) -- (3.35,2.85);
    \draw (3.9,2.85) -- (3.58,2.85);
    \draw (3.35,2.85) -- (3.48,2.85);
    \draw (3.48,2.67) -- (3.48,3.03);
    \draw[ultra thick] (3.58,2.74) -- (3.58,2.96);

    % E_3 voltage arrow between top and bottom rails
    \draw[<->, thin] (4.72,0) -- (4.72,3.6) node[midway, fill=white, inner sep=1.5pt] {$E_3$};

    % Right vacuum tube diode at x=5.6, center (5.6, 1.55)
    \draw (5.6,1.55) circle (0.62);
    % Anode at top
    \draw (5.6,3.6) -- (5.6,1.87);
    \draw (5.38,1.87) -- (5.82,1.87);
    % Cathode & heater at bottom
    \draw (5.6,0) -- (5.6,0.55);
    \fill (5.6,0.55) circle (1.7pt);
    \draw (5.6,0.55) -- (5.6,1.43) -- (5.82,1.43);
    % Heater loop & filament battery at bottom-right of diode
    \draw (5.82,1.43) -- (5.82,1.27) -- (5.44,1.27) -- (5.44,0.85) -- (6.15,0.85) -- (6.15,0.55);
    \draw (5.6,0.55) -- (5.92,0.55);
    \draw (6.15,0.55) -- (6.02,0.55);
    \draw (5.92,0.37) -- (5.92,0.73);
    \draw[ultra thick] (6.02,0.44) -- (6.02,0.66);
\end{tikzpicture}

\end{center}
\end{problem}
\clearpage


\begin{problem}{9.4.14}
图 9.91 所示的系统中电容器 C 的上极板悬在劲度系数为 k 的线圈弹簧的下端(弹簧上端固定)，质量为 m，上极板在重力、弹簧力和极板间的电场力作用下在竖直方向运动，下极板是固定的，电容器的电容 $C=\frac{A}{s-x}$ ，其中 A 是与极板面积及所用的单位有关的常数，s 是电容器未充电情况下上极板处于平衡位置两极板的间距。用拉格朗日方程给出运动微分方程。
\begin{center}
\begin{tikzpicture}[>=Stealth, thick,
    res/.style={draw, fill=white, minimum width=0.68cm, minimum height=0.20cm, inner sep=0pt}]
    % Circuit path
    \draw (0,1.78) -- (0,3.2) -- (4.3,3.2) -- (4.3,2.7);
    \draw (0,1.22) -- (0,0) -- (4.3,0) -- (4.3,0.45);

    % Left AC voltage source terminals
    \draw[fill=white] (0,1.78) circle (2.2pt);
    \draw[fill=white] (0,1.22) circle (2.2pt);
    \node[right, font=\small] at (0.06,1.78) {$+$};
    \node[right, font=\small] at (0.06,1.22) {$-$};
    \node[left] at (-0.08,1.5) {$E = E_0\sin\omega t$};

    % Top branch: R, L, current q dot
    \node[res, label=below:$R$] at (1.1,3.2) {};
    \draw[fill=white, draw=none] (1.9,3.0) rectangle (2.75,3.4);
    \draw[decorate, decoration={coil, aspect=0.45, segment length=2.1mm, amplitude=1.4mm}] (1.9,3.2) -- (2.75,3.2);
    \node[below] at (2.32,3.0) {$L$};
    \draw[->] (3.0,3.2) -- (3.45,3.2);
    \node[below] at (3.22,3.15) {$\dot{q}$};

    % Upper fixed support for spring
    \draw (3.55,2.7) -- (5.05,2.7);
    \foreach \x in {3.65,3.85,...,4.95} {
        \draw[thin] (\x,2.7) -- (\x-0.14,2.56);
    }

    % Spring k
    \draw (4.3,2.7) -- (4.3,2.45);
    \draw[decorate, decoration={zigzag, segment length=2.4mm, amplitude=1.3mm}] (4.3,2.45) -- (4.3,1.65);
    \draw (4.3,1.65) -- (4.3,1.35);
    \node[right] at (4.48,2.05) {$k$};

    % Equilibrium position dashed line
    \draw[dashed, thin] (2.6,1.82) -- (5.15,1.82);
    \node[font=\small] (eq) at (3.35,2.35) {平衡位置};
    \draw[->, thin] (eq.south) -- (3.75,1.85);

    % Movable upper plate (mass m)
    \fill (3.45,1.25) rectangle (5.15,1.37);
    \node[above left] at (3.95,1.35) {$m$};

    % Fixed lower plate
    \draw (3.45,0.45) -- (5.15,0.45);
    \foreach \x in {3.55,3.75,...,5.05} {
        \draw[thin] (\x,0.45) -- (\x-0.14,0.31);
    }

    % Dimension arrows s and x
    \draw[<->, thin] (2.85,0.45) -- (2.85,1.82) node[midway, left] {$s$};
    \draw[thin] (2.7,0.45) -- (3.25,0.45);

    \draw[<->, thin] (4.95,1.37) -- (4.95,1.82) node[midway, right] {$x$};

    % Capacitance formula
    \node[font=\small] at (4.3,0.86) {$C = \dfrac{A}{s - x}$};
\end{tikzpicture}

\end{center}
\end{problem}
\clearpage


\begin{problem}{9.4.15}
图 9.92 中可变电容器的两极板可沿竖直直线 ab 运动, 无转动, 极板受弹簧力、重力和两极板间电场力的作用, $l_{10}$ 、 $l_{20}$ 是极板不带电时两极板处于平衡时两弹簧的长度, s、A 均为常量, $m_{1}$ 、 $m_{2}$ 是两极板的质量. 求系统的拉格朗日函数及运动微分方程.
\begin{center}
\begin{tikzpicture}[>=Stealth, thick,
    vres/.style={draw, fill=white, minimum width=0.20cm, minimum height=0.68cm, inner sep=0pt}]
    % Left circuit loop
    \draw (3.8,4.2) -- (3.8,4.4) -- (0,4.4) -- (0,0) -- (3.8,0) -- (3.8,0.35);

    % Left branch: L and R
    \draw[fill=white, draw=none] (-0.2,2.7) rectangle (0.2,3.55);
    \draw[decorate, decoration={coil, aspect=0.45, segment length=2.1mm, amplitude=1.4mm}] (0,3.55) -- (0,2.7);
    \node[right] at (0.18,3.12) {$L$};
    \node[vres, label=right:$R$] at (0,1.25) {};

    % Bottom branch: DC battery E and current q dot
    \draw[fill=white, draw=none] (1.72,-0.2) rectangle (1.96,0.2);
    \draw[ultra thick] (1.76,-0.14) -- (1.76,0.14);
    \draw (1.92,-0.26) -- (1.92,0.26);
    \node[above right] at (1.92,0.08) {$E$};

    \draw[->] (2.5,0) -- (2.95,0);
    \node[below] at (2.72,-0.05) {$\dot{q}$};

    % Top fixed support
    \draw[fill=white] (3.1,4.10) rectangle (4.5,4.20);
    \foreach \x in {3.2,3.4,...,4.45} {
        \draw[thin] (\x,4.20) -- (\x+0.14,4.34);
    }

    % Upper spring k_1
    \draw (3.8,4.10) -- (3.8,3.88);
    \draw[decorate, decoration={zigzag, segment length=2.3mm, amplitude=1.2mm}] (3.8,3.88) -- (3.8,3.15);
    \draw (3.8,3.15) -- (3.8,2.95);
    \node[right] at (3.96,3.58) {$k_1$};

    % Upper equilibrium dashed line
    \draw[dashed, thin] (2.0,3.42) -- (4.9,3.42);

    % Upper plate m_1
    \draw[fill=white] (2.8,2.86) rectangle (4.8,2.95);
    \node[right] at (4.82,2.90) {$m_1$};

    % Capacitance formula
    \node[font=\small] at (3.8,2.30) {$C = \dfrac{A}{s - x_1 - x_2}$};

    % Lower plate m_2
    \draw[fill=white] (2.8,1.72) rectangle (4.8,1.81);
    \node[right] at (4.82,1.76) {$m_2$};

    % Lower equilibrium dashed line
    \draw[dashed, thin] (2.0,1.12) -- (4.9,1.12);

    % Lower spring k_2
    \draw (3.8,1.72) -- (3.8,1.48);
    \draw[decorate, decoration={zigzag, segment length=2.3mm, amplitude=1.2mm}] (3.8,1.48) -- (3.8,0.75);
    \draw (3.8,0.75) -- (3.8,0.45);
    \node[right] at (3.96,1.02) {$k_2$};

    % Bottom fixed support
    \draw[fill=white] (3.1,0.35) rectangle (4.5,0.45);
    \foreach \x in {3.2,3.4,...,4.45} {
        \draw[thin] (\x,0.35) -- (\x-0.14,0.21);
    }

    % Dimension arrows: l_10, x_1, s, x_2, l_20
    \draw[<->, thin] (3.32,3.42) -- (3.32,4.10) node[midway, left] {$l_{10}$};
    \draw[<->, thin] (4.55,2.95) -- (4.55,3.42) node[midway, right] {$x_1$};
    \draw[<->, thin] (2.25,1.12) -- (2.25,3.42) node[midway, left] {$s$};
    \draw[<->, thin] (4.55,1.12) -- (4.55,1.72) node[midway, right] {$x_2$};
    \draw[<->, thin] (3.32,0.45) -- (3.32,1.12) node[midway, left] {$l_{20}$};
\end{tikzpicture}

\end{center}
\end{problem}
\clearpage


\begin{problem}{9.4.16}
图 9.93 中半圆柱 A 通过一根扭转常数为 k (转单位角度时受到恢复力矩为 k) 的细弹性棒钉在位于通过柱心的固定的垂直轴 O 上, 可在固定的半圆柱面 B 内转动, 构成可变电容器, 假定其电容为 $C = C_{0} \left(1 - \frac{\theta}{\pi}\right)$ , $\theta = \theta_{0}$ 时, 弹性棒未被扭曲, 写出拉格朗日函数和运动微分方程.
\begin{center}
\begin{tikzpicture}[>=Stealth, thick, font=\small]
  % Center O
  \coordinate (O) at (3.8, 1.5);

  % Outer fixed semicircular shell B (0 to 180 deg)
  \fill[pattern=north east lines]
    ($(O)+(0:1.15)$) arc (0:180:1.15) --
    ($(O)+(180:1.38)$) arc (180:0:1.38) -- cycle;
  \draw
    ($(O)+(0:1.15)$) arc (0:180:1.15)
    ($(O)+(0:1.38)$) arc (0:180:1.38)
    ($(O)+(-1.38,0)$) -- ($(O)+(1.38,0)$);

  % Inner rotating semicircular shell A (theta to theta+180 deg, theta=42 deg)
  \def\ang{42}
  \fill[pattern=north east lines]
    ($(O)+(\ang:0.85)$) arc (\ang:\ang+180:0.85) --
    ($(O)+(\ang+180:1.03)$) arc (\ang+180:\ang:1.03) -- cycle;
  \draw
    ($(O)+(\ang:0.85)$) arc (\ang:\ang+180:0.85)
    ($(O)+(\ang:1.03)$) arc (\ang:\ang+180:1.03)
    ($(O)+(\ang:1.03)$) -- ($(O)+(\ang+180:1.03)$);

  % Angle theta
  \draw[->, thin] ($(O)+(0:0.48)$) arc (0:\ang:0.48);
  \node at ($(O)+(21:0.70)$) {$\theta$};

  % Labels on capacitor
  \node at ($(O)+(145:1.62)$) {$B$};
  \node at ($(O)+(-38:0.25)$) {$O$};
  \node at ($(O)+(198:0.58)$) {$A$};
  \node at ($(O)+(233:0.75)$) {$I$};

  % Circuit connections
  % Top wire with inductor L
  \draw ($(O)+(90:1.38)$) -- (3.8, 3.2) -- (2.15, 3.2);
  \draw[decorate, decoration={coil, aspect=0.45, segment length=2.2mm, amplitude=1.4mm}]
    (2.15, 3.2) -- (0.95, 3.2);
  \node[below=4pt] at (1.55, 3.2) {$L$};
  \draw (0.95, 3.2) -- (0.2, 3.2) -- (0.2, 1.78);

  % DC source E
  \draw (0.0, 1.78) -- (0.4, 1.78);
  \draw[line width=1.6pt] (0.09, 1.62) -- (0.31, 1.62);
  \node[right=3pt] at (0.4, 1.70) {$E$};

  % Bottom wire with resistor R
  \draw (0.2, 1.62) -- (0.2, 0.2) -- (1.15, 0.2);
  \draw (1.15, 0.08) rectangle (1.85, 0.32);
  \node[above=2pt] at (1.50, 0.32) {$R$};
  \draw (1.85, 0.2) -- (3.8, 0.2) -- (O);
\end{tikzpicture}

\end{center}
\end{problem}
\clearpage


\begin{problem}{9.4.17}
可变电容器与墙式电流计如图 9.94 所示连接，设电流计中的磁场沿径向，磁感应强度为 B，线圈有 n 匝，每匝面积为 $A_{1}$ ，线圈绕轴的转动惯量为 $I, k_{1}$ 是扭转弹簧的扭转常数, 可变电容器的电容 $C = \frac{A_2}{d - x}, A_2, d$ 均为常量, $l_0$ 是极板不带电时, 质量为 $m$ 的极板平衡时的弹簧长度, $n$ 匝线圈的自感系数为 $L_1$ , 写出系统的拉格朗日函数和运动微分方程.
\begin{center}
\begin{tikzpicture}[>=Stealth, thick, font=\small, scale=0.84, transform shape]
  % Top-left ceiling support
  \fill[pattern=north east lines] (0.0, 3.72) rectangle (1.55, 3.90);
  \draw (0.0, 3.60) rectangle (1.55, 3.72);

  % Suspension wire and rotation arrow theta
  \draw (0.78, 3.60) -- (0.78, 2.95);
  \fill (0.78, 2.95) circle (1.6pt);
  \draw[->, thin] (0.60, 3.28) arc (175:395:0.18 and 0.08);
  \node[right=1pt] at (0.97, 3.25) {$\theta$};

  % Rectangular coil
  \draw (0.38, 2.00) rectangle (1.18, 2.95);
  \draw (0.47, 2.09) rectangle (1.09, 2.86);
  \node at (0.63, 2.26) {$I$};
  \node at (0.92, 2.26) {$L_1$};

  % Magnet pole N (left)
  \path[pattern=north east lines]
    (0.20, 1.92) -- (-0.30, 1.92) -- (-0.38, 2.08) -- (-0.26, 2.22) --
    (-0.38, 2.38) -- (-0.26, 2.54) -- (-0.38, 2.70) -- (-0.26, 2.86) --
    (-0.34, 3.03) -- (0.20, 3.03) -- cycle;
  \draw
    (-0.30, 1.92) -- (0.20, 1.92) -- (0.20, 3.03) -- (-0.34, 3.03)
    (-0.30, 1.92) -- (-0.38, 2.08) -- (-0.26, 2.22) -- (-0.38, 2.38) --
    (-0.26, 2.54) -- (-0.38, 2.70) -- (-0.26, 2.86) -- (-0.34, 3.03);
  \node[fill=white, inner sep=1.5pt] at (-0.07, 2.48) {$\mathrm{N}$};

  % Magnet pole S (right)
  \path[pattern=north east lines]
    (1.36, 1.92) -- (1.86, 1.92) -- (1.94, 2.08) -- (1.82, 2.22) --
    (1.94, 2.38) -- (1.82, 2.54) -- (1.94, 2.70) -- (1.82, 2.86) --
    (1.90, 3.03) -- (1.36, 3.03) -- cycle;
  \draw
    (1.86, 1.92) -- (1.36, 1.92) -- (1.36, 3.03) -- (1.90, 3.03)
    (1.86, 1.92) -- (1.94, 2.08) -- (1.82, 2.22) -- (1.94, 2.38) --
    (1.82, 2.54) -- (1.94, 2.70) -- (1.82, 2.86) -- (1.90, 3.03);
  \node[fill=white, inner sep=1.5pt] at (1.63, 2.48) {$\mathrm{S}$};

  % Radial B pointer (clear of middle wire at x=3.15)
  \draw[->, thin] (1.50, 3.35) to[out=195, in=80] (1.27, 2.88);
  \node[right=0pt] at (1.50, 3.37) {$B\,(\text{沿径向})$};

  % Torsion spring k1 and resistor R1
  \draw (0.78, 2.00) -- (0.78, 1.80);
  \draw[decorate, decoration={coil, aspect=0.45, segment length=1.6mm, amplitude=1.1mm}]
    (0.78, 1.80) -- (0.78, 1.30);
  \node[right=3pt] at (0.78, 1.55) {$k_1$};
  \draw (0.78, 1.30) -- (0.78, 1.15);
  \draw (0.66, 0.65) rectangle (0.90, 1.15);
  \node[right=2pt] at (0.90, 0.90) {$R_1$};
  \draw (0.78, 0.65) -- (0.78, 0.25);

  % Bottom wire with E1, E2 and currents q1_dot, q3_dot
  \draw (0.78, 0.25) -- (1.80, 0.25);
  \draw (1.80, 0.07) -- (1.80, 0.43);
  \draw[line width=1.6pt] (1.94, 0.14) -- (1.94, 0.36);
  \node[above=2pt] at (1.87, 0.43) {$E_1$};
  \draw (1.94, 0.25) -- (3.15, 0.25);
  \draw[->, thin] (2.35, 0.25) -- (2.65, 0.25);
  \node[below=2pt] at (2.50, 0.25) {$\dot{q}_1$};

  \draw (3.15, 0.25) -- (4.05, 0.25);
  \draw[->, thin] (3.80, 0.25) -- (3.50, 0.25);
  \node[below=2pt] at (3.65, 0.25) {$\dot{q}_3$};
  \draw (4.05, 0.07) -- (4.05, 0.43);
  \draw[line width=1.6pt] (4.19, 0.14) -- (4.19, 0.36);
  \node[above=2pt] at (4.12, 0.43) {$E_2$};
  \draw (4.19, 0.25) -- (5.45, 0.25);

  % Middle vertical branch (R2, L2, q2_dot)
  \fill (3.15, 0.25) circle (1.8pt);
  \fill (3.15, 3.66) circle (1.8pt);
  \draw (3.15, 0.25) -- (3.15, 1.05);
  \draw[->, thin] (3.15, 0.55) -- (3.15, 0.80);
  \node[right=2pt] at (3.15, 0.68) {$\dot{q}_2$};
  \draw (3.03, 1.05) rectangle (3.27, 1.55);
  \node[right=2pt] at (3.27, 1.30) {$R_2$};
  \draw (3.15, 1.55) -- (3.15, 2.20);
  \draw[decorate, decoration={coil, aspect=0.45, segment length=1.8mm, amplitude=1.3mm}]
    (3.15, 2.95) -- (3.15, 2.20);
  \node[right=4pt] at (3.15, 2.58) {$L_2$};
  \draw (3.15, 2.95) -- (3.15, 3.66);

  % Top wire to right capacitor support
  \draw (1.55, 3.66) -- (5.45, 3.66) -- (5.45, 3.45);

  % Right ceiling support
  \fill[pattern=north east lines] (4.75, 3.45) rectangle (6.15, 3.63);
  \draw (4.75, 3.33) rectangle (6.15, 3.45);

  % Spring k2 and dashed equilibrium line
  \draw[decorate, decoration={coil, aspect=0.45, segment length=1.6mm, amplitude=1.2mm}]
    (5.45, 3.33) -- (5.45, 2.65);
  \node[right=4pt] at (5.45, 2.99) {$k_2$};
  \draw[dashed] (3.95, 2.65) -- (6.15, 2.65);

  % Dimension l0
  \draw[<->, thin] (4.78, 3.33) -- (4.78, 2.65);
  \node[left=1pt] at (4.78, 2.99) {$l_0$};

  % Movable plate m
  \draw (5.45, 2.65) -- (5.45, 2.15);
  \draw (4.55, 2.05) rectangle (6.35, 2.15);
  \node[right=2pt] at (6.35, 2.10) {$m$};

  % Dimension x
  \draw[<->, thin] (5.92, 2.65) -- (5.92, 2.15);
  \node[right=1pt] at (5.92, 2.40) {$x$};

  % Fixed lower plate
  \draw (3.95, 1.35) -- (6.40, 1.35);
  \fill[pattern=north east lines] (4.45, 1.18) rectangle (6.40, 1.35);

  % Dimension d
  \draw[<->, thin] (4.10, 2.65) -- (4.10, 1.35);
  \node[left=1pt] at (4.10, 2.00) {$d$};

  % Capacitance label
  \node at (5.32, 1.72) {$C=\dfrac{A_2}{d-x}$};

  % Wire from lower plate to bottom wire
  \draw (5.45, 1.35) -- (5.45, 0.25);
\end{tikzpicture}

\end{center}
\end{problem}
\clearpage


\begin{problem}{9.4.18}
半径为 r 的圆盘的中心安装在水平的弹性棒的一端, 弹性棒的另一端是固定的, 圆盘转动 $\theta$ 角时, 弹性棒给予圆盘的力矩为 $-k_{1}\theta$ , 图 9.95 中弹簧的劲度系数为 $k_{2}$ , $m_{1}$ 、 $m_{2}$ 是两个金属圆柱的质量, 用绝缘的不可伸长的轻绳和弹簧连接并跨在圆盘上, 圆盘转动时, 圆盘与绳子间无相对滑动, 圆盘绕弹性棒的转动惯量为 I, 圆柱体可分别在两个固定的导电圆筒内作竖直运动, 圆柱体与圆筒构成可变电容器

\[
\begin{array}{l l} C _ {1} = \left\{ \begin{array}{l l} C _ {1 0} \left(1 - \frac {y _ {1}}{l _ {1}}\right) & y _ {1} \geqslant 0 \\ C _ {1 0} & y _ {1} <   0 \end{array} \right. \\ C _ {2} = \left\{ \begin{array}{l l} C _ {2 0} \left(1 - \frac {y _ {2}}{l _ {2}}\right) & y _ {2} \geqslant 0 \\ C _ {2 0} & y _ {2} <   0 \end{array} \right. \end{array}
\]

由互感将两个电路耦合，用 $q_{1}, q_{2}, y_{1}, y_{2}$ 为广义坐标，写出系统的动能、势能和运动微分方程.
\begin{center}
\begin{tikzpicture}[>=Stealth, thick, font=\small, scale=0.78, transform shape]
  % Disk centered at (3.1, 4.7)
  \coordinate (C) at (3.1, 4.7);
  \draw (C) circle (1.10);
  \draw (C) circle (0.14);

  % Pointer arrows for elastic rod and k1
  \draw[->, thin] (4.35, 5.45) -- (3.22, 4.80);
  \node[right=1pt] at (4.35, 5.48) {弹性棒};
  \draw[->, thin] (3.62, 4.46) -- (3.23, 4.65);
  \node[right=0pt] at (3.60, 4.42) {$k_1$};

  % Radius r and moment of inertia I
  \draw[->, thin] (C) -- ($(C)+(225:1.10)$);
  \node[above left=-2pt] at ($(C)+(225:0.65)$) {$r$};
  \node at (3.45, 4.02) {$I$};

  % Rotation arrow theta
  \draw[->, thin] ($(C)+(168:1.28)$) arc (168:124:1.28);
  \node at ($(C)+(148:1.52)$) {$\theta$};

  % Gravity g
  \draw[->, thin] (1.10, 5.05) -- (1.10, 4.15);
  \node[right=2pt] at (1.10, 4.60) {$g$};

  % Left string to m2
  \draw (2.00, 4.70) -- (2.00, 3.00);

  % Right string + spring k2 to m1
  \draw (4.20, 4.70) -- (4.20, 3.65);
  \draw[decorate, decoration={coil, aspect=0.45, segment length=1.6mm, amplitude=1.1mm}]
    (4.20, 3.65) -- (4.20, 2.92);
  \node[left=4pt] at (4.20, 3.28) {$k_2$};
  \draw (4.20, 2.92) -- (4.20, 2.68);
  \fill (4.20, 2.68) circle (1.6pt);

  % Dashed reference line at rim of fixed cylinders
  \draw[dashed] (1.25, 2.32) -- (5.15, 2.32);

  % Cylinder m2 (left)
  \draw (1.92, 1.82) rectangle (2.62, 3.00);
  \node at (2.38, 2.55) {$m_2$};
  \draw[<->, thin] (2.08, 1.82) -- (2.08, 3.00);
  \node[right=0pt] at (2.08, 2.12) {$l_2$};

  % Dimension y2
  \draw[thin] (2.62, 3.00) -- (3.02, 3.00);
  \draw[<->, thin] (2.90, 2.32) -- (2.90, 3.00);
  \node[right=1pt] at (2.90, 2.66) {$y_2$};

  % Fixed conducting cylinder (left)
  \draw (1.68, 2.32) -- (1.80, 2.32) -- (1.80, 1.38) -- (2.74, 1.38) -- (2.74, 2.32) -- (2.86, 2.32);
  \draw (1.80, 1.38) -- (1.80, 0.88);
  \fill (1.80, 0.88) circle (1.6pt);

  % Cylinder m1 (right)
  \draw (4.05, 1.55) rectangle (4.70, 2.68);
  \node at (4.48, 2.30) {$m_1$};
  \draw[<->, thin] (4.20, 1.55) -- (4.20, 2.68);
  \node[right=0pt] at (4.20, 1.86) {$l_1$};

  % Dimension y1
  \draw[thin] (3.62, 2.68) -- (4.05, 2.68);
  \draw[<->, thin] (3.74, 2.32) -- (3.74, 2.68);
  \node[left=1pt] at (3.74, 2.50) {$y_1$};

  % Fixed conducting cylinder (right)
  \draw (3.81, 2.32) -- (3.93, 2.32) -- (3.93, 1.38) -- (4.82, 1.38) -- (4.82, 2.32) -- (4.94, 2.32);
  \draw (4.82, 1.38) -- (4.82, 0.88);
  \fill (4.82, 0.88) circle (1.6pt);

  % Label for fixed conducting cylinders (clear of wires below)
  \node[align=center] at (3.34, 1.56) {固定的\\导电圆筒};
  \draw[->, thin] (2.82, 1.30) to[out=195, in=-90] (2.35, 1.38);
  \draw[->, thin] (3.86, 1.30) to[out=-15, in=-90] (4.30, 1.38);

  % Left flexible wire + DC source E2 + R2
  \draw[decorate, decoration={snake, amplitude=1.4mm, segment length=4.5mm}]
    (1.92, 2.95) -- (0.35, 2.95);
  \draw (0.35, 2.95) -- (0.35, 1.72);
  \draw (0.17, 1.72) -- (0.53, 1.72);
  \draw[line width=1.6pt] (0.24, 1.58) -- (0.46, 1.58);
  \node[right=3pt] at (0.53, 1.65) {$E_2$};
  \draw (0.35, 1.58) -- (0.35, 1.05);
  \draw (0.23, 0.55) rectangle (0.47, 1.05);
  \node[right=2pt] at (0.47, 0.80) {$R_2$};
  \draw (0.35, 0.55) -- (0.35, 0.10) -- (2.98, 0.10);
  \draw[->, thin] (1.40, 0.10) -- (1.70, 0.10);
  \node[above=2pt] at (1.55, 0.10) {$\dot{q}_2$};

  % Right flexible wire + AC source E1 + R1
  \draw[decorate, decoration={snake, amplitude=1.4mm, segment length=4.5mm}]
    (4.70, 2.63) -- (6.25, 2.63);
  \draw[->, thin] (5.35, 3.18) -- (5.20, 2.78);
  \node[above=1pt] at (5.45, 3.16) {软导线};
  \draw (6.25, 2.63) -- (6.25, 2.02);
  \draw (6.25, 1.97) circle (1.8pt);
  \draw (6.25, 1.48) circle (1.8pt);
  \node[fill=white, inner sep=1pt] at (6.15, 1.72) {$E_1=E_0\sin\omega t$};
  \draw (6.25, 1.43) -- (6.25, 1.05);
  \draw (6.13, 0.55) rectangle (6.37, 1.05);
  \node[right=2pt] at (6.37, 0.80) {$R_1$};
  \draw (6.25, 0.55) -- (6.25, 0.10) -- (3.62, 0.10);
  \draw[->, thin] (5.15, 0.10) -- (4.85, 0.10);
  \node[above=2pt] at (5.00, 0.10) {$\dot{q}_1$};

  % Mutual inductors L2 and L1
  \draw (2.98, 0.10) -- (2.98, 0.25);
  \draw[decorate, decoration={coil, aspect=0.45, segment length=2.0mm, amplitude=1.3mm}]
    (2.98, 0.25) -- (2.98, 0.88);
  \node[left=4pt] at (2.98, 0.55) {$L_2$};
  \draw (2.98, 0.88) -- (1.80, 0.88);

  \draw (3.62, 0.10) -- (3.62, 0.25);
  \draw[decorate, decoration={coil, aspect=0.45, segment length=2.0mm, amplitude=1.3mm}]
    (3.62, 0.88) -- (3.62, 0.25);
  \node[right=4pt] at (3.62, 0.55) {$L_1$};
  \draw (3.62, 0.88) -- (4.82, 0.88);

  \node at (3.30, 0.72) {$M$};
\end{tikzpicture}

\end{center}
\end{problem}
\clearpage


\begin{problem}{9.4.19}
图 9.96(a)、(b)、(c) 分别画了三个系统. 在 (a) 中, E 提供恒定的电压, 并假定线圈间无互感；在(b)中， $F$ 是恒定的外力，每个木块受到黏性力 $-\alpha_{i} \dot{x}_{i}$ ；在(c)中， $\tau$ 是恒定的外力矩，每个圆盘受到制动器的黏性力矩 $-b_{i} r_{i} \dot{\theta}_{i}, k_{1}, k_{2}$ 是弹簧的扭转常数。写出三个系统的拉格朗日函数，并证明：适当选择坐标，对应的物理量取相同的值，三个系统的运动微分方程相同。

(a)

(b)

(c)
\begin{center}
\begin{tikzpicture}[>=Stealth, thick, font=\small, scale=0.72, transform shape]
  % Shaft 1 and Disk I1
  \draw (0.3, 1.42) rectangle (2.5, 1.58);
  \draw[fill=white] (1.75, 0.55) rectangle (2.00, 2.45);
  \node[above=2pt] at (1.875, 2.45) {$I_1$};

  % Bearings for Disk 1
  \foreach \bx in {1.35, 2.12} {
    \fill[pattern=north east lines] (\bx, 1.58) rectangle (\bx+0.28, 1.88);
    \draw (\bx, 1.88) -- (\bx, 1.58) -- (\bx+0.28, 1.58);
    \fill[pattern=north east lines] (\bx, 1.12) rectangle (\bx+0.28, 1.42);
    \draw (\bx, 1.12) -- (\bx, 1.42) -- (\bx+0.28, 1.42);
  }

  % Torque tau and angular velocity theta1_dot
  \draw[->, thin] (0.50, 1.25) arc (240:60:0.22 and 0.30);
  \node[left=2pt] at (0.18, 1.55) {$\tau$};
  \draw[->, thin] (0.95, 1.25) arc (240:60:0.16 and 0.26);
  \node[above=2pt] at (0.95, 1.72) {$\dot{\theta}_1$};

  % Damper b1
  \draw (1.50, 0.42) -- (2.25, 0.42);
  \fill[pattern=north east lines] (1.50, 0.24) rectangle (2.25, 0.42);
  \draw[->, thin] (2.35, 0.28) to[out=135, in=0] (1.92, 0.48);
  \node[right=1pt] at (2.35, 0.28) {$b_1$};

  % Spring k1 (3 clear loops)
  \draw (2.50, 1.50) -- (2.62, 1.50);
  \draw[decorate, decoration={coil, aspect=0.45, segment length=1.6mm, amplitude=1.2mm}]
    (2.62, 1.50) -- (3.30, 1.50);
  \draw (3.30, 1.50) -- (3.42, 1.50);
  \node[above=4pt] at (2.96, 1.50) {$k_1$};

  % Shaft 2 and Disk I2
  \draw (3.42, 1.42) rectangle (5.32, 1.58);
  \draw[fill=white] (4.07, 0.25) rectangle (4.32, 2.75);
  \node[above=2pt] at (4.195, 2.75) {$I_2$};

  % Bearings for Disk 2
  \foreach \bx in {3.67, 4.44} {
    \fill[pattern=north east lines] (\bx, 1.58) rectangle (\bx+0.28, 1.88);
    \draw (\bx, 1.88) -- (\bx, 1.58) -- (\bx+0.28, 1.58);
    \fill[pattern=north east lines] (\bx, 1.12) rectangle (\bx+0.28, 1.42);
    \draw (\bx, 1.12) -- (\bx, 1.42) -- (\bx+0.28, 1.42);
  }

  % Angular velocity theta2_dot
  \draw[->, thin] (4.94, 1.25) arc (240:60:0.16 and 0.26);
  \node[above=2pt] at (4.94, 1.72) {$\dot{\theta}_2$};

  % Damper b2
  \draw (3.82, 0.12) -- (4.57, 0.12);
  \fill[pattern=north east lines] (3.82, -0.06) rectangle (4.57, 0.12);
  \draw[->, thin] (4.67, -0.02) to[out=135, in=0] (4.24, 0.18);
  \node[right=1pt] at (4.67, -0.02) {$b_2$};

  % Spring k2 (3 clear loops)
  \draw (5.32, 1.50) -- (5.44, 1.50);
  \draw[decorate, decoration={coil, aspect=0.45, segment length=1.6mm, amplitude=1.2mm}]
    (5.44, 1.50) -- (6.12, 1.50);
  \draw (6.12, 1.50) -- (6.24, 1.50);
  \node[above=4pt] at (5.78, 1.50) {$k_2$};

  % Shaft 3 and Disk I3
  \draw (6.24, 1.42) rectangle (7.89, 1.58);
  \draw[fill=white] (6.89, 0.45) rectangle (7.14, 2.55);
  \node[above=2pt] at (7.015, 2.55) {$I_3$};

  % Bearings for Disk 3
  \foreach \bx in {6.49, 7.26} {
    \fill[pattern=north east lines] (\bx, 1.58) rectangle (\bx+0.28, 1.88);
    \draw (\bx, 1.88) -- (\bx, 1.58) -- (\bx+0.28, 1.58);
    \fill[pattern=north east lines] (\bx, 1.12) rectangle (\bx+0.28, 1.42);
    \draw (\bx, 1.12) -- (\bx, 1.42) -- (\bx+0.28, 1.42);
  }

  % Angular velocity theta3_dot
  \draw[->, thin] (7.66, 1.25) arc (240:60:0.16 and 0.26);
  \node[above=2pt] at (7.66, 1.72) {$\dot{\theta}_3$};

  % Damper b3
  \draw (6.64, 0.32) -- (7.39, 0.32);
  \fill[pattern=north east lines] (6.64, 0.14) rectangle (7.39, 0.32);
  \draw[->, thin] (7.49, 0.18) to[out=135, in=0] (7.06, 0.38);
  \node[right=1pt] at (7.49, 0.18) {$b_3$};
\end{tikzpicture}

\end{center}
\begin{center}
\begin{tikzpicture}[>=Stealth, thick, font=\small, scale=0.80, transform shape]
  % Left vertical branch with DC source E and current q1_dot
  \draw (0.4, 0.3) -- (0.4, 1.23);
  \draw (0.20, 1.37) -- (0.60, 1.37);
  \draw[line width=1.6pt] (0.28, 1.23) -- (0.52, 1.23);
  \node[left=3pt] at (0.20, 1.30) {$E$};
  \draw (0.4, 1.37) -- (0.4, 2.3);
  \draw[->, thin] (0.4, 1.65) -- (0.4, 1.95);
  \node[right=2pt] at (0.4, 1.80) {$\dot{q}_1$};

  % Top wire with L1, L2, q2_dot, L3
  \draw (0.4, 2.3) -- (1.05, 2.3);
  \draw[decorate, decoration={coil, aspect=0.45, segment length=2.2mm, amplitude=1.4mm}]
    (1.05, 2.3) -- (1.95, 2.3);
  \node[below=4pt] at (1.50, 2.3) {$L_1$};
  \draw (1.95, 2.3) -- (3.15, 2.3);
  \draw[decorate, decoration={coil, aspect=0.45, segment length=2.2mm, amplitude=1.4mm}]
    (3.15, 2.3) -- (4.05, 2.3);
  \node[below=4pt] at (3.60, 2.3) {$L_2$};
  \draw (4.05, 2.3) -- (5.35, 2.3);
  \draw[->, thin] (4.25, 2.3) -- (4.55, 2.3);
  \node[below=2pt] at (4.40, 2.3) {$\dot{q}_2$};
  \draw[decorate, decoration={coil, aspect=0.45, segment length=2.2mm, amplitude=1.4mm}]
    (5.35, 2.3) -- (6.25, 2.3);
  \node[below=4pt] at (5.80, 2.3) {$L_3$};
  \draw (6.25, 2.3) -- (6.9, 2.3);

  % Right vertical branch with q3_dot
  \draw (6.9, 2.3) -- (6.9, 0.3);
  \draw[->, thin] (6.9, 1.45) -- (6.9, 1.15);
  \node[left=2pt] at (6.9, 1.30) {$\dot{q}_3$};

  % Bottom wire with R1, R2, R3
  \draw (0.4, 0.3) -- (1.15, 0.3);
  \draw (1.15, 0.18) rectangle (1.85, 0.42);
  \node[above=2pt] at (1.50, 0.42) {$R_1$};
  \draw (1.85, 0.3) -- (3.35, 0.3);
  \draw (3.35, 0.18) rectangle (4.05, 0.42);
  \node[above=2pt] at (3.70, 0.42) {$R_2$};
  \draw (4.05, 0.3) -- (5.45, 0.3);
  \draw (5.45, 0.18) rectangle (6.15, 0.42);
  \node[above=2pt] at (5.80, 0.42) {$R_3$};
  \draw (6.15, 0.3) -- (6.9, 0.3);

  % Middle vertical branch 1 (q4_dot and C1)
  \fill (2.6, 2.3) circle (1.8pt);
  \fill (2.6, 0.3) circle (1.8pt);
  \draw (2.6, 2.3) -- (2.6, 1.24);
  \draw[->, thin] (2.6, 1.95) -- (2.6, 1.65);
  \node[right=2pt] at (2.6, 1.80) {$\dot{q}_4$};
  \draw (2.38, 1.24) -- (2.82, 1.24);
  \draw (2.38, 1.08) -- (2.82, 1.08);
  \node[right=3pt] at (2.82, 1.16) {$C_1$};
  \draw (2.6, 1.08) -- (2.6, 0.3);

  % Middle vertical branch 2 (C2 and q5_dot)
  \fill (4.8, 2.3) circle (1.8pt);
  \fill (4.8, 0.3) circle (1.8pt);
  \draw (4.8, 2.3) -- (4.8, 1.48);
  \draw (4.58, 1.48) -- (5.02, 1.48);
  \draw (4.58, 1.32) -- (5.02, 1.32);
  \node[right=3pt] at (5.02, 1.40) {$C_2$};
  \draw (4.8, 1.32) -- (4.8, 0.3);
  \draw[->, thin] (4.8, 0.95) -- (4.8, 0.65);
  \node[right=2pt] at (4.8, 0.80) {$\dot{q}_5$};
\end{tikzpicture}

\end{center}
\begin{center}
\begin{tikzpicture}[>=Stealth, thick, font=\small, scale=0.80, transform shape]
  % Reference wall and ground
  \draw (0, 2.60) -- (0, 0) -- (7.2, 0);

  % Blocks m1, m2, m3
  \draw (1.30, 0.14) rectangle (2.20, 0.82);
  \node at (1.75, 0.48) {$m_1$};
  \draw (3.60, 0.14) rectangle (4.50, 0.82);
  \node at (4.05, 0.48) {$m_2$};
  \draw (5.90, 0.14) rectangle (6.80, 0.82);
  \node at (6.35, 0.48) {$m_3$};

  % Force F
  \draw[->, thin] (0.55, 0.48) -- (1.30, 0.48);
  \node[above=1pt] at (0.92, 0.48) {$F$};

  % Springs k1, k2
  \draw (2.20, 0.48) -- (2.45, 0.48);
  \draw[decorate, decoration={zigzag, segment length=2.4mm, amplitude=1.3mm}]
    (2.45, 0.48) -- (3.35, 0.48);
  \draw (3.35, 0.48) -- (3.60, 0.48);
  \node[above=4pt] at (2.90, 0.48) {$k_1$};

  \draw (4.50, 0.48) -- (4.75, 0.48);
  \draw[decorate, decoration={zigzag, segment length=2.4mm, amplitude=1.3mm}]
    (4.75, 0.48) -- (5.65, 0.48);
  \draw (5.65, 0.48) -- (5.90, 0.48);
  \node[above=4pt] at (5.20, 0.48) {$k_2$};

  % Viscous friction arrows alpha1, alpha2, alpha3
  \draw[->, thin] (2.25, -0.22) to[out=135, in=-10] (1.75, 0.10);
  \node[right=1pt] at (2.25, -0.22) {$\alpha_1$};
  \draw[->, thin] (4.55, -0.22) to[out=135, in=-10] (4.05, 0.10);
  \node[right=1pt] at (4.55, -0.22) {$\alpha_2$};
  \draw[->, thin] (6.85, -0.22) to[out=135, in=-10] (6.35, 0.10);
  \node[right=1pt] at (6.85, -0.22) {$\alpha_3$};

  % Dimension lines x1, x2, x3
  \draw[thin] (1.30, 1.15) -- (1.30, 1.50);
  \draw[<->, thin] (0, 1.34) -- (1.30, 1.34);
  \node[above=1pt] at (0.65, 1.34) {$x_1$};

  \draw[dashed, thin] (3.60, 0.82) -- (3.60, 1.65);
  \draw[thin] (3.60, 1.65) -- (3.60, 2.00);
  \draw[<->, thin] (0, 1.84) -- (3.60, 1.84);
  \node[above=1pt] at (1.80, 1.84) {$x_2$};

  \draw[dashed, thin] (5.90, 0.82) -- (5.90, 2.15);
  \draw[thin] (5.90, 2.15) -- (5.90, 2.50);
  \draw[<->, thin] (0, 2.34) -- (5.90, 2.34);
  \node[above=1pt] at (2.95, 2.34) {$x_3$};
\end{tikzpicture}

\end{center}
\end{problem}
\clearpage


\begin{problem}{9.4.20}
列出图 9.97(a)、(b) 系统的运动微分方程, 质量为 $m_{1}$ 、 $m_{2}$ 的两木块底面受到黏性阻力, 黏度分别为 $\alpha_{1}$ 、 $\alpha_{2}$ , 减震器也提供一个黏性阻力, 黏度为 $\alpha_{3}$ , 证明两个系统可以是等效的.

(a)

(b)
\begin{center}
\begin{tikzpicture}[>=Stealth, thick, font=\small, scale=0.85, transform shape]
  % Left vertical branch with AC source E = E0 sin omega t
  \draw (0.4, 0.3) -- (0.4, 1.29);
  \draw (0.4, 2.01) -- (0.4, 3.0);
  \draw (0.4, 1.95) circle (1.8pt);
  \draw (0.4, 1.35) circle (1.8pt);
  \node[right=3pt] at (0.4, 1.95) {$+$};
  \node[right=3pt] at (0.4, 1.35) {$-$};
  \node[fill=white, inner sep=1.5pt] at (0.50, 1.65) {$E=E_0\sin\omega t$};

  % Top wire with currents q1_dot and q2_dot
  \draw (0.4, 3.0) -- (5.8, 3.0);
  \draw[->, thin] (2.25, 3.0) -- (1.95, 3.0);
  \node[below=2pt] at (2.10, 3.0) {$\dot{q}_1$};
  \fill (3.8, 3.0) circle (1.8pt);
  \draw[->, thin] (4.95, 3.0) -- (4.65, 3.0);
  \node[below=2pt] at (4.80, 3.0) {$\dot{q}_2$};

  % Middle vertical branch (R3, C3, q3_dot)
  \draw (3.8, 3.0) -- (3.8, 2.60);
  \draw (3.68, 2.10) rectangle (3.92, 2.60);
  \node[left=2pt] at (3.68, 2.35) {$R_3$};
  \draw (3.8, 2.10) -- (3.8, 1.45);
  \draw (3.58, 1.45) -- (4.02, 1.45);
  \draw (3.58, 1.29) -- (4.02, 1.29);
  \node[left=3pt] at (3.58, 1.37) {$C_3$};
  \draw (3.8, 1.29) -- (3.8, 0.3);
  \draw[->, thin] (3.8, 0.65) -- (3.8, 0.95);
  \node[right=2pt] at (3.8, 0.80) {$\dot{q}_3$};
  \fill (3.8, 0.3) circle (1.8pt);

  % Bottom wire (R1, C1, L1)
  \draw (0.4, 0.3) -- (0.95, 0.3);
  \draw (0.95, 0.18) rectangle (1.55, 0.42);
  \node[above=2pt] at (1.25, 0.42) {$R_1$};
  \draw (1.55, 0.3) -- (2.02, 0.3);
  \draw (2.02, 0.08) -- (2.02, 0.52);
  \draw (2.18, 0.08) -- (2.18, 0.52);
  \node[above=3pt] at (2.10, 0.52) {$C_1$};
  \draw (2.18, 0.3) -- (2.60, 0.3);
  \draw[decorate, decoration={coil, aspect=0.45, segment length=2.2mm, amplitude=1.4mm}]
    (2.60, 0.3) -- (3.35, 0.3);
  \node[above=4pt] at (2.98, 0.3) {$L_1$};
  \draw (3.35, 0.3) -- (5.8, 0.3);

  % Right vertical branch (L2, C2, R2)
  \draw (5.8, 3.0) -- (5.8, 2.70);
  \draw[decorate, decoration={coil, aspect=0.45, segment length=2.2mm, amplitude=1.4mm}]
    (5.8, 2.70) -- (5.8, 2.00);
  \node[right=4pt] at (5.8, 2.35) {$L_2$};
  \draw (5.8, 2.00) -- (5.8, 1.65);
  \draw (5.58, 1.65) -- (6.02, 1.65);
  \draw (5.58, 1.49) -- (6.02, 1.49);
  \node[right=3pt] at (6.02, 1.57) {$C_2$};
  \draw (5.8, 1.49) -- (5.8, 1.15);
  \draw (5.68, 0.65) rectangle (5.92, 1.15);
  \node[right=2pt] at (5.92, 0.90) {$R_2$};
  \draw (5.8, 0.65) -- (5.8, 0.3);
\end{tikzpicture}

\end{center}
\begin{center}
\begin{tikzpicture}[>=Stealth, thick, font=\small, scale=0.78, transform shape]
  % Fixed walls and floor
  \fill[pattern=north east lines] (-0.18, 0) rectangle (0, 2.25);
  \fill[pattern=north east lines] (6.80, 0) rectangle (6.98, 2.25);
  \draw (0, 3.25) -- (0, 0) -- (6.80, 0) -- (6.80, 2.25);

  % Blocks m1 and m2
  \draw (2.05, 0.28) rectangle (2.95, 1.90);
  \node at (2.50, 1.09) {$m_1$};
  \node at (2.50, 0.14) {$\alpha_1$};

  \draw (4.25, 0.28) rectangle (5.45, 1.15);
  \node at (4.85, 0.72) {$m_2$};
  \node at (4.85, 0.14) {$\alpha_2$};

  % Springs k1, k3, k2
  \draw (0, 0.72) -- (0.45, 0.72);
  \draw[decorate, decoration={zigzag, segment length=2.4mm, amplitude=1.3mm}]
    (0.45, 0.72) -- (1.60, 0.72);
  \draw (1.60, 0.72) -- (2.05, 0.72);
  \node[above=4pt] at (1.02, 0.72) {$k_1$};

  \draw (2.95, 0.72) -- (3.15, 0.72);
  \draw[decorate, decoration={zigzag, segment length=2.4mm, amplitude=1.3mm}]
    (3.15, 0.72) -- (4.05, 0.72);
  \draw (4.05, 0.72) -- (4.25, 0.72);
  \node[above=4pt] at (3.60, 0.72) {$k_3$};

  \draw (5.45, 0.72) -- (5.65, 0.72);
  \draw[decorate, decoration={zigzag, segment length=2.4mm, amplitude=1.3mm}]
    (5.65, 0.72) -- (6.60, 0.72);
  \draw (6.60, 0.72) -- (6.80, 0.72);
  \node[above=4pt] at (6.12, 0.72) {$k_2$};

  % Driving force F = F0 sin omega t
  \draw[->, thin] (0.85, 1.25) -- (2.05, 1.25);
  \node[above=1pt] at (1.05, 1.28) {$F=F_0\sin\omega t$};

  % Dashpot between m1 and m2
  \draw (2.95, 1.50) -- (4.85, 1.50);
  \draw (4.85, 1.22) -- (4.85, 1.78);
  \draw (4.45, 1.70) -- (4.45, 1.85) -- (5.25, 1.85) -- (5.25, 1.15);
  \draw (4.45, 1.15) -- (4.45, 1.30);
  \foreach \px/\py in {4.98/1.68, 5.08/1.55, 4.98/1.40, 5.12/1.30, 5.02/1.23, 5.15/1.72}
    \fill (\px, \py) circle (0.7pt);

  % Dimension lines x1 and x2
  \draw[dashed, thin] (2.95, 1.90) -- (2.95, 2.35);
  \draw[thin] (2.95, 2.35) -- (2.95, 2.65);
  \draw[<->, thin] (0, 2.50) -- (2.95, 2.50);
  \node[above=1pt] at (1.48, 2.50) {$x_1$};

  \draw[dashed, thin] (5.45, 1.15) -- (5.45, 2.95);
  \draw[thin] (5.45, 2.95) -- (5.45, 3.25);
  \draw[<->, thin] (0, 3.10) -- (5.45, 3.10);
  \node[above=1pt] at (2.72, 3.10) {$x_2$};
\end{tikzpicture}

\end{center}
\end{problem}
\clearpage

